%========================================
%. CONFIGURATION ARTICLE  COURS  SEC 101
%========================================
% 				CONTRIBUTIONS
%========================================

%************************************************************
% Chargement des variables du modèle
%************************************************************
%-------------------------------------------------------------
%               FR CYBERDEF SECOPS COURSE
%                          CONFIG NOTES File
%          SET BEAMER or Article/book Mode
%                    PRZMODE OR DOCMODE
%                             2020 eduf@ction
%-------------------------------------------------------------

% insert InterTitlle for each file compose main "Notes"
%\newcommand{\INTERTITLE}{OK}

 %------------------------------------------
 % Use global project CYBERDEF 
 %            Configuration file
 %------------------------------------------
%-----------------------------------------------
%               FR CYBERDEF SECOPS COURSE
%                         GLOBAL CONFIG  File
%                             2020 eduf@ction
%-----------------------------------------------

% just used for quick access to samples on project editing tool
%%-------------------------------------------------------------
%               FR CYBERDEF SECOPS COURSE
%              SAMPLES for edx & u  commands
%                             2020 eduf@ction
%-------------------------------------------------------------


% 												I M A G E S   
%-------------------------------------------------------------------------
%\upicture{../Latex/Sources/EDU/Pictures/img-liovar-incidents}{Incidents}{01}{lb:mtbf}
%-------------------------------------------------------------------------


%\mode<all>{\picframe{../Latex/Sources/EDU/Pictures/img-cycle}{Cycle de vie de gouvernance Cyberdef}{0.3}{lbl:cycle}}


% Planche PRZ 
% begin frame \mode<all>{\texframe{Déroulement}{présentation Gestion de la menace}
%{

%}} % end frame



%\begin{table}[!hbtp]
%\resizebox{\textwidth}{!}{%
%\begin{tabular}{ l|l l}
%%\begin{tabular}{l*{3}{S[table-format=2.1]}}
%\toprule
%\textbf{Package\_name }                 & \textbf{C\&C }                                        & \textbf{sha256}                                                            \\ \hline
%caracal.raceinspace.astronaut & http://api.lulquid.xyz                       & f1d32c17a169574369088...\\ \hline
%
%\bottomrule
%\end{tabular}
%}
%\caption{Exemple de  données de TI (Threat Intelligence)}
%\label{tab:ExempleTI}
%\end{table}

%----------------------------
% Path
\def\upath{../Latex}

%-------------------------------------------------------------
%			 FR CYBERDEF SECOPS COURSE
% 					INIT Commands File
% 		Select BEAMER or Article/book Mode
% 			with sub extension .doc or .prz
%                           2020 eduf@ction
%-------------------------------------------------------------
\RequirePackage{xstring}
\RequirePackage[realmainfile]{currfile}

\newcommand{\udocname}{\jobname}

\IfSubStr*{\currfilebase}{.prz}{\newcommand{\PRZMODE}{YES}}{\newcommand{\DOCMODE}{YES}}
   
\IfSubStr*{\currfilebase}{.edt}{\newcommand{\EDTMODE}{YES}}{\newcommand{\NOMODE}{YES}}

\input{../Latex/Templates/Magic/commons.inc/edx.library.tex}

\def\utemplatepath{../Latex/Templates/Magic}

\def\edxlibspath{\utemplatepath/edxlibs}
%\input{\edxlibspath/load.tex} à mettre dans les class article, book ...

\def\uauthorphoto{\utemplatepath/commons.inc/CommonsPictures/missionauthorphoto.png}


\def\utemplatetheme{cnam}
\def\ubeamermodel{edx}
\def\uproject{CNAMSEC101}


\newcommand{\upictureext}{pdf}

 %----------------------------
 % Basic Var
 
 %\setcounter{cntx}{\year+1}
\newcommand{\INFODistrib}{Notes de cours SECOPS 2022-2023}
\newcommand{\uJournalInfo}{Orange Cyber et Cloud avec CNAM Bretagne, Cybersécurité SEC101, edufaction}

\newcommand{\udescription}{La sécurité opérationnelle au coeur de la cyberdéfense d'entreprise}

\newcommand{\uUniversity}{Conservatoire National des Arts et Métiers}
\newcommand{\uCompany}{Orange Campus Cyber}
 %\newcommand{\uJournalInfo}{CNAM Bretagne, Cybersécurité SEC101, eduf@ction}
\newcommand{\uinstituteshort}{CNAM}
\newcommand{\uinstitute}{\uUniversity}
\newcommand{\uchaire}{Chaire de Cybersécurité}
\newcommand{\uversion}{Cnam Bretagne}





%----------------------------
% All other variables are defined in
 % file : *.lesson.tex
\newcommand{\ushorttitle}{CYBERDEF SEC101}
\newcommand{\uCoursetittle}{Eléments de Sécurité Opérationnelle}
\newcommand{\uMinilogo}{Cyberdéfense d'entreprise}
\newcommand{\uCourseLongName} {Eléments de sécurité opérationnelle en cyberdéfense d'entreprise}
%\newcommand{\utitle}{Cyberdéfense d'entreprise}

%----------------------------
% Book versus Course Notes
%----------------------------
% Use "e document' "e cours"
\newcommand{\edoc}{e document\xspace}
\newcommand{\ecours}{e cours\xspace} 
\newcommand{\etitle}{éléments de sécurité opérationnelle en cyberdéfense d'entreprise} 
\newcommand{\fichetech}{une fiche TECHNO}


%%----------------------------
%% Spécifique BOOK
% %----------------------------
\newcommand{\printer}{Publication limitée pour le \uUniversity \xspace et Orange Campus Cyber}

%%----------------------------
% Attention : Use hard coded Pictures Path
% (PicframeS)
\newcommand{\picpath}{../Latex/Sources/EDU/SRC1/Pictures}
%-------------------------------------------------------------

%%----------------------------
%% Spécifique COURSE
% %----------------------------

\newcommand {\parttitlecourse}{TITRE PARTIE (config)}
\newcommand{\coursenumber}{1}

\newcommand{\uGlossaryAcronyms}{
\loadglsentries{../Latex/Sources/commons/acronymes.tex}
\loadglsentries{../Latex/Sources/commons/glossaire.tex}
}

\newcommand{\ucourseresume}{\normalsize Aborder la sécurité des systèmes d’information sous l’angle d'une sécurité dynamique est un axe qui depuis quelques années apporte de nouvelle manière d’aborder la protection, la défense, et la résilience des systèmes d’information. La transformation digitale de l'entreprise modifie et rend plus flous les périmètres des systèmes d’informations. Cela nécessite une approche élargie du risque numérique et des nouvelles architectures de systèmes d'information intégrant des technologies de cloud et de cyberdéfense. Malgré la mise en place de mesures et de technologies de protection de plus en plus élaborées, l'impact d'une attaque ayant franchi ces frontières poreuses à considérablement augmenté.  Cette compilation des notes de cours élaborée dans le cadre d'un cours d'introduction à la cybersécurité et la cyberdéfense d'entreprise aborde les grands éléments fondamentaux permettant d'appréhender le domaine. Il permet d'en comprendre les enjeux, les codes, et acteurs. Protéger l'ensemble de l'entreprise alors qu'il est complexe de définir ses frontières est illusoire. Identifier les actifs essentiels ou vitaux et mettre en place les moyens adaptés à leur protection et leur défense est une démarche tactique qui permet de graduellement réduire ses cyber-risques tout en assurant un pilotage globale de la gouvernance de la sécurité sur ses volets de conformité.}

 %------------------------------------------
 
\renewcommand{\uMinilogo}{\includegraphics[width=0.04\paperwidth]{\utemplatepath/commons.inc/CommonsPictures/\utemplatetheme.theme/shield-20.pdf}}
 
\newcommand {\Ucontribute}
{
%%-------------------------------------------------------------
%               FR CYBERDEF SECOPS COURSE
%                            Contributions Rules
%                              2020 eduf@ction
%-------------------------------------------------------------

% \section{Contributions}

% \subsection{Comment contribuer}

Les notes et les présentations sont réalisées sous \hologo{LaTeX}. 

Vous pouvez contribuer au projet du cours CYBERDEF101. Les contributions peuvent se faire sous deux formes :

\begin{itemize}
  \item Corriger, amender, améliorer les notes publiées. A chaque session des modifications et évolutions sont apportées pour tenir compte des corrections de fond et de formes. 
  \item Ajouter, compléter, modifier des parties de notes sur la base de votre lecture du cours et de vos expertises dans chacun des domaines évoqués.
\end{itemize}

 Les fichiers sources sont publiés sur GITHUB dans l'espace : \ulink{https://github.com/edufaction/CYBERDEF101}{(edufaction/CYBERDEF101)}.
 
%%-------------------------------------------------------------
%               FR CYBERDEF SECOPS COURSE
%                Contributions DATAS & NAMES
%                              2020 eduf@ction
%-------------------------------------------------------------

\subsection{Les contributeurs/auteurs du cours}

\subsubsection{co-auteurs}

%-------------------------------------------------------------
%               FR CYBERDEF SECOPS COURSE
%                Contributions DATAS & NAMES
%                              2020 eduf@ction
%-------------------------------------------------------------
%           CO AUTEURS - Rédacteur de chapitre
%-------------------------------------------------------------

\coauthor{David  BATANY}{2019-2020}{Architecture et fonctionnement des Botnets}{Cnam SEC101}

\subsubsection{contributeurs}

%-------------------------------------------------------------
%               FR CYBERDEF SECOPS COURSE
%                Contributions DATAS & NAMES
%                              2020 eduf@ction
%-------------------------------------------------------------
%                 CONTRIBUTIONS AU COURS
%-------------------------------------------------------------


\coauthor{Céline JUBY}{2020}{Contributions d'amélioration et relectures}{Orange Cyberdefense}


} 

%------------------------------------------
 %     BEAMER OU ARTICLE MODEL
 %------------------------------------------
  
\ifdefined\PRZMODE 
	\providecommand{\umainload}{%-------------------------------------------------------------
%               FR CYBERDEF SECOPS COURSE
%                          BEAMER MAIN FILE
%                             2020 eduf@ction
%-------------------------------------------------------------

%\begin{frame}
%\frametitle<presentation>{TITRE DIAPO}
% frame text
%\end{frame}

%-------------------------------------------------------------
\documentclass[ignorenonframetext, 8pt, allowframebreaks]{beamer}

%-----------------------------------------
% Load Edxlibs in MAGIC path
%-----------------------------------------
% Need : \edxlibspath defined

\input{\edxlibspath/edxlib-placeholder.tex}
\input{\edxlibspath/edxlib-epigraphe.tex}
\input{\edxlibspath/edxlib-quizzmoodle.tex}


\usepackage[utf8]{inputenc}
\usepackage[T1]{fontenc}
\usepackage{csquotes}
\usepackage[french]{babel}
\usepackage{lmodern}
\usepackage{multicol}

\input{\utemplatepath/commons.inc/CommonsPictures/\utemplatetheme.theme/colors.inc} % set uMainColor

%-------------------------------------------------------------
%               FR CYBERDEF SECOPS COURSE
%                                  EDX PACKAGES
%.    Commons packages for BOOK, NOTES (.doc) 
%                            and BEAMER (.prz)
%                              2020 eduf@ction
%-------------------------------------------------------------

%\usepackage[hyphens,spaces,obeyspaces]{url}
%\usepackage[hyphens,spaces,obeyspaces]{url}



\PassOptionsToPackage{hyphens,spaces,obeyspaces}{url}

\usepackage[utf8]{inputenc}
\usepackage[french]{babel}
\usepackage[T1]{fontenc}
%\usepackage{markdown}



\usepackage{ccicons}

\usepackage{tikz} % Required for drawing custom shapes
\usetikzlibrary{tikzmark} 
\usetikzlibrary{mindmap,trees, backgrounds}


\usepackage{verbatim}
\usepackage{wrapfig}
\usepackage{xpatch}

\usepackage{hologo}

\usepackage[figure]{totalcount}

\usepackage{fontawesome}

\usepackage[skins]{tcolorbox} % don't forget SKINS option

\usepackage{listings}
\usepackage{upquote}

\lstset{
upquote=true,columns=flexible, basicstyle=\ttfamily,
language=HTML, 
frameround=tttt,
commentstyle=\color{gray},
identifierstyle=\color{blue},
keywordstyle=\color{uMainColor}\bfseries,
xleftmargin=2em,
xrightmargin=2em,
aboveskip=\topsep,
belowskip=\topsep, 
frame=single,
rulecolor=\color{uMainColor},
backgroundcolor=\color{uMainColor!5},
breaklines,
breakindent=1.5em,
showspaces=false,
showstringspaces=false,
showtabs=false,
xleftmargin=2em,
xrightmargin=2em,
aboveskip=1em,
belowskip=1em,
}



\usepackage{graphicx} % Required for including pictures
%\graphicspath{{Pictures/}} % Specifies the directory where pictures are stored

\usepackage{booktabs} % Required for nicer horizontal rules in tables

\usepackage{setspace}
\setstretch{1,2}

\usepackage{csquotes}

\usepackage{fancyhdr}

\usepackage{makeidx} 

\usepackage{xspace} 

\usepackage{floatflt} 

\usepackage{hyperref}
\usepackage{cleveref}

%pdfdisplaydoctitle={\utitle}, Attention le titre ne doit pas contenir de saut de ligne 

\hypersetup{
pdftitle={\uCoursetittle},
pdfauthor={\uauthor},
hidelinks,
backref=true,
pagebackref=true,
hyperindex=true,
pdfcenterwindow = true,
colorlinks=true,
breaklinks=true,
bookmarks=true,
allcolors=uMainColor,
pdfsubject = {\udescription},
pdfcreator = {MACTEX on TEXPAD},
bookmarksopen=true}


% urlcolor=cnam,
% runcolor=uMainColor,
%\pdfstringdefDisableCommands{%
%  \def\\{}%
%  \def\texttt#1{<#1>}%
%}

\usepackage[style=numeric,citestyle=numeric,sorting=nyt,sortcites=true,autopunct=true,hyperref=true,abbreviate=false,backref=true,backend=biber]{biblatex}

%\usepackage{cleveref} old position

\usepackage{multicol}

%--------------------------------------------------------
% pour la génération du QRCODE en page de garde d'article
\usepackage{qrcode} 

%--------------------------------------------------------
% POur la date d'impression en premiere page d'article
\usepackage[useregional]{datetime2}

%--------------------------------------------------------
% Gestion des glossaires et Sources/EDU/acronymes
% Un fichier acroymes.tex centralisé
\usepackage[acronym]{glossaries}

%\makeglossaries
\makenoidxglossaries 

\loadglsentries{Latex/Sources/commons/acronymes.tex}
\loadglsentries{Latex/Sources/commons/glossaire.tex}


%\renewcommand{\bibname}{R\'ef\'erences}

%--------------------------------------------------------
% Pour exclure des "environnement", ici techworkbox par exemple, travaux pour étudiants, non affiché dans le BOOK.
\usepackage{comment}


%--------------------------------------------------------
% pour filigramme sur les pages
\usepackage{eso-pic}

%--------------------------------------------------------
% for \gappto{\UrlBreaks}{\UrlOrds}
% rather than 
%\makeatletter
% \g@addto@macro{\UrlBreaks}{\UrlOrds}
% \makeatother

\usepackage{etoolbox}

\usepackage{breakurl}


\newcommand{\ucite}[1]{\cite{#1}}
%\newcommand{\ucite}[1]{REF BIB #1}
\newcommand{\ugls}[1]{\gls{#1}}
%\newcommand{\ugls}[1]{GLOSSAIRE (#1)}
%\usepackage[Rejne]{fncychap}

%\usepackage{pythontex} % need shell escape

%
%\makeatletter
%\DeclareRobustCommand{\afilename}[1]{%
% \begingroup
%  % \lstname seems to change hyphens into \textendash
%  \def\textendash{-}%
%  \afilename@parse{#1}%
%  \edef\afilename@base{\detokenize\expandafter{\afilename@base}}%
%  \texttt{\afilename@base.\afilename@ext}%
% \endgroup
%}
%\makeatother

%-------------------------------------------------------------
%               FR CYBERDEF SECOPS COURSE
%                                  EDX PACKAGES
%.    Commons COMMANDS for BOOK, NOTES (.doc) 
%                            and BEAMER (.prz)
%                              2020 eduf@ction
%-------------------------------------------------------------

\newcommand{\ucite}[1]{\cite{#1}}
%\newcommand{\ucite}[1]{REF BIB #1}
\newcommand{\ugls}[1]{\gls{#1}}
%\newcommand{\ugls}[1]{GLOSSAIRE (#1)}
\newcommand{\uac}[1]{\acrshort{#1}}



% \newcommand{\ExtractBeforeLastSlash}[1]{\directlua{extract_before_last_slash("#1")}}
% \newcommand{\ExtractAfterLastSlash}[1]{\directlua{extract_after_last_slash("#1")}}
% \newcommand{\RemoveExtension}[1]{\directlua{remove_extension("#1")}}
% \newcommand{\ExtractExtension}[1]{\directlua{extract_extension("#1")}}


\newcommand{\includepicturesubpath}[1]
{#1}

%{\ExtractBeforeLastSlash{#1}/\upicturethemepath\RemoveExtension{\ExtractAfterLastSlash{#1}}.\ExtractExtension{#1}}
    

%---------------------------------------------------
%										 	AUDIO NOTE 
%---------------------------------------------------

\newcommand\audionote[2]{%
  \note{#2}
  \immediate\write\audioexport{==AudioSlide=============}
  \immediate\write\audioexport{#1}
  \immediate\write\audioexport{\insertframenumber}
  \immediate\write\audioexport{---------------------------------}
  \immediate\write\audioexport{#2}
  \immediate\write\audioexport{---------------------------------}
}%

%---------------------------------------------------
%										 	TEXT STYLE SHORTCUT 
%---------------------------------------------------

\newcommand{\tcb}[1]{\textbf{\textcolor{uMainColor}{#1}}}

\newcommand{\voir}[1]{(Schéma~\ref{#1},~page~\pageref{#1}}

\newcommand{\graphicxextension}{pdf}

\newcommand\myemptypage{
   
    \thispagestyle{empty}
    \addtocounter{page}{-1}
    \newpage
  \null
    %CECI EST UNE PAGE VIDE
       \vfill
    }


\newcommand{\uNA}[2]
{
\newacronym{a#1}{#1}{#2}
}

\newcommand{\coauthor}[4]
{
\noindent (#2) \textbf{#1} - #4 : \textit{#3}
}


\newcommand\UKword[1]{\emph{#1}\xspace}
\newcommand\textUK[1]{\textit{\textbf{#1}}\xspace}
\newcommand\Ucenter[1]{\begin{center}{#1}\end{center}}

\newcommand\head[1]{\textbf{#1}}
\newcommand\tb[1]{\textbf{#1}}
\newcommand{\g}[1]{\og #1 \fg{}}
\newcommand{\uload}[1] {\input{../Latex/#1}}
\newcommand{\ucurcolor}{black}

\newcommand{\mak}[1]{\faicon{\aArrowCircleORight} #1}


%-----------------------------------------------------
%										 	\INFOGitHub 
%---------------------------------------------------


\newcommand{\GITfilename}{https://github.com/edufaction/CYBERDEF101/raw/main/Release/Builder/\jobname.pdf}

%\newcommand{\safeqrcode}[2][]{%
%\qrcode[#1]{\detokenize{#2}} }

\newcommand{\safeqrcode}[2][]{\qrcode[#1]{#2}} 

\newcommand{\INFOGitHub}
{
\begin{center} 
\setstretch{1.5}
 {Vérifiez la disponibilité d'une version plus récente de} \\
 {\ulink{\GITfilename}{\textbf{\jobname.pdf }sur GITHUB CYBERDEF} }   \\ 
{{\huge\ccbyncndeu}}  \\  
{Publication en \href{https://creativecommons.org/licenses/by-nc-nd/2.0/fr/}{Creative Common BY-NC-ND} by eduf@ction}    \\  % Copyright notice
{\safeqrcode[height=4 cm,padding]{\GITfilename}}  
\end{center} 
}

\newcommand{\INFOGitHubS}
{
\noindent{\safeqrcode[height=3 cm]{\GITfilename}} \\
\noindent {\tiny Téléchargez une version à jour} 
}



%-----------------------------------------------------
%										 	\upspicture 
%---------------------------------------------------
\newcommand{\upspicture}[3]
{
\renewcommand{\ucurcolor}{#3}
\resizebox{#2\textwidth}{!}{\input{../#1}}
}


%-----------------------------------------------------
%										 	\uindex 
%---------------------------------------------------


\newcommand{\uindex}[1]
{
#1\index{#1}\xspace
}

\newcommand{\ulindex}[1]{\index{#1}}

%-----------------------------------------------------
%										 	\includer 
%---------------------------------------------------

\newcommand{\includer}[1]
{
\input{../Latex/Sources/EDU/SRC1/Chapters/#1}
}

%-----------------------------------------------------
%										 	\edxdico 
%---------------------------------------------------

\newcommand{\edxdico}[2] {\textbf{#1}:\textit{#2} \index{#1}}

%-----------------------------------------------------
%  command DEF : InTextx
%-----------------------------------------------------
\newcommand\InTexta{no text}
\newcommand\InTextb{no text}
\newcommand\InTextc{no text}
\newcommand\rnc{\renewcommand}


%-----------------------------------------------------
%										 	ENV warningbox
%---------------------------------------------------

\newtcolorbox{warningbox}[2][]
{
  colframe = uMainColor!25,
  colback  = black!05,
  coltitle = uMainColor!20!black,
  fonttitle = \bfseries, 
  title    = #2,
  #1,
}


\newtcolorbox{techworkbox}[2][]{sharp corners,colframe = uMainColor!25,colback  = black!05,coltitle = uMainColor!20!black,fonttitle = \bfseries, title={\textcolor{uMainColor}{\faFile~\small {Sujet de mémoire SECOPS}}} : #2,#1,}


\newtcolorbox{notebox2}[2][]
{
sharp corners, 
colback = uMainColor!5!white, 
colframe = uMainColor!75!black,
fonttitle = \bfseries, 
colbacktitle= uMainColor!85!black,
title=#2,#1
}

\newtcolorbox{notebox}[2][]{colback=uMainColor!5!white,
colframe=uMainColor!75!black,fonttitle=\bfseries,
colbacktitle=uMainColor!85!black,enhanced,
attach boxed title to top right={yshift=+1mm},
title=#2,#1}


\newtcolorbox{toolsboxin}[2][]{sharp corners, colback=uMainColor!0!white,
colframe= black!25,fonttitle=\bfseries,
colbacktitle=black!60,enhanced,
attach boxed title to top left={yshift=+1mm},
title=#2,#1}


\newtcolorbox{techworkin}[2][]{sharp corners, colback=uMainColor!5!white,
colframe= uMainColor!100, fonttitle=\bfseries,
colbacktitle=black!5, enhanced,
 coltitle=black!50,
attach boxed title to top right ={yshift=+1mm},
title=#2,#1}



%-----------------------------------------------------
%										 	ENV ToolsBox
%---------------------------------------------------


\newcommand\tooleditor{No EDITOR in Datas.tex}
\newcommand\toolurl{No URL INSERTED in Datas.tex}
\newcommand\toolclass{No EDITOR in Datas.tex}
\newcommand\toolname{No NAME in Datas.tex}
\newcommand\toolanalyst{No ANALYST in Datas.tex}

\newcommand{\toolsbox}[2]
{
\input{../Latex/Sources/EDU/SRC1/SecComponents/#2.#1/datas.tex} 
\begin{toolsboxin}[sidebyside,righthand width=.7\textwidth]{#1}
\includegraphics[width=0.9\textwidth]{../Latex/Sources/EDU/SRC1/SecComponents/#2.#1/icon.png} \\
\tcblower
\input{../Latex/Sources/EDU/SRC1/SecComponents/#2.#1/description.tex}
\end{toolsboxin}
\noindent \small {\faCogs~Classe : \uppercase{\textbf{#2}}, Site de référence : \textbf{\ulink{\toolurl}{#1}}}
\newline
\small {\faGlobe~Editeur : \textbf{\tooleditor}}
\small {\faEye~Analyste : \textbf{\toolanalyst}}
}

%\renewcommand\toolsclass{SIEM}
%\renewcommand\toolsname{RSA Netwitness}
%\renewcommand\toolseditor{RSA}
%\renewcommand\toolsurl{https://www.rsa.com/fr-fr/products/threat-detection-response}

%-----------------------------------------------------
%										 	ENV  TechWork
%---------------------------------------------------


\newcommand{\techwork}[2]
{
\begin{techworkin} [righthand width=.8\textwidth]{\faCogs~\small {Fiche TECHNO} : #1}#2
\end{techworkin}
}

%-----------------------------------------------------
%										 	\uchap
%---------------------------------------------------
\ifdefined\INTERTITLE
\providecommand{\uchap}[1]
	{
	\begin{LARGE}
	\textbf{\Ucenter{#1 \\ -oOo-}}	
	\end{LARGE}
	}
\else
	\providecommand{\uchap}[1]{}
\fi

%-----------------------------------------------------
%											\uexpand
%---------------------------------------------------
\newcommand {\uexpand}[1]{
\mode<all>\input{../Latex/Sources/EDU/SRC1/Chapters/#1}
}


%-----------------------------------------------------
%											\spicture 
% direct include graphic with Scrivener
%-----------------------------------------------------

\newcommand{\spicture}[3]{
\begin{figure}[!hbtp] %h !hbtp
  \begin{center}
	 \includegraphics[width=0.8\textwidth]{#1}
  \end{center}
\caption{#3 - (slbl-#2)\label{slbl-#2}}
\end{figure}
}

%\StrBefore{#2}{.pdf}

%-----------------------------------------------------
%											\upicture 
%---------------------------------------------------

\ifdefined\PRZMODE

\providecommand{\upicture}[4]{
\framesubtitle{#2}
%\begin{figure}[!h] %h
  \begin{center}
	 \includegraphics[height=#3\textheight, width=#3\textwidth, keepaspectratio ]{\includepicturesubpath{#1.pdf}}
  \end{center}
%  \end{figure}
}

\else

\providecommand{\upicture}[4]{
\begin{figure}[!hbtp] %h !hbtp
  \begin{center}
	 \includegraphics[width=#3\textwidth]{\includepicturesubpath{#1.pdf}}
  \end{center}
\caption{#2 \label{#4}}
\end{figure}
}
	
\fi


\ifdefined\PRZMODE

\providecommand{\uwpicture}[4]{
\framesubtitle{#2}
%\begin{figure}[!h] %h
  \begin{center}
	 \includegraphics[height=0.8\textheight, width=0.8\textwidth, keepaspectratio ]{\includepicturesubpath{#1.pdf}}
  \end{center}
%  \end{figure}
}
\else
\providecommand{\uwpicture}[4]
{
\begin{wrapfigure}{O}{#3\textwidth} % inside / Outside
  \begin{center}
	 \includegraphics[width=#3\textwidth]{\includepicturesubpath{#1.pdf}}
\end{center}
\caption{#2 \label{#4}}
\end{wrapfigure}
}
\fi

%-----------------------------------------------------
%	\picframe for BEAMER (Image included in Beamer Slide)
%---------------------------------------------------

\ifdefined\PRZMODE

\providecommand{\picframe}[4]
{
\begin{frame}
\frametitle<presentation>{#2}
  \begin{center}

   \begin{tikzpicture}[remember picture, overlay]
    \node at (current page.center) {\includegraphics[width=#3\paperwidth, height=#3\paperheight,keepaspectratio]{\includepicturesubpath{#1.pdf}}};
    \end{tikzpicture}

  \end{center}
\end{frame}

}
\else
\providecommand{\picframe}[4]
{
\begin{figure}[!hbtp] % !hbtp ou !h 
\centering 
%{\includegraphics[width=#3\paperwidth, height=#3\paperheight,keepaspectratio]{\includepicturesubpath{#1.pdf}}}
\includegraphics[width=#3\textwidth]{\includepicturesubpath{#1.pdf}}
\caption{#2 \label{#4}}
\end{figure}
}
\fi

%-----------------------------------------------------
%.  for BEAMER (Image included in Beamer Slide)
%	\picframeS Simplified with Name and Label from One Param
%---------------------------------------------------


\ifdefined\PRZMODE

\providecommand{\picframeS}[2]
{
\begin{frame}
\frametitle<presentation>{#2}
  \begin{center}
	 \includegraphics[height=0.95\textheight, width=0.95\textwidth, keepaspectratio ]{\picpath/img-#1.\graphicxextension}
  \end{center}
\end{frame}

}
\else
\providecommand{\picframeS}[2]
{
\begin{figure}[!hbtp] % !hbtp ou !h 
\centering \includegraphics[width=0.9\textwidth]{\picpath/img-#1.\graphicxextension}
\caption{(s) #2 \label{lbl:#1}}
\end{figure}
}

\fi


%{figure}[!h]
%{wrapfigure}{O}{#3\textwidth}
%-----------------------------------------------------
%	FRAME			\wpicframe 
%---------------------------------------------------

\ifdefined\PRZMODE

\providecommand{\wpicframe}[4]
{
\begin{frame}
\frametitle<presentation>{#2}
  \begin{center}
	 \includegraphics[height=0.95\textheight, width=0.95\textwidth, keepaspectratio ]{#1.\graphicxextension}
  \end{center}
\end{frame}

}
\else
\providecommand{\wpicframe}[4]
{
\begin{wrapfigure}{O}{#3\textwidth}
\centering \includegraphics[width=#3\textwidth]{#1.\graphicxextension}
\caption{#2 \label{#4}}
\end{wrapfigure}
}
\fi

%-----------------------------------------------------
%	FRAME										\texframe 
%---------------------------------------------------
\ifdefined\PRZMODE

\providecommand{\texframe}[3]
{
\begin{frame}
\frametitle<presentation>{#1}
\framesubtitle<presentation>{#2}
#3
\end{frame}
}
\else
\providecommand{\texframe}[3]
{
#3
}

\fi

%-----------------------------------------------------
%											\uref 
%---------------------------------------------------
%\newcommand{\uref}[2]{(Voir~#1~\cref{#2} page~\pageref{#2})}

\newcommand{\uref}[2]{(Voir~#1~\ref{#2} page~\pageref{#2})}
%-----------------------------------------------------
%											\ubg
%---------------------------------------------------
\newcommand{\ubg}[1]{
\textbf{\g{\uppercase{#1}}}
}
%-----------------------------------------------------
%											\updfimage 
%---------------------------------------------------

\newcommand{\updfimage}[3]
{
\begin{wrapfigure}{R}{#3\textwidth}
  \begin{center}
	 \resizebox{#3\textwidth}{!}{\includegraphics{\includepicturesubpath{#1.pdf}}}
  \end{center}
\caption{#2}
\end{wrapfigure}
}


%************************PACKAGE***

%\usepackage{tikz}% Required for drawing custom shapes
\usetikzlibrary{shadows}


%-----------------------------------------------------
%											 		\rem 
% insertion d'une remarque avec R
%---------------------------------------------------
\newcommand\rem[1]{%
   \marginpar{%
   \tikzpicture[baseline={(title.base)}]
      \node[inner sep=5pt,text width=4cm,drop shadow={shadow yshift=-5pt,shadow xshift=5pt,uMainColor},fill=white] (box) {\vskip5pt \nointerlineskip #1};
      \node[right=10pt,inner sep=0pt,outer sep=10pt] (title) at (box.north west) {\bfseries\color{uMainColor}Remarque};
      \draw[draw=uMainColor,very thick](title.west)--(box.north west)--(box.south west)--(box.south east)--(box.north east)--(title.east);
      \fill[uMainColor]([yshift=-10pt]box.north west)--+(-5pt,-5pt)--+(0pt,-10pt);
   \endtikzpicture}%
}
%-----------------------------------------------------
%											 \utikzimage 
% insertion d'une image sur fichier tkz.tex
%---------------------------------------------------

\newcommand{\utikzimage}[3]
{
\begin{wrapfigure}{R}{#3\textwidth}
  \begin{center}
	 \resizebox{#3\textwidth}{!}{\input{../#1.tkz.tex}}
  \end{center}
\caption{#2}
\end{wrapfigure}
}

%-----------------------------------------------------
%														\link 
% Url et footnote du lien
%-----------------------------------------------------

% TT every character and hyphenate after it
\def\hyphenateAndTtWholeString #1{\xHyphenate#1$\wholeString\unskip}

\def\xHyphenate#1#2\wholeString {\if#1$%
    \else\transform{#1}%
    \takeTheRest#2\ofTheString\fi}

\def\takeTheRest#1\ofTheString\fi
{\fi \xHyphenate#1\wholeString}

\def\transform#1{\url{#1}\hskip 0pt plus 1pt}

% Define the \urlx command which works like \url, but with line brakes
\def\urlx #1{\href{#1}{\hyphenateAndTtWholeString{#1}}}
  

 %-----------------------------------------------------
%														\link 
%-----------------------------------------------------  
  
  \newcommand{\ulink}[2]
  {\href{#1}
  {#2~\raisebox{-0.2ex}{\faExternalLink}\footnote{\url{#1}}}\xspace}
  
  \newcommand{\wikipedia}[1]
  { \href{https://fr.wikipedia.org/wiki/#1}{\textcolor{grey}{#1}}\raisebox{0.5ex}{\tiny{\textcolor{uMainColor}{\faWikipediaW}}}}
  


 %-----------------------------------------------------
%														\utodo 
%-----------------------------------------------------

  \newcommand{\utodo}{\begin{center}{\noindent \textcolor{cyan}{\faGears~{En cours de rédaction, DRAFT non publiable} \faGears}}\end{center}}
  
    \newcommand{\utocomplete}{\begin{center}{\noindent \textcolor{blue}{\faClipboard~{à compléter, les éléments ne donnent qu'une vue trop réduite ou trop parcellaire du sujet}}}\end{center}}
  

 





%-------------------------------------------------------
\usepackage{\utemplatepath/LessonModel/Beamer/\ubeamermodel.model/edxbeamer}
%-------------------------------------------------------


% a charger après hyperref, pour avoir les guillemets sous beamer
\usepackage{ae}
\usepackage{aecompl}


\usepackage{environ}
\newwrite\audioexport



%\renewcommand{\familydefault}{pag}

%-------------------------------------------------------
% Optional: a subtitle to be displayed on the title slide

\title[\ushorttitle]{\lessonID~-~\utitle}
\subtitle{\tiny\uCourseLongName}

%-------------------------------------------------------
% The author(s) of the presentation:
%  - again first a short version to be displayed at the bottom;
%  - next the full list of authors, which may include contact information;

\author[\uauthorwriter]{ \LARGE{\uauthor} \\ \medskip {\small \url{\umaila} \and \url{\umailb}\\ {\small \url{\usitea}}}}

%-------------------------------------------------------
% The institute:
%  - to start the name of the university as displayed on the top of each slide
%    this can be adjusted such that you can also create a Dutch version
%  - next the institute information as displayed on the title slide

\institute[\uinstituteshort]{{\uinstitute \\ \uchaire}}


%-------------------------------------------------------
% Add a date and possibly the name of the event to the slides
%  - again first a short version to be shown at the bottom of each slide
%  - second the full date and event name for the title slide
\date[\uversion]{\tiny{Publication \INFODistrib ~du \\ \DTMnow}}


\setbeamertemplate{navigation symbols}{%
\insertslidenavigationsymbol
\insertframenavigationsymbol
\insertsubsectionnavigationsymbol
\insertsectionnavigationsymbol
\insertdocnavigationsymbol
\insertbackfindforwardnavigationsymbol
}

%\mode<all>{
%%-------------------------------------------------------------
%               FR CYBERDEF SECOPS COURSE
%
%                 DEFINITION DES Sources/EDU/acronymes

%                      $File :  Sources/EDU/acronymes.tex
%                             2020 eduf@ction
%-------------------------------------------------------------
% Require : glossaries package
%-------------------------------------------------------------
%\newacronym{XXX}{XXX}{XXX}
%-----------------------------------------------------------

%  Organismes

\newacronym{aCNIL}{CNIL}{Commission National Informatique et Liberté}
\newacronym{aOTAN}{OTAN}{Organisation du Traité de l'Atlantique Nord}
\newacronym{aANSSI}{ANSSI}{Agence Nationale de la Sécurité des systèmes d'information}
\newacronym{aCERT}{CERT}{Computer Emergency Response Team}
\newacronym{aAFNOR}{AFNOR}{Association Française de Normalisation}
\newacronym{aSIEM}{SIEM}{System Incident and Event Management}
\newacronym{aSIM}{SIM}{System Incident Management}
\newacronym{aSEM}{SEM}{System Event Management}
\newacronym{aIRP}{IRP}{Incident Response Plan, Plan de Réponse à Incident}
\newacronym{aISO}{ISO}{International Organization for Standardization}
\newacronym{aDICT}{DICT}{Disponibilité, Intégrité, Confidentialité, Traçabilité}
\newacronym{aPDCA}{PDCA}{Plan, Do, Check, Act (Roue de Deming)}
\newacronym{aGDPR}{GDPR}{General Data Protection Regulation ou RGPD Règlement général sur la protection des données}

%-----------------------------------------------------------

% ROLES ET METIERS

\newacronym{aOS}{OS}{Officier de sécurité}
\newacronym{aRSSI}{RSSI}{Responsable de la sécurité des systèmes d'information}
\newacronym{aDPO}{DPO}{Data Protection Officier}
\newacronym{aCIL}{CIL}{ Correspondant Informatique et Liberté}
\newacronym{aSSI}{SSI}{Sécurité des Systèmes d’information}
\newacronym{aGRC}{GRC}{Gouvernance Risques et Conformité}
\newacronym{aDTF}{DTF}{Defend The Flag}
\newacronym{aCTF}{CTF}{Capture The Flag}
\newacronym{aOIV}{OIV}{Opérateur d'Infrastructures Vitales}
\newacronym{aOSE}{OSE}{Opérateur de Services Essentiels}
\newacronym{aLPM}{LPM}{Loi de Programmation Militaire}
\newacronym{aEBIOS}{EBIOS}{Expression des Besoins et Identification des Objectifs de Sécurité}
\newacronym{aDLL}{DLL}{Dynamic Loads Library}
\newacronym{aACID}{ACID}{Authentification, Confidentialité, Intégrité, Disponibilité}

%-------------------------------------------------------
% Sources/EDU/acronymes DATABASE XLS Generation (cf. RESSOURCES DIR)
%-------------------------------------------------------
% DO NOT MODIFY DIRECTLY, USE XLS File (BEGIN)

\newacronym{a2FA}{2FA}{\textit{Two-Factor Authentication} : Authentification à deux facteurs}
\newacronym{aAAA}{AAA}{\textit{Authentication, Authorization, Accounting} : Authentification, autorisation, comptabilisation}
\newacronym{aAC}{AC}{\textit{Access Control} : Accès conditionnel}
\newacronym{aAES}{AES}{\textit{Advanced Encryption Standard} : Algorithme de chiffrement avancé}
\newacronym{aAPT}{APT}{\textit{Advanced Persistent Threat Group} : Groupe de cyber-espionnage avancé}
\newacronym{aASLR}{ASLR}{\textit{Address Space Layout Randomization} : Randomisation de l'espace d'adressage}
\newacronym{aAV}{AV}{\textit{Antivirus Software} : Logiciel antivirus}
\newacronym{aBOT}{BOT}{\textit{Robot (Malicious)} : Robot (malveillant)}
\newacronym{aCC}{C\&C}{\textit{Command And Control} : Commande et contrôle}
\newacronym{aCSRF}{CSRF}{\textit{Cross-Site Request Forgery Attack} : Attaque de confiance de site à site}
\newacronym{aDDOS}{DDOS}{\textit{Distributed Denial Of Service Attack} : Attaque de déni de service distribué}
\newacronym{aDES}{DES}{\textit{Data Encryption Standard} : Algorithme de chiffrement à clé de 64 bits}
\newacronym{aDLP}{DLP}{\textit{Data Loss Prevention} : Gestion de la protection des données sensibles}
\newacronym{aEDR}{EDR}{\textit{Endpoint Detection And Response} : Détection et réponse aux comportements anormaux d'un terminal}
\newacronym{aEK}{EK}{\textit{Exploit Kit} : Kit d'exploit}
\newacronym{aFTP}{FTP}{\textit{File Transfer Protocol} : Protocole de transfert de fichier}
\newacronym{aHTTPS}{HTTPS}{\textit{Protocol For Securing Communications On The Internet} : Protocole de sécurisation HTTP}
\newacronym{aIAM}{IAM}{\textit{Identity And Access Management} : Identité et Gestion des Accés}
\newacronym{aIDS}{IDS}{\textit{Intrusion Detection System} : Système de détection d'intrusion}
\newacronym{aIIS}{IIS}{\textit{Internet Information Services Server}}
\newacronym{aIPS}{IPS}{\textit{Intrusion Prevention System} : Système de prévention d'incident}
\newacronym{aISA}{ISA}{\textit{Internet Security And Acceleration Server} : Serveur de sécurité et d'accélération Internet (Microsoft)}
\newacronym{aLDAP}{LDAP}{\textit{Lightweight Directory Access Protocol} : Protocole de gestion d'annuaire léger}
\newacronym{aMD5}{MD5}{\textit{Message-Digest Algorithm 5}}
\newacronym{aMFA}{MFA}{\textit{Multi-Factor Authentication} : Authentification multi-facteurs}
\newacronym{aNAC}{NAC}{\textit{Network Access Control}}
\newacronym{aOAuth}{OAuth}{\textit{Open Authorization Protocol}}
\newacronym{aPAM}{PAM}{\textit{Privilege Access Management}}
\newacronym{aPKI}{PKI}{\textit{Public Key Infrastructure} : Infrastructure de clefs publiques}
\newacronym{aRSA}{RSA}{\textit{Rivest-Shamir-Adleman Encryption Algorithm}}
\newacronym{aSAML}{SAML}{\textit{Security Assertion Markup Language}}
\newacronym{aSMS}{SMS}{\textit{Short Message Service}}
\newacronym{aSQLI}{SQLI}{\textit{Sql Injection} : Injection de script de commande SQL}
\newacronym{aSSH}{SSH}{\textit{Secure Shell Protocol}}
\newacronym{aSSL}{SSL}{\textit{Protocol For Securing Communications On The Internet} : Protocole de sécurisation des communications sur Internet}
\newacronym{aTKIP}{TKIP}{\textit{Temporal Key Integrity Protocol}}
\newacronym{aTLS}{TLS}{\textit{Transport Layer Security}}
\newacronym{aURL}{URL}{\textit{Uniform Resource Locator}}
\newacronym{aVPN}{VPN}{\textit{Virtual Private Network} : Réseau privé virtuel}
\newacronym{aWAF}{WAF}{\textit{Web Application Firewall}}
\newacronym{aWPA}{WPA}{\textit{Wifi Protected Access}}
\newacronym{aXSS}{XSS}{\textit{Cross-Site Scripting Attack}}

% DO NOT MODIFY (END)
%-------------------------------------------------------




%%-------------------------------------------------------------
%          FR CYBERDEF SECOPS COURSE
%
%                        GLOSSAIRE
%
%                 $File :  Sources/EDU/glossaire.tex
%                    2020 eduf@ction
%-------------------------------------------------------------
% Require : glossaries package
%-------------------------------------------------------------

%\newglossaryentry{xxx}
%{
%name={xxx},
%description={xxx}
%}

\newglossaryentry{gLPM}
{
name={Loi de programmation militaire},
description={Une loi de programmation militaire (LPM), est une loi visant à établir une programmation pluriannuelle des dépenses et donc du budget que l'État français consacre à ses forces armées. Les lois de programmation militaire ont des durées d'application de quatre, cinq ou six ans. La cyberdéfense fait partie des axes importants des lois de programmation}
}

\newglossaryentry{gCTF}
{
name={Capture The Flag},
description={Challenge consistant à exploiter ses connaissance et les vulnérabilités affectant des logiciels  et systèmes numériques de manière à s’introduire sur des réseaux , des ordinateurs, pour récupérer (capturer) des preuves de l’intrusion réussie dénommées G\g{des drapeaux (\textit{Flags} en anglais)}.}
}

\newglossaryentry{gDTF}
{
name={Defend The Flag},
description={Challenge consistant à détecter, résister à des attaquants qui tentent de s'introduire ou sont introduits dans des systèmes numériques de manière à défendre le bien considéré comme le plus sensible convoité par l'attaquant dénommé G\g{drapeau (\textit{Flags} en anglais)}.}
}


\newglossaryentry{gROOTKIT}
{
name={Rootkit},
description={ou simplement « kit » (est un « outil de dissimulation d'activité » ou  « code furtif » qui est une sorte de trousse à outils pour un pirate informatique). Il est constitué  de techniques mises en œuvre par un ou plusieurs logiciels, dont le but est d'obtenir et de pérenniser un accès ceci avec en déployant des techniques de furtivité et de dissimulation.}
}

%-------------------------------------------------------
% GLS DATABASE XLS Generation (cf. RESSOURCES DIR)
%-------------------------------------------------------
% DO NOT MODIFY DIRECTLY, USE XLS File (BEGIN)

\newglossaryentry{gAffiliate}{name={\textit{Affiliates}},description={(Affiliés): Cybercriminels qui mènent des cyberattaques, parfois sophistiquées, par le biais de rançongiciels qui sont mis à disposition par des développeurs ou concepteurs. Les gains générés par les affiliés sont partagés avec les développeurs via un système de commissions.}}

% DO NOT MODIFY (END)
%-------------------------------------------------------


%}

%++++++++++++++++++++++++++++++++
\begin{document}



  %-------------------------------------------------------
  % AudioNote Open File
  %-------------------------------------------------------
\mode<all>{\immediate\openout\audioexport=\jobname.txt}

%================================
\begin{frame}[plain]
 					\begin{figure}
  					 \centering
  					 \begin{columns}[c]
  \begin{column}{0.4\textwidth}
  
    \begin{figure}
        \begin{flushleft}
  			\includegraphics[width=0.9\textwidth]{\utemplatepath/commons.inc/CommonsPictures/\utemplatetheme.theme/maincourselogo.pdf}
  \end{flushleft}
  \end{figure}
      \end{column}
  
  \begin{column}{0.6\textwidth}
    \begin{figure}
        \begin{flushright}
 		    \includegraphics[width=0.7\textwidth]{\utemplatepath/commons.inc/CommonsPictures/\utemplatetheme.theme/maininstitutelogo.pdf}
    \end{flushright}
  \end{figure}
  \end{column}
\end{columns}
   							  
   							  \end{figure}
							
				  				\titlepage
\end{frame}
%================================



%================================
\begin{frame}
	\frametitle{Abstract}
	\begin{techworkin} [righthand width=.8\textwidth]{\faCogs~\small {Hashtags} : \ukeywords}\uabstract
	\end{techworkin}
	\end{frame}
%================================

%================================
\begin{frame}
 \frametitle{Sommaire}
      \begin{multicols}{2}
         \tableofcontents[hideallsubsections]
     \end{multicols}
 \end{frame}
 %================================

%-------------------------------------------------------
%  BEAMER MAIN TEX configured with ubody
%-------------------------------------------------------
\mode<all>{\ubody}

%-------------------------------------------------------
% Last Frame
%-------------------------------------------------------
\begin{frame}
 \frametitle{des questions ?}
  \framesubtitle{contacter \umaila}
  					\begin{figure}
  					 \centering
   						 \includegraphics[width=0.2\textwidth]{\utemplatepath/commons.inc/CommonsPictures/\utemplatetheme.theme/mainlogo.pdf}\\
				  \end{figure}
				 	 \centering \textit{Tous les documents publiés dans le cadre de ce cours sont perfectibles, \\ne pas hésiter à m'envoyer vos remarques !}
 \end{frame}

%-------------------------------------------------------
% Contributions
%-------------------------------------------------------
\begin{frame}
 \frametitle{Contributions}
\begin{columns}
		\begin{column}{0.1\textwidth}

 					\begin{figure}
  						  \centering
   						 \includegraphics[width=0.9\textwidth]{../Latex/Sources/EDU/SRC1/Contribs/contribute/ContribPictures/github.pdf}\\
   						 	 \includegraphics[width=0.9\textwidth]{../Latex/Sources/EDU/SRC1/Contribs/contribute/ContribPictures/overleaf.pdf}
				  \end{figure}
		\end{column}
		\begin{column}{0.9\textwidth}
							%-------------------------------------------------------------
%               FR CYBERDEF SECOPS COURSE
%                            Contributions Rules
%                              2020 eduf@ction
%-------------------------------------------------------------

% \section{Contributions}

% \subsection{Comment contribuer}

Les notes et les présentations sont réalisées sous \hologo{LaTeX}. 

Vous pouvez contribuer au projet du cours CYBERDEF101. Les contributions peuvent se faire sous deux formes :

\begin{itemize}
  \item Corriger, amender, améliorer les notes publiées. A chaque session des modifications et évolutions sont apportées pour tenir compte des corrections de fond et de formes. 
  \item Ajouter, compléter, modifier des parties de notes sur la base de votre lecture du cours et de vos expertises dans chacun des domaines évoqués.
\end{itemize}

 Les fichiers sources sont publiés sur GITHUB dans l'espace : \ulink{https://github.com/edufaction/CYBERDEF101}{(edufaction/CYBERDEF101)}.
 
		\end{column}
\end{columns}
 \end{frame}
 
 %-------------------------------------------------------
% Link to download file and copyright
%-------------------------------------------------------
\begin{frame}
\frametitle{Mises à jour régulières}
\framesubtitle{Eduf@ction~\umaila}
%\INFOGitHub
{\begin{center} 
\setstretch{2}
 {Vérifiez la disponibilité d'une version plus récente de} \\
 {\ulink{\GITfilename}{\textbf{\udocname.pdf }sur GITHUB CYBERDEF} }   \\ 
{{\huge\ccbyncndeu}}  \\  
{\the\year~eduf@ction - Publication en Creative Common BY-NC-ND }    \\  % Copyright notice
{\safeqrcode[padding]{\GITfilename}}  
\end{center} }
\end{frame}


  %-------------------------------------------------------
  % AudioNote Open File
  %-------------------------------------------------------
  \mode<all>{\immediate\closeout\audioexport}


\end{document}




}
\else

\ifdefined\EDTMODE 
	\providecommand{\umainload}{%-------------------------------------------------------------
%               FR CYBERDEF SECOPS COURSE
%                          ARTICLE MAIN FILE
%                             2020 eduf@ction
%-------------------------------------------------------------

\documentclass[10pt,fleqn,twoside]{\utemplatepath/LessonModel/Article/edxstyle} 
%edxstyle
\usepackage{beamerarticle}

\usepackage{ae}
\usepackage[french]{babel}
\usepackage[utf8]{inputenc}
\usepackage[T1]{fontenc}
\usepackage{lmodern}




\usepackage{titletoc} % Required for manipulating the table of contents

\input{\utemplatepath/commons.inc/CommonsPictures/\utemplatetheme.theme/colors.inc}

%-------------------------------------------------------------
%               FR CYBERDEF SECOPS COURSE
%                                  EDX PACKAGES
%.    Commons packages for BOOK, NOTES (.doc) 
%                            and BEAMER (.prz)
%                              2020 eduf@ction
%-------------------------------------------------------------

%\usepackage[hyphens,spaces,obeyspaces]{url}
%\usepackage[hyphens,spaces,obeyspaces]{url}



\PassOptionsToPackage{hyphens,spaces,obeyspaces}{url}

\usepackage[utf8]{inputenc}
\usepackage[french]{babel}
\usepackage[T1]{fontenc}
%\usepackage{markdown}



\usepackage{ccicons}

\usepackage{tikz} % Required for drawing custom shapes
\usetikzlibrary{tikzmark} 
\usetikzlibrary{mindmap,trees, backgrounds}


\usepackage{verbatim}
\usepackage{wrapfig}
\usepackage{xpatch}

\usepackage{hologo}

\usepackage[figure]{totalcount}

\usepackage{fontawesome}

\usepackage[skins]{tcolorbox} % don't forget SKINS option

\usepackage{listings}
\usepackage{upquote}

\lstset{
upquote=true,columns=flexible, basicstyle=\ttfamily,
language=HTML, 
frameround=tttt,
commentstyle=\color{gray},
identifierstyle=\color{blue},
keywordstyle=\color{uMainColor}\bfseries,
xleftmargin=2em,
xrightmargin=2em,
aboveskip=\topsep,
belowskip=\topsep, 
frame=single,
rulecolor=\color{uMainColor},
backgroundcolor=\color{uMainColor!5},
breaklines,
breakindent=1.5em,
showspaces=false,
showstringspaces=false,
showtabs=false,
xleftmargin=2em,
xrightmargin=2em,
aboveskip=1em,
belowskip=1em,
}



\usepackage{graphicx} % Required for including pictures
%\graphicspath{{Pictures/}} % Specifies the directory where pictures are stored

\usepackage{booktabs} % Required for nicer horizontal rules in tables

\usepackage{setspace}
\setstretch{1,2}

\usepackage{csquotes}

\usepackage{fancyhdr}

\usepackage{makeidx} 

\usepackage{xspace} 

\usepackage{floatflt} 

\usepackage{hyperref}
\usepackage{cleveref}

%pdfdisplaydoctitle={\utitle}, Attention le titre ne doit pas contenir de saut de ligne 

\hypersetup{
pdftitle={\uCoursetittle},
pdfauthor={\uauthor},
hidelinks,
backref=true,
pagebackref=true,
hyperindex=true,
pdfcenterwindow = true,
colorlinks=true,
breaklinks=true,
bookmarks=true,
allcolors=uMainColor,
pdfsubject = {\udescription},
pdfcreator = {MACTEX on TEXPAD},
bookmarksopen=true}


% urlcolor=cnam,
% runcolor=uMainColor,
%\pdfstringdefDisableCommands{%
%  \def\\{}%
%  \def\texttt#1{<#1>}%
%}

\usepackage[style=numeric,citestyle=numeric,sorting=nyt,sortcites=true,autopunct=true,hyperref=true,abbreviate=false,backref=true,backend=biber]{biblatex}

%\usepackage{cleveref} old position

\usepackage{multicol}

%--------------------------------------------------------
% pour la génération du QRCODE en page de garde d'article
\usepackage{qrcode} 

%--------------------------------------------------------
% POur la date d'impression en premiere page d'article
\usepackage[useregional]{datetime2}

%--------------------------------------------------------
% Gestion des glossaires et Sources/EDU/acronymes
% Un fichier acroymes.tex centralisé
\usepackage[acronym]{glossaries}

%\makeglossaries
\makenoidxglossaries 

\loadglsentries{Latex/Sources/commons/acronymes.tex}
\loadglsentries{Latex/Sources/commons/glossaire.tex}


%\renewcommand{\bibname}{R\'ef\'erences}

%--------------------------------------------------------
% Pour exclure des "environnement", ici techworkbox par exemple, travaux pour étudiants, non affiché dans le BOOK.
\usepackage{comment}


%--------------------------------------------------------
% pour filigramme sur les pages
\usepackage{eso-pic}

%--------------------------------------------------------
% for \gappto{\UrlBreaks}{\UrlOrds}
% rather than 
%\makeatletter
% \g@addto@macro{\UrlBreaks}{\UrlOrds}
% \makeatother

\usepackage{etoolbox}

\usepackage{breakurl}


\newcommand{\ucite}[1]{\cite{#1}}
%\newcommand{\ucite}[1]{REF BIB #1}
\newcommand{\ugls}[1]{\gls{#1}}
%\newcommand{\ugls}[1]{GLOSSAIRE (#1)}
%\usepackage[Rejne]{fncychap}

%\usepackage{pythontex} % need shell escape

%
%\makeatletter
%\DeclareRobustCommand{\afilename}[1]{%
% \begingroup
%  % \lstname seems to change hyphens into \textendash
%  \def\textendash{-}%
%  \afilename@parse{#1}%
%  \edef\afilename@base{\detokenize\expandafter{\afilename@base}}%
%  \texttt{\afilename@base.\afilename@ext}%
% \endgroup
%}
%\makeatother

%-------------------------------------------------------------
%               FR CYBERDEF SECOPS COURSE
%                                  EDX PACKAGES
%.    Commons COMMANDS for BOOK, NOTES (.doc) 
%                            and BEAMER (.prz)
%                              2020 eduf@ction
%-------------------------------------------------------------

\newcommand{\ucite}[1]{\cite{#1}}
%\newcommand{\ucite}[1]{REF BIB #1}
\newcommand{\ugls}[1]{\gls{#1}}
%\newcommand{\ugls}[1]{GLOSSAIRE (#1)}
\newcommand{\uac}[1]{\acrshort{#1}}



% \newcommand{\ExtractBeforeLastSlash}[1]{\directlua{extract_before_last_slash("#1")}}
% \newcommand{\ExtractAfterLastSlash}[1]{\directlua{extract_after_last_slash("#1")}}
% \newcommand{\RemoveExtension}[1]{\directlua{remove_extension("#1")}}
% \newcommand{\ExtractExtension}[1]{\directlua{extract_extension("#1")}}


\newcommand{\includepicturesubpath}[1]
{#1}

%{\ExtractBeforeLastSlash{#1}/\upicturethemepath\RemoveExtension{\ExtractAfterLastSlash{#1}}.\ExtractExtension{#1}}
    

%---------------------------------------------------
%										 	AUDIO NOTE 
%---------------------------------------------------

\newcommand\audionote[2]{%
  \note{#2}
  \immediate\write\audioexport{==AudioSlide=============}
  \immediate\write\audioexport{#1}
  \immediate\write\audioexport{\insertframenumber}
  \immediate\write\audioexport{---------------------------------}
  \immediate\write\audioexport{#2}
  \immediate\write\audioexport{---------------------------------}
}%

%---------------------------------------------------
%										 	TEXT STYLE SHORTCUT 
%---------------------------------------------------

\newcommand{\tcb}[1]{\textbf{\textcolor{uMainColor}{#1}}}

\newcommand{\voir}[1]{(Schéma~\ref{#1},~page~\pageref{#1}}

\newcommand{\graphicxextension}{pdf}

\newcommand\myemptypage{
   
    \thispagestyle{empty}
    \addtocounter{page}{-1}
    \newpage
  \null
    %CECI EST UNE PAGE VIDE
       \vfill
    }


\newcommand{\uNA}[2]
{
\newacronym{a#1}{#1}{#2}
}

\newcommand{\coauthor}[4]
{
\noindent (#2) \textbf{#1} - #4 : \textit{#3}
}


\newcommand\UKword[1]{\emph{#1}\xspace}
\newcommand\textUK[1]{\textit{\textbf{#1}}\xspace}
\newcommand\Ucenter[1]{\begin{center}{#1}\end{center}}

\newcommand\head[1]{\textbf{#1}}
\newcommand\tb[1]{\textbf{#1}}
\newcommand{\g}[1]{\og #1 \fg{}}
\newcommand{\uload}[1] {\input{../Latex/#1}}
\newcommand{\ucurcolor}{black}

\newcommand{\mak}[1]{\faicon{\aArrowCircleORight} #1}


%-----------------------------------------------------
%										 	\INFOGitHub 
%---------------------------------------------------


\newcommand{\GITfilename}{https://github.com/edufaction/CYBERDEF101/raw/main/Release/Builder/\jobname.pdf}

%\newcommand{\safeqrcode}[2][]{%
%\qrcode[#1]{\detokenize{#2}} }

\newcommand{\safeqrcode}[2][]{\qrcode[#1]{#2}} 

\newcommand{\INFOGitHub}
{
\begin{center} 
\setstretch{1.5}
 {Vérifiez la disponibilité d'une version plus récente de} \\
 {\ulink{\GITfilename}{\textbf{\jobname.pdf }sur GITHUB CYBERDEF} }   \\ 
{{\huge\ccbyncndeu}}  \\  
{Publication en \href{https://creativecommons.org/licenses/by-nc-nd/2.0/fr/}{Creative Common BY-NC-ND} by eduf@ction}    \\  % Copyright notice
{\safeqrcode[height=4 cm,padding]{\GITfilename}}  
\end{center} 
}

\newcommand{\INFOGitHubS}
{
\noindent{\safeqrcode[height=3 cm]{\GITfilename}} \\
\noindent {\tiny Téléchargez une version à jour} 
}



%-----------------------------------------------------
%										 	\upspicture 
%---------------------------------------------------
\newcommand{\upspicture}[3]
{
\renewcommand{\ucurcolor}{#3}
\resizebox{#2\textwidth}{!}{\input{../#1}}
}


%-----------------------------------------------------
%										 	\uindex 
%---------------------------------------------------


\newcommand{\uindex}[1]
{
#1\index{#1}\xspace
}

\newcommand{\ulindex}[1]{\index{#1}}

%-----------------------------------------------------
%										 	\includer 
%---------------------------------------------------

\newcommand{\includer}[1]
{
\input{../Latex/Sources/EDU/SRC1/Chapters/#1}
}

%-----------------------------------------------------
%										 	\edxdico 
%---------------------------------------------------

\newcommand{\edxdico}[2] {\textbf{#1}:\textit{#2} \index{#1}}

%-----------------------------------------------------
%  command DEF : InTextx
%-----------------------------------------------------
\newcommand\InTexta{no text}
\newcommand\InTextb{no text}
\newcommand\InTextc{no text}
\newcommand\rnc{\renewcommand}


%-----------------------------------------------------
%										 	ENV warningbox
%---------------------------------------------------

\newtcolorbox{warningbox}[2][]
{
  colframe = uMainColor!25,
  colback  = black!05,
  coltitle = uMainColor!20!black,
  fonttitle = \bfseries, 
  title    = #2,
  #1,
}


\newtcolorbox{techworkbox}[2][]{sharp corners,colframe = uMainColor!25,colback  = black!05,coltitle = uMainColor!20!black,fonttitle = \bfseries, title={\textcolor{uMainColor}{\faFile~\small {Sujet de mémoire SECOPS}}} : #2,#1,}


\newtcolorbox{notebox2}[2][]
{
sharp corners, 
colback = uMainColor!5!white, 
colframe = uMainColor!75!black,
fonttitle = \bfseries, 
colbacktitle= uMainColor!85!black,
title=#2,#1
}

\newtcolorbox{notebox}[2][]{colback=uMainColor!5!white,
colframe=uMainColor!75!black,fonttitle=\bfseries,
colbacktitle=uMainColor!85!black,enhanced,
attach boxed title to top right={yshift=+1mm},
title=#2,#1}


\newtcolorbox{toolsboxin}[2][]{sharp corners, colback=uMainColor!0!white,
colframe= black!25,fonttitle=\bfseries,
colbacktitle=black!60,enhanced,
attach boxed title to top left={yshift=+1mm},
title=#2,#1}


\newtcolorbox{techworkin}[2][]{sharp corners, colback=uMainColor!5!white,
colframe= uMainColor!100, fonttitle=\bfseries,
colbacktitle=black!5, enhanced,
 coltitle=black!50,
attach boxed title to top right ={yshift=+1mm},
title=#2,#1}



%-----------------------------------------------------
%										 	ENV ToolsBox
%---------------------------------------------------


\newcommand\tooleditor{No EDITOR in Datas.tex}
\newcommand\toolurl{No URL INSERTED in Datas.tex}
\newcommand\toolclass{No EDITOR in Datas.tex}
\newcommand\toolname{No NAME in Datas.tex}
\newcommand\toolanalyst{No ANALYST in Datas.tex}

\newcommand{\toolsbox}[2]
{
\input{../Latex/Sources/EDU/SRC1/SecComponents/#2.#1/datas.tex} 
\begin{toolsboxin}[sidebyside,righthand width=.7\textwidth]{#1}
\includegraphics[width=0.9\textwidth]{../Latex/Sources/EDU/SRC1/SecComponents/#2.#1/icon.png} \\
\tcblower
\input{../Latex/Sources/EDU/SRC1/SecComponents/#2.#1/description.tex}
\end{toolsboxin}
\noindent \small {\faCogs~Classe : \uppercase{\textbf{#2}}, Site de référence : \textbf{\ulink{\toolurl}{#1}}}
\newline
\small {\faGlobe~Editeur : \textbf{\tooleditor}}
\small {\faEye~Analyste : \textbf{\toolanalyst}}
}

%\renewcommand\toolsclass{SIEM}
%\renewcommand\toolsname{RSA Netwitness}
%\renewcommand\toolseditor{RSA}
%\renewcommand\toolsurl{https://www.rsa.com/fr-fr/products/threat-detection-response}

%-----------------------------------------------------
%										 	ENV  TechWork
%---------------------------------------------------


\newcommand{\techwork}[2]
{
\begin{techworkin} [righthand width=.8\textwidth]{\faCogs~\small {Fiche TECHNO} : #1}#2
\end{techworkin}
}

%-----------------------------------------------------
%										 	\uchap
%---------------------------------------------------
\ifdefined\INTERTITLE
\providecommand{\uchap}[1]
	{
	\begin{LARGE}
	\textbf{\Ucenter{#1 \\ -oOo-}}	
	\end{LARGE}
	}
\else
	\providecommand{\uchap}[1]{}
\fi

%-----------------------------------------------------
%											\uexpand
%---------------------------------------------------
\newcommand {\uexpand}[1]{
\mode<all>\input{../Latex/Sources/EDU/SRC1/Chapters/#1}
}


%-----------------------------------------------------
%											\spicture 
% direct include graphic with Scrivener
%-----------------------------------------------------

\newcommand{\spicture}[3]{
\begin{figure}[!hbtp] %h !hbtp
  \begin{center}
	 \includegraphics[width=0.8\textwidth]{#1}
  \end{center}
\caption{#3 - (slbl-#2)\label{slbl-#2}}
\end{figure}
}

%\StrBefore{#2}{.pdf}

%-----------------------------------------------------
%											\upicture 
%---------------------------------------------------

\ifdefined\PRZMODE

\providecommand{\upicture}[4]{
\framesubtitle{#2}
%\begin{figure}[!h] %h
  \begin{center}
	 \includegraphics[height=#3\textheight, width=#3\textwidth, keepaspectratio ]{\includepicturesubpath{#1.pdf}}
  \end{center}
%  \end{figure}
}

\else

\providecommand{\upicture}[4]{
\begin{figure}[!hbtp] %h !hbtp
  \begin{center}
	 \includegraphics[width=#3\textwidth]{\includepicturesubpath{#1.pdf}}
  \end{center}
\caption{#2 \label{#4}}
\end{figure}
}
	
\fi


\ifdefined\PRZMODE

\providecommand{\uwpicture}[4]{
\framesubtitle{#2}
%\begin{figure}[!h] %h
  \begin{center}
	 \includegraphics[height=0.8\textheight, width=0.8\textwidth, keepaspectratio ]{\includepicturesubpath{#1.pdf}}
  \end{center}
%  \end{figure}
}
\else
\providecommand{\uwpicture}[4]
{
\begin{wrapfigure}{O}{#3\textwidth} % inside / Outside
  \begin{center}
	 \includegraphics[width=#3\textwidth]{\includepicturesubpath{#1.pdf}}
\end{center}
\caption{#2 \label{#4}}
\end{wrapfigure}
}
\fi

%-----------------------------------------------------
%	\picframe for BEAMER (Image included in Beamer Slide)
%---------------------------------------------------

\ifdefined\PRZMODE

\providecommand{\picframe}[4]
{
\begin{frame}
\frametitle<presentation>{#2}
  \begin{center}

   \begin{tikzpicture}[remember picture, overlay]
    \node at (current page.center) {\includegraphics[width=#3\paperwidth, height=#3\paperheight,keepaspectratio]{\includepicturesubpath{#1.pdf}}};
    \end{tikzpicture}

  \end{center}
\end{frame}

}
\else
\providecommand{\picframe}[4]
{
\begin{figure}[!hbtp] % !hbtp ou !h 
\centering 
%{\includegraphics[width=#3\paperwidth, height=#3\paperheight,keepaspectratio]{\includepicturesubpath{#1.pdf}}}
\includegraphics[width=#3\textwidth]{\includepicturesubpath{#1.pdf}}
\caption{#2 \label{#4}}
\end{figure}
}
\fi

%-----------------------------------------------------
%.  for BEAMER (Image included in Beamer Slide)
%	\picframeS Simplified with Name and Label from One Param
%---------------------------------------------------


\ifdefined\PRZMODE

\providecommand{\picframeS}[2]
{
\begin{frame}
\frametitle<presentation>{#2}
  \begin{center}
	 \includegraphics[height=0.95\textheight, width=0.95\textwidth, keepaspectratio ]{\picpath/img-#1.\graphicxextension}
  \end{center}
\end{frame}

}
\else
\providecommand{\picframeS}[2]
{
\begin{figure}[!hbtp] % !hbtp ou !h 
\centering \includegraphics[width=0.9\textwidth]{\picpath/img-#1.\graphicxextension}
\caption{(s) #2 \label{lbl:#1}}
\end{figure}
}

\fi


%{figure}[!h]
%{wrapfigure}{O}{#3\textwidth}
%-----------------------------------------------------
%	FRAME			\wpicframe 
%---------------------------------------------------

\ifdefined\PRZMODE

\providecommand{\wpicframe}[4]
{
\begin{frame}
\frametitle<presentation>{#2}
  \begin{center}
	 \includegraphics[height=0.95\textheight, width=0.95\textwidth, keepaspectratio ]{#1.\graphicxextension}
  \end{center}
\end{frame}

}
\else
\providecommand{\wpicframe}[4]
{
\begin{wrapfigure}{O}{#3\textwidth}
\centering \includegraphics[width=#3\textwidth]{#1.\graphicxextension}
\caption{#2 \label{#4}}
\end{wrapfigure}
}
\fi

%-----------------------------------------------------
%	FRAME										\texframe 
%---------------------------------------------------
\ifdefined\PRZMODE

\providecommand{\texframe}[3]
{
\begin{frame}
\frametitle<presentation>{#1}
\framesubtitle<presentation>{#2}
#3
\end{frame}
}
\else
\providecommand{\texframe}[3]
{
#3
}

\fi

%-----------------------------------------------------
%											\uref 
%---------------------------------------------------
%\newcommand{\uref}[2]{(Voir~#1~\cref{#2} page~\pageref{#2})}

\newcommand{\uref}[2]{(Voir~#1~\ref{#2} page~\pageref{#2})}
%-----------------------------------------------------
%											\ubg
%---------------------------------------------------
\newcommand{\ubg}[1]{
\textbf{\g{\uppercase{#1}}}
}
%-----------------------------------------------------
%											\updfimage 
%---------------------------------------------------

\newcommand{\updfimage}[3]
{
\begin{wrapfigure}{R}{#3\textwidth}
  \begin{center}
	 \resizebox{#3\textwidth}{!}{\includegraphics{\includepicturesubpath{#1.pdf}}}
  \end{center}
\caption{#2}
\end{wrapfigure}
}


%************************PACKAGE***

%\usepackage{tikz}% Required for drawing custom shapes
\usetikzlibrary{shadows}


%-----------------------------------------------------
%											 		\rem 
% insertion d'une remarque avec R
%---------------------------------------------------
\newcommand\rem[1]{%
   \marginpar{%
   \tikzpicture[baseline={(title.base)}]
      \node[inner sep=5pt,text width=4cm,drop shadow={shadow yshift=-5pt,shadow xshift=5pt,uMainColor},fill=white] (box) {\vskip5pt \nointerlineskip #1};
      \node[right=10pt,inner sep=0pt,outer sep=10pt] (title) at (box.north west) {\bfseries\color{uMainColor}Remarque};
      \draw[draw=uMainColor,very thick](title.west)--(box.north west)--(box.south west)--(box.south east)--(box.north east)--(title.east);
      \fill[uMainColor]([yshift=-10pt]box.north west)--+(-5pt,-5pt)--+(0pt,-10pt);
   \endtikzpicture}%
}
%-----------------------------------------------------
%											 \utikzimage 
% insertion d'une image sur fichier tkz.tex
%---------------------------------------------------

\newcommand{\utikzimage}[3]
{
\begin{wrapfigure}{R}{#3\textwidth}
  \begin{center}
	 \resizebox{#3\textwidth}{!}{\input{../#1.tkz.tex}}
  \end{center}
\caption{#2}
\end{wrapfigure}
}

%-----------------------------------------------------
%														\link 
% Url et footnote du lien
%-----------------------------------------------------

% TT every character and hyphenate after it
\def\hyphenateAndTtWholeString #1{\xHyphenate#1$\wholeString\unskip}

\def\xHyphenate#1#2\wholeString {\if#1$%
    \else\transform{#1}%
    \takeTheRest#2\ofTheString\fi}

\def\takeTheRest#1\ofTheString\fi
{\fi \xHyphenate#1\wholeString}

\def\transform#1{\url{#1}\hskip 0pt plus 1pt}

% Define the \urlx command which works like \url, but with line brakes
\def\urlx #1{\href{#1}{\hyphenateAndTtWholeString{#1}}}
  

 %-----------------------------------------------------
%														\link 
%-----------------------------------------------------  
  
  \newcommand{\ulink}[2]
  {\href{#1}
  {#2~\raisebox{-0.2ex}{\faExternalLink}\footnote{\url{#1}}}\xspace}
  
  \newcommand{\wikipedia}[1]
  { \href{https://fr.wikipedia.org/wiki/#1}{\textcolor{grey}{#1}}\raisebox{0.5ex}{\tiny{\textcolor{uMainColor}{\faWikipediaW}}}}
  


 %-----------------------------------------------------
%														\utodo 
%-----------------------------------------------------

  \newcommand{\utodo}{\begin{center}{\noindent \textcolor{cyan}{\faGears~{En cours de rédaction, DRAFT non publiable} \faGears}}\end{center}}
  
    \newcommand{\utocomplete}{\begin{center}{\noindent \textcolor{blue}{\faClipboard~{à compléter, les éléments ne donnent qu'une vue trop réduite ou trop parcellaire du sujet}}}\end{center}}
  

 




%-------------------------------------------------------------
%               FR CYBERDEF SECOPS COURSE
%
%                                    EDX STYLE
%
%                              2020 eduf@ction
%-------------------------------------------------------------
\usepackage{enumitem}

\usepackage{parskip}
\setlength{\parindent}{0em}
\setlength{\parskip}{5pt} % 1ex plus 0.5ex minus 0.2ex}

\def\uCoverColor{}

%---------------------------------------------------------------
%	MARGINS
%---------------------------------------------------------------

\usepackage{geometry}

\geometry{
	paper=a4paper, 
%	top=1cm, 
%	bottom=1cm, 
%	left=1cm, 
%	right=1cm, 
%	headheight=20pt, 
%	footskip=1.4cm, 
%	headsep=10pt, 
%	showframe, 
	showcrop 
}

%---------------------------------------------------------------
%	FOOTNOTE (format)
%---------------------------------------------------------------

\makeatletter
\long\def\@makefntextFB#1{%
    \ifx\thefootnote\ftnISsymbol
        \@makefntextORI{#1}%
    \else
        \rule\z@\footnotesep
        \setbox\@tempboxa\hbox{\@thefnmark}%
            \ifdim\wd\@tempboxa>\z@
                \kern2em\llap{\@thefnmark.\kern0.5em}%
            \fi
        \hangindent2em\hangafter\@ne#1
    \fi}
\makeatother

%-----------------------------------------------------
%	COLOR
%-----------------------------------------------------
%\usepackage{xcolor} % Required for specifying colors by name
%\definecolor{uMainColor}{RGB}{160,0,0} 
\colorlet{edxcolorcover}{uCoverColor}
%\definecolor{edxcolorcover}{RGB}{160,0,0} 

%\definecolor{cnam}{RGB}{128, 0, 32} 
                 

  \definecolor{comments}{rgb}{0.7,0,0}    % rouge foncé
  \definecolor{link}{rgb}{0,0.4,0.6}      % ~RoyalBlue de dvips
  \definecolor{url}{rgb}{0.6,0,0}         % rouge-brun
  \definecolor{citation}{rgb}{0,0.5,0}    % vert foncé
  \definecolor{ULlinkcolor}{rgb}{0,0,0.3} % de ulthese.cls
  \definecolor{rouge}{rgb}{0.85,0,0.07}   % rouge bandeau identitaire
  \definecolor{or}{rgb}{1,0.8,0}          % or bandeau identitaire

%-----------------------------------------------------
%	FONTS
%-----------------------------------------------------

%%\usepackage{avant} % Use the Avantgarde font for headings
%\usepackage{times} % Use the Times font for headings
%\usepackage{mathptmx} % Use the Adobe Times Rgman as the default text font together with math symbols from the Sym­bol, Chancery and Com­puter Modern fonts

%%\usepackage{microtype} % Slightly tweak font spacing for aesthetics
%%\usepackage[utf8]{inputenc} % Required for including letters with accents
%\usepackage[T1]{fontenc} % Use 8-bit encoding that has 256 glyphs

% Using AVANT GARDE Family

%\renewcommand{\familydefault}{\sfdefault}

%\fontfamily{pag}\selectfont
%\renewcommand{\familydefault}{pag} % sfdefault pag lmss

%bch         Charter
%lmr         Latin Modern Roman
%lmss        Latin Modern Sans Serif
%lmssq       Latin Modern Sans Serif extended
%lmtt        Latin Modern Typewriter
%lmvtt       Latin Modern Typewriter proportional
%pag         Avant Garde
%pbk         Bookman
%pcr         Courier
%phv         Helvetica
%pnc         New Century Schoolbook
%ppl         Palatino
%ptm         Times
%put         Utopia
 
%\renewcommand{\rmdefault}{pag} % text 
%\renewcommand{\sfdefault}{pag} % titre


%\renewcommand{\rmdefault}{pag} % text 
%\renewcommand{\sfdefault}{phv} % titre


%-----------------------------------------------------
%	BIBLIOGRAPHY AND INDEX
%-----------------------------------------------------


\usepackage{calc} % For simpler calculation - used for spacing the index letter headings correctly

\usepackage{imakeidx}
%\makeindex % Tells LaTeX to create the files required for indexing

\makeindex[columns=2,intoc=true,options={-s \upath/indexstyle.ist}]

%-----------------------------------------------------
%	HEADERS AND FOOTERS
%-----------------------------------------------------


\usepackage{fancyhdr} % Required for header and footer configuration
\pagestyle{fancy} % Enable the custom headers and footers

\ifcsname chapter \endcsname   % verify if command name exists in the CLASS (Book, report, article ...)
% BEGIN IF
\renewcommand{\chaptermark}[1]{\markboth{\sffamily\normalsize\bfseries\chaptername\ \thechapter.\ #1}{}} % Styling for the current chapter in the header
  \else
%nothing
\fi
% END IF

\renewcommand{\sectionmark}[1]{\markright{\sffamily\normalsize\thesection\hspace{5pt}#1}{}} % Styling for the current section in the header

\fancyhf{} % Clear default headers and footers
\fancyhead[LE,RO]{\sffamily\normalsize\thepage} % Styling for the page number in the header
\fancyhead[LO]{\sffamily\normalsize\rightmark} % Print the nearest section name on the left side of odd pages
\fancyhead[RE]{\sffamily\normalsize\leftmark} % Print the current chapter name on the right side of even pages
\fancyfoot[RE,LO]{\sffamily\normalsize\uCoursetittle} % Uncomment to include a footer
\fancyfoot[LE,RO]{\hyperlink{toc}{\sffamily\normalsize\uMinilogo}} % Uncomment to include a footer
% \hyperlink{toc}
\renewcommand{\headrulewidth}{1.0pt} % Thickness of the rule under the header

\fancypagestyle{plain}{% Style for when a plain pagestyle is specified
	\fancyhead{}\renewcommand{\headrulewidth}{0pt}} 

% Removes the header from odd empty pages at the end of chapters
\makeatletter
\renewcommand{\cleardoublepage}{
\clearpage\ifodd\c@page\else
\hbox{}
\vspace*{\fill}
\thispagestyle{empty}
\newpage
\fi}


%----------------------------------------------------------------------------------------
%	MAIN TABLE OF CONTENTS
%----------------------------------------------------------------------------------------
%\usepackage{titletoc} % Required for manipulating the table of contents

\contentsmargin{0cm} % Removes the default margin

% Part text styling (this is mostly taken care of in the PART HEADINGS section of this file)
\titlecontents{part}
	[0cm] % Left indentation
	{\addvspace{20pt}\bfseries} % Spacing and font options for parts
	{}
	{}
	{}

% Chapter text styling
\titlecontents{chapter}
	[1.25cm] % Left indentation
	{\addvspace{12pt}\large\sffamily\bfseries} % Spacing and font options for chapters
	{\color{uMainColor!60}\contentslabel[\Large\thecontentslabel]{1.25cm}\color{uMainColor}} % Formatting of numbered sections of this type
	{\color{uMainColor}} % Formatting of numberless sections of this type
	{\color{uMainColor!60}\normalsize\;\titlerule*[.5pc]{.}\;\thecontentspage} % Formatting of the filler to the right of the heading and the page number

% Section text styling
\titlecontents{section}
	[1.25cm] % Left indentation
	{\addvspace{3pt}\sffamily\bfseries} % Spacing and font options for sections
	{\contentslabel[\thecontentslabel]{1.25cm}} % Formatting of numbered sections of this type
	{} % Formatting of numberless sections of this type
	{\hfill\color{black}\thecontentspage} % Formatting of the filler to the right of the heading and the page number

% Subsection text styling
\titlecontents{subsection}
	[1.25cm] % Left indentation
	{\addvspace{1pt}\sffamily\small} % Spacing and font options for subsections
	{\contentslabel[\thecontentslabel]{1.25cm}} % Formatting of numbered sections of this type
	{} % Formatting of numberless sections of this type
	{\ \titlerule*[.5pc]{.}\;\thecontentspage} % Formatting of the filler to the right of the heading and the page number

%\titlecontents{subsubsection}
%	[1.25cm] % Left indentation
%	{\addvspace{1pt}\sffamily\tiny} % Spacing and font options for subsections
%	{\contentslabel[\thecontentslabel]{1.25cm}} % Formatting of numbered sections of this type
%	{} % Formatting of numberless sections of this type
%	{\ \titlerule*[.5pc]{.}\;\thecontentspage} % Formatting of the filler to the right of the heading and the page number

\titlecontents{subsubsection}
[1.25cm]
  {\addvspace{1pt}\sffamily\tiny}
 {\contentslabel[\thecontentslabel]{1.25cm}}
  {}
  {}

% Figure text styling
\titlecontents{figure}
	[1.25cm] % Left indentation
	{\addvspace{1pt}\sffamily\small} % Spacing and font options for figures
	{\thecontentslabel\hspace*{1em}} % Formatting of numbered sections of this type
	{} % Formatting of numberless sections of this type
	{\ \titlerule*[.5pc]{.}\;\thecontentspage} % Formatting of the filler to the right of the heading and the page number

% Table text styling
\titlecontents{table}
	[1.25cm] % Left indentation
	{\addvspace{1pt}\sffamily\small} % Spacing and font options for tables
	{\thecontentslabel\hspace*{1em}} % Formatting of numbered sections of this type
	{} % Formatting of numberless sections of this type
	{\ \titlerule*[.5pc]{.}\;\thecontentspage} % Formatting of the filler to the right of the heading and the page number

%----------------------------------------------------------------------------------------
%	MINI TABLE OF CONTENTS IN PART HEADS
%----------------------------------------------------------------------------------------

% Chapter text styling
\titlecontents{lchapter}
	[0em] % Left indentation
	{\addvspace{15pt}\large\sffamily\bfseries} % Spacing and font options for chapters
	{\color{uMainColor}\contentslabel[\Large\thecontentslabel]{1.25cm}\color{uMainColor}} % Chapter number
	{}  
	{\color{uMainColor}\normalsize\sffamily\bfseries\;\titlerule*[.5pc]{.}\;\thecontentspage} % Page number

% Section text styling
\titlecontents{lsection}
	[0em] % Left indentation
	{\sffamily\small} % Spacing and font options for sections
	{\contentslabel[\thecontentslabel]{1.25cm}} % Section number
	{}
	{}

% Subsection text styling (note these aren't shown by default, display them by searchings this file for tgcdepth and reading the commented text)
%\titlecontents{lsubsection}
%	[.5em] % Left indentation
%	{\sffamily\footnotesize} % Spacing and font options for subsections
%	{\contentslabel[\thecontentslabel]{1.25cm}}
%	{}
%	{}


\usepackage[]{ccicons}


%----------------------------------------------------------------------------------------
%	THEOREM STYLES
%----------------------------------------------------------------------------------------

\usepackage{amsmath,amsfonts,amssymb,amsthm} % For math equations, theorems, symbols, etc

\newcommand{\intoo}[2]{\mathopen{]}#1\,;#2\mathclose{[}}
\newcommand{\ud}{\mathop{\mathrm{{}d}}\mathopen{}}
\newcommand{\intff}[2]{\mathopen{[}#1\,;#2\mathclose{]}}
\renewcommand{\qedsymbol}{$\blacksquare$}

\ifcsname chapter \endcsname   % verify if command name exists in the CLASS (Book, report, article ...)
	\newtheorem{notation}{Notation}[chapter]
 \else
	\newtheorem{notation}{Notation}[section]
\fi


% Boxed/framed environments
\newtheoremstyle{ocrenumbox}% Theorem style name
{0pt}% Space above
{0pt}% Space below
{\normalfont}% Body font
{}% Indent amount
{\small\bf\sffamily\color{uMainColor}}% Theorem head font
{\;}% Punctuation after theorem head
{0.25em}% Space after theorem head
{\small\sffamily\color{uMainColor}\thmname{#1}\nobreakspace\thmnumber{\@ifnotempty{#1}{}\@upn{#2}}% Theorem text (e.g. Theorem 2.1)
\thmnote{\nobreakspace\the\thm@notefont\sffamily\bfseries\color{black}---\nobreakspace#3.}} % Optional theorem note

\newtheoremstyle{ocrenumbox}% Theorem style name
{0pt}% Space above
{0pt}% Space below
{\normalfont}% Body font
{}% Indent amount
{\small\bf\sffamily\color{uMainColor}}% Theorem head font
{\;}% Punctuation after theorem head
{0.25em}% Space after theorem head
{\small\sffamily\color{uMainColor}\thmname{#1}\nobreakspace\thmnumber{\@ifnotempty{#1}{}\@upn{#2}}% Theorem text (e.g. Theorem 2.1)
\thmnote{\nobreakspace\the\thm@notefont\sffamily\bfseries\color{black}---\nobreakspace#3.}} % Optional theorem note

\newtheoremstyle{blacknumex}% Theorem style name
{5pt}% Space above
{5pt}% Space below
{\normalfont}% Body font
{} % Indent amount
{\small\bf\sffamily}% Theorem head font
{\;}% Punctuation after theorem head
{0.25em}% Space after theorem head
{\small\sffamily{\tiny\ensuremath{\blacksquare}}\nobreakspace\thmname{#1}\nobreakspace\thmnumber{\@ifnotempty{#1}{}\@upn{#2}}% Theorem text (e.g. Theorem 2.1)
\thmnote{\nobreakspace\the\thm@notefont\sffamily\bfseries---\nobreakspace#3.}}% Optional theorem note


\newtheoremstyle{blacknumboxS}% name of the style to be used
{5pt}% measure of space to leave above the theorem. E.g.: 3pt
{5pt}% measure of space to leave below the theorem. E.g.: 3pt
{\normalfont}% name of font to use in the body of the theorem
{}% measure of space to indent
{\small\bf\sffamily}% name of head font
{}% punctuation between head and body
{ }% space after theorem head; " " = normal interword space
{\color{black}\faEye~  \color{uMainColor}\small\sffamily\thmnote{#3 : }}

% Non-boxed/non-framed environments
\newtheoremstyle{ocrenum}% Theorem style name
{5pt}% Space above
{5pt}% Space below
{\normalfont}% Body font
{}% Indent amount
{\small\bf\sffamily\color{uMainColor}}% Theorem head font
{\;}% Punctuation after theorem head
{0.25em}% Space after theorem head
{\small\sffamily\color{uMainColor}\thmname{#1}\nobreakspace\thmnumber{\@ifnotempty{#1}{}\@upn{#2}}% Theorem text (e.g. Theorem 2.1)
\thmnote{\nobreakspace\the\thm@notefont\sffamily\bfseries\color{black}---\nobreakspace#3.}} % Optional theorem note
\makeatother

% Defines the theorem text style for each type of theorem to one of the three styles above



\ifcsname chapter \endcsname

\newcounter{dummy} 

\numberwithin{dummy}{section}
	\theoremstyle{ocrenumbox}
		\newtheorem{theoremeT}[dummy]{Proposition}
		\newtheorem{problem}{Problem}[chapter]
		\newtheorem{exerciseT}{Outillage}[chapter]
	\theoremstyle{blacknumex}
		\newtheorem{exampleT}{Example}[chapter]
	\theoremstyle{blacknumbox}
		\newtheorem{vocabulary}{Vocabulary}[chapter]
		\newtheorem{definitionT}{Definition}[section]
		\newtheorem{corollaryT}[dummy]{Concept}
	\theoremstyle{blacknumboxS}
		\newtheorem{remarqueST}[]{}
		\newtheorem{remarqueT}[dummy]{Remarque}
	\theoremstyle{ocrenum}
	\newtheorem{proposition}[dummy]{Proposition}
  
 \else
 
\newcounter{dummy} 
\numberwithin{dummy}{section}

	\theoremstyle{ocrenumbox}
		\newtheorem{theoremeT}[dummy]{Proposition}
		%\newtheorem{problem}{Problem}[section]
		\newtheorem{exerciseT}{Outillage}[section]
	\theoremstyle{blacknumex}
		\newtheorem{exampleT}{Example}[section]
	\theoremstyle{blacknumbox}
		\newtheorem{vocabulary}{Vocabulary}[section]
		\newtheorem{definitionT}{Definition}[section]
		\newtheorem{corollaryT}[dummy]{concept}
	\theoremstyle{blacknumboxS}
		\newtheorem{remarqueST}[]{}
		\newtheorem{remarqueT}[dummy]{Remarque}
	\theoremstyle{ocrenum}
		\newtheorem{proposition}[dummy]{Proposition}
\fi



%----------------------------------------------------------------------------------------
%	DEFINITION OF COLORED BOXES
%----------------------------------------------------------------------------------------

\RequirePackage[framemethod=default]{mdframed} % Required for creating the theorem, definition, exercise and corollary boxes

% Theorem box
\newmdenv[skipabove=7pt,
skipbelow=7pt,
backgroundcolor=black!5,
linecolor=uMainColor,
innerleftmargin=5pt,
innerrightmargin=5pt,
innertopmargin=5pt,
leftmargin=0cm,
rightmargin=0cm,
innerbottommargin=5pt]{tBox}

% Exercise box	  
\newmdenv[skipabove=7pt,
skipbelow=7pt,
rightline=false,
leftline=true,
topline=false,
bottomline=false,
backgroundcolor=uMainColor!10,
linecolor=uMainColor,
innerleftmargin=5pt,
innerrightmargin=5pt,
innertopmargin=5pt,
innerbottommargin=5pt,
leftmargin=0cm,
rightmargin=0cm,
linewidth=4pt]{eBox}	

% Definition box
\newmdenv[skipabove=7pt,
skipbelow=7pt,
rightline=false,
leftline=true,
topline=false,
bottomline=false,
linecolor=uMainColor,
innerleftmargin=5pt,
innerrightmargin=5pt,
innertopmargin=0pt,
leftmargin=0cm,
rightmargin=0cm,
linewidth=4pt,
innerbottommargin=0pt]{dBox}	

% Corollary box
\newmdenv[skipabove=7pt,
skipbelow=7pt,
rightline=false,
leftline=true,
topline=false,
bottomline=false,
linecolor=gray,
backgroundcolor=black!5,
innerleftmargin=5pt,
innerrightmargin=5pt,
innertopmargin=5pt,
leftmargin=0cm,
rightmargin=0cm,
linewidth=4pt,
innerbottommargin=5pt]{cBox}

% Creates an environment for each type of theorem and assigns it a theorem text style from the "Theorem Styles" section above and a colored box from above
%\newenvironment{theorem}{\begin{tBox}\begin{theoremeT}}{\end{theoremeT}\end{tBox}}
\newenvironment{exercise}{\begin{eBox}\begin{exerciseT}}{\hfill{\color{uMainColor}\tiny\ensuremath{\blacksquare}}\end{exerciseT}\end{eBox}}				  
%\newenvironment{definition}{\begin{dBox}\begin{definitionT}}{\end{definitionT}\end{dBox}}	
%\newenvironment{example}{\begin{exampleT}}{\hfill{\tiny\ensuremath{\blacksquare}}\end{exampleT}}		
%\newenvironment{corollary}{\begin{cBox}\begin{corollaryT}}{\end{corollaryT}\end{cBox}}	
\newenvironment{nota}{\begin{cBox}\begin{remarqueT}}{\end{remarqueT}\end{cBox}}	

%----------------------------------------------------------------------------------------
%	REMARK ENVIRONMENT
%----------------------------------------------------------------------------------------

\newenvironment{remark}{\par\vspace{10pt}\small % Vertical white space above the remark and smaller font size
\begin{list}{}{
\leftmargin=35pt % Indentation on the left
\rightmargin=25pt}\item\ignorespaces % Indentation on the right
\makebox[-2.5pt]{\begin{tikzpicture}[overlay]
\node[draw=uMainColor!60,line width=1pt,circle,fill=uMainColor!25,font=\sffamily\bfseries,inner sep=2pt,outer sep=0pt] at (-15pt,0pt){\textcolor{uMainColor}{i}};\end{tikzpicture}} % Orange R in a circle
\advance\baselineskip -1pt}{\end{list}\vskip5pt} % Tighter line spacing and white space after remark

%----------------------------------------------------------------------------------------
%	SECTION NUMBERING IN THE MARGIN
%----------------------------------------------------------------------------------------

\makeatletter
% \renewcommand{\@seccntformat}[1]{\textcolor{uMainColor}{\csname the#1\endcsname}\quad}   

%\makeatother

%----------------------------------------------------------------------------------------
%	PART HEADINGS
%----------------------------------------------------------------------------------------

% Numbered part in the table of contents
\newcommand{\@mypartnumtocformat}[2]{%
	\setlength\fboxsep{0pt}%
	\noindent\colorbox{uMainColor!20}{\strut\parbox[c][.7cm]{\ecart}{\color{uMainColor!70}\Large\sffamily\bfseries\centering#1}}\hskip\esp\colorbox{uMainColor!40}{\strut\parbox[c][.7cm]{\linewidth-\ecart-\esp}{\Large\sffamily\centering#2}}%
}

% Unnumbered part in the table of contents
\newcommand{\@myparttocformat}[1]{%
	\setlength\fboxsep{0pt}%
	\noindent\colorbox{uMainColor!40}{\strut\parbox[c][.7cm]{\linewidth}{\Large\sffamily\centering#1}}%
}

\newlength\esp
\setlength\esp{4pt}
\newlength\ecart
\setlength\ecart{1.2cm-\esp}
\newcommand{\thepartimage}{}%
\newcommand{\partimage}[1]{\renewcommand{\thepartimage}{#1}}%
\def\@part[#1]#2{%
\ifnum \c@secnumdepth >-2\relax%
\refstepcounter{part}%
\addcontentsline{toc}{part}{\texorpdfstring{\protect\@mypartnumtocformat{\thepart}{#1}}{\partname~\thepart\ ---\ #1}}
\else%
\addcontentsline{toc}{part}{\texorpdfstring{\protect\@myparttocformat{#1}}{#1}}%
\fi%
\startcontents%
\markboth{}{}%
{\thispagestyle{empty}%
\begin{tikzpicture}[remember picture,overlay]%
\node at (current page.north west){\begin{tikzpicture}[remember picture,overlay]%	
\fill[uMainColor!20](0cm,0cm) rectangle (\paperwidth,-\paperheight);
\node[anchor=north] at (4cm,-3.25cm){\color{uMainColor!40}\fontsize{220}{100}\sffamily\bfseries\thepart}; 
\node[anchor=south east] at (\paperwidth-1cm,-\paperheight+1cm){\parbox[t][][t]{8.5cm}{
\printcontents{l}{0}{\setcounter{tocdepth}{1}}% The depth to which the Part mini table of contents displays headings; 0 for chapters only, 1 for chapters and sections and 2 for chapters, sections and subsections
}};
\node[anchor=north east] at (\paperwidth-1.5cm,-3.25cm){\parbox[t][][t]{15cm}{\strut\raggedleft\color{white}\fontsize{30}{30}\sffamily\bfseries#2}};
\end{tikzpicture}};
\end{tikzpicture}}%
\@endpart}
\def\@spart#1{%
\startcontents%
\phantomsection
{\thispagestyle{empty}%
\begin{tikzpicture}[remember picture,overlay]%
\node at (current page.north west){\begin{tikzpicture}[remember picture,overlay]%	
\fill[uMainColor!20](0cm,0cm) rectangle (\paperwidth,-\paperheight);
\node[anchor=north east] at (\paperwidth-1.5cm,-3.25cm){\parbox[t][][t]{15cm}{\strut\raggedleft\color{white}\fontsize{30}{30}\sffamily\bfseries#1}};
\end{tikzpicture}};
\end{tikzpicture}}
\addcontentsline{toc}{part}{\texorpdfstring{%
\setlength\fboxsep{0pt}%
\noindent\protect\colorbox{uMainColor!40}{\strut\protect\parbox[c][.7cm]{\linewidth}{\Large\sffamily\protect\centering #1\quad\mbox{}}}}{#1}}%
\@endpart}
\def\@endpart{\vfil\newpage
\if@twoside
\if@openright
\null
\thispagestyle{empty}%
\newpage
\fi
\fi
\if@tempswa
\twocolumn
\fi}

%----------------------------------------------------------------------------------------
%	CHAPTER HEADINGS
%-------------------------------------------------%--------------------------------------

%
%
%% A switch to conditionally include a picture, implemented by Christian Hupfer
\newif\ifusechapterimage
\usechapterimagetrue
\newcommand{\thechapterimage}{}%
\newcommand{\chapterimage}[1]{\ifusechapterimage\renewcommand{\thechapterimage}{#1}\fi}%
\newcommand{\autodot}{.}

\def\@makechapterhead#1{%
{\parindent \z@ \raggedright \normalfont
\ifnum \c@secnumdepth >\m@ne
\if@mainmatter
\begin{tikzpicture}[remember picture,overlay]
	
\node at (current page.north west)
			{\begin{tikzpicture}[remember picture,overlay]
			\node[anchor=north west,inner sep=0pt] at (0,0) {\ifusechapterimage\includegraphics[width=\paperwidth]{\thechapterimage}\fi};
			\draw[anchor=west] (\Gm@lmargin,-9cm) node [line width=2pt,rounded corners=15pt,draw=uMainColor,fill=white,fill opacity=1,inner sep=15pt]{\strut\makebox[22cm]{}};
			
			\draw[anchor=west] (\Gm@lmargin+.3cm,-9cm) node {\fontsize{20}{30}\selectfont\sffamily\bfseries\color{uMainColor}~#1\strut};
			%{\Huge\sffamily\bfseries\color{uMainColor}\thechapter\autodot~#1\strut};
			\node at (current page.north east)  [xshift=-3.95cm, yshift=-4.05cm, text opacity=1]  {{\color{white}\centering\fontsize{200}{30}\selectfont \bfseries\sffamily\thechapter\strut} };
			\node at (current page.north east)  [xshift=-4cm, yshift=-4cm, text opacity=1]  {{\color{uMainColor}\centering\fontsize{200}{30}\selectfont \bfseries\sffamily\thechapter\strut} };
			
			\end{tikzpicture}};
	\end{tikzpicture}

\else

\begin{tikzpicture}[remember picture,overlay]
\node at (current page.north west)
{\begin{tikzpicture}[remember picture,overlay]
\node[anchor=north west,inner sep=0pt] at (0,0) {\ifusechapterimage\includegraphics[width=\paperwidth]{\thechapterimage}\fi};
\draw[anchor=west] (\Gm@lmargin,-9cm) node [line width=2pt,rounded corners=15pt,draw=uMainColor,fill=white,fill opacity=1,inner sep=15pt]{\strut\makebox[22cm]{}};

\draw[anchor=east] (\Gm@lmargin,-9cm) node [line width=2pt,rounded corners=15pt,draw=uMainColor,fill=white,fill opacity=1,inner sep=15pt]{\strut\makebox[22cm]{}};
\draw[anchor=west] (\Gm@lmargin+.3cm,-9cm) node {\huge\sffamily\bfseries\color{uMainColor}#1\strut};

\end{tikzpicture}};
 
\end{tikzpicture}
\fi\fi\par\vspace*{270\p@}}}

\makeatother
%-------------------------------------------

%\def\@makeschapterhead#1{%
%\begin{tikzpicture}[remember picture,overlay]
%\node at (current page.north west)
%{\begin{tikzpicture}[remember picture,overlay]
%\node[anchor=north west,inner sep=0pt] at (0,0) {\ifusechapterimage\includegraphics[width=\paperwidth]{\thechapterimage}\fi};
%\draw[anchor=west] (\Gm@lmargin,-9cm) node [line width=2pt,rounded corners=15pt,draw=uMainColor,fill=white, fill opacity=0.8,inner sep=15pt]{\strut\makebox[22cm]{}};
%\draw[anchor=west] (\Gm@lmargin+.3cm,-9cm) node {\huge\sffamily\bfseries\color{uMainColor}#1\strut};
%\end{tikzpicture}};
%\end{tikzpicture}
%\par\vspace*{270\p@}}




%----------------------------------------------------------------------------------------
%	LINKS
%----------------------------------------------------------------------------------------


\usepackage{bookmark}
\bookmarksetup{
open,
numbered,
addtohook={%
\ifnum\bookmarkget{level}=0 % chapter
\bookmarksetup{bold}%
\fi
\ifnum\bookmarkget{level}=-1 % part
\bookmarksetup{color=uMainColor,bold}%
\fi
}
}

%----------------------------------------------------------------------------------------
%	FANCY CHAPTER Header
%----------------------------------------------------------------------------------------





%\definecolor{uMainColor}{RGB}{255,128,0}   

%-------------------------------------------------
%	BIBLIOGRAPHY
%-------------------------------------------------

\addbibresource{Latex/Sources/commons/bibliography.bib} % BibTeX bibliography file
\defbibheading{bibempty}{}

%-------------------------------------------------
%	COLORS & BORDERS
%-------------------------------------------------
\setlength{\fboxrule}{0.75pt} % Width of the border around the abstract
%\definecolor{uMainColor}{RGB}{255,128,0} % Color of the article title and section
\definecolor{color2}{RGB}{220,220,220} % Color of the boxes behind the abstract and headings

%-------------------------------------------------
% ITEMS DEFINITION
%-------------------------------------------------

\setlist [itemize,1]{label=\color{uMainColor}\faCaretRight }

\begin{document}


\addtocontents{toc}{\protect\hypertarget{toc}{}}

%-------------------------------------------------
%	ARTICLE INFORMATION
%-------------------------------------------------

\Abstract {\uabstract%===========================
% COURS "INTRO CYBERDEF"
% Abstract général des articles
%===========================


Il fait partie du cours introductif aux fondamentaux de la sécurité des systèmes d'information vue sous deux prismes quelques fois opposés dans la littérature : la gouvernance et la gestion opérationnelle de la sécurité.
Le cours est constitué d'un ensemble de notes de synthèse indépendantes compilées en un document unique, mais édité par chapitre dans le cadre de ce cours.\\
C\edoc ne constitue pas à lui seul le référentiel du cours CYBERDEF101 (SEC101 du Cnam). Il compile des notes de cours mises à disposition de l'auditeur comme support pédagogique partiel à ce cours introductif à la cyberdéfense d'entreprise. }
\JournalInfo{\uJournalInfo} % Journal information
\Archive{Notes de cours éditées le  \DTMnow} % Additional notes (e.g. copyright, DOI, review/research article)
\PaperTitle{\utitle} % Article title
\Authors{\uauthor\textsuperscript{1,}\textsuperscript{2}*} % Authors
\affiliation{\textsuperscript{1}\textit{\uproa}} % Author affiliation
\affiliation{\textsuperscript{2}\textit{\uprob}} % Author affiliation
\affiliation{*\textbf{email}: \umaila\ -- \umailb} % Corresponding author
\Keywords{\ukeywords} % Keywords 
\newcommand{\keywordname}{Mots clefs} % Defines the keywords heading name

%-------------------------------------------------
%	MAKE TITLE
%-------------------------------------------------

%\maketitle

%\begin{tikzpicture}[remember picture,overlay]\node at (current page.north west)[xshift=4cm, yshift=-4cm, text opacity=1.0]{\includegraphics[width=0.1\paperwidth]{\utemplatepath/commons.inc/CommonsPictures/\utemplatetheme.theme/shield-20.pdf}}; 
%\end{tikzpicture} 
%
%
%\AddToShipoutPicture{%
%\AtPageCenter{%
%\begin{tikzpicture}[remember picture,overlay]
%\node [rotate=60,scale=5,text opacity=0.1] at (current page.center) { \textit{\color{uAlertColor}\INFODistrib} };
%\end{tikzpicture} 
%%\begin{tikzpicture}[remember picture,overlay]
%%\node [rotate=60,scale=10,text opacity=0.3]
%%at (current page.center) {moi};
%%\end{tikzpicture}
%}%
%}


%-------------------------------------------------	
%\begin{center}
%	\textcolor{uMainColor}{\textbf{\INFODistrib}} % Distribution information
%\end{center}
%
%\INFOGitHub % Link to download file and copyright
%%-------------------------------------------------

%%-------------------------------------------------------------
%               FR CYBERDEF SECOPS COURSE
%
%                 DEFINITION DES Sources/EDU/acronymes

%                      $File :  Sources/EDU/acronymes.tex
%                             2020 eduf@ction
%-------------------------------------------------------------
% Require : glossaries package
%-------------------------------------------------------------
%\newacronym{XXX}{XXX}{XXX}
%-----------------------------------------------------------

%  Organismes

\newacronym{aCNIL}{CNIL}{Commission National Informatique et Liberté}
\newacronym{aOTAN}{OTAN}{Organisation du Traité de l'Atlantique Nord}
\newacronym{aANSSI}{ANSSI}{Agence Nationale de la Sécurité des systèmes d'information}
\newacronym{aCERT}{CERT}{Computer Emergency Response Team}
\newacronym{aAFNOR}{AFNOR}{Association Française de Normalisation}
\newacronym{aSIEM}{SIEM}{System Incident and Event Management}
\newacronym{aSIM}{SIM}{System Incident Management}
\newacronym{aSEM}{SEM}{System Event Management}
\newacronym{aIRP}{IRP}{Incident Response Plan, Plan de Réponse à Incident}
\newacronym{aISO}{ISO}{International Organization for Standardization}
\newacronym{aDICT}{DICT}{Disponibilité, Intégrité, Confidentialité, Traçabilité}
\newacronym{aPDCA}{PDCA}{Plan, Do, Check, Act (Roue de Deming)}
\newacronym{aGDPR}{GDPR}{General Data Protection Regulation ou RGPD Règlement général sur la protection des données}

%-----------------------------------------------------------

% ROLES ET METIERS

\newacronym{aOS}{OS}{Officier de sécurité}
\newacronym{aRSSI}{RSSI}{Responsable de la sécurité des systèmes d'information}
\newacronym{aDPO}{DPO}{Data Protection Officier}
\newacronym{aCIL}{CIL}{ Correspondant Informatique et Liberté}
\newacronym{aSSI}{SSI}{Sécurité des Systèmes d’information}
\newacronym{aGRC}{GRC}{Gouvernance Risques et Conformité}
\newacronym{aDTF}{DTF}{Defend The Flag}
\newacronym{aCTF}{CTF}{Capture The Flag}
\newacronym{aOIV}{OIV}{Opérateur d'Infrastructures Vitales}
\newacronym{aOSE}{OSE}{Opérateur de Services Essentiels}
\newacronym{aLPM}{LPM}{Loi de Programmation Militaire}
\newacronym{aEBIOS}{EBIOS}{Expression des Besoins et Identification des Objectifs de Sécurité}
\newacronym{aDLL}{DLL}{Dynamic Loads Library}
\newacronym{aACID}{ACID}{Authentification, Confidentialité, Intégrité, Disponibilité}

%-------------------------------------------------------
% Sources/EDU/acronymes DATABASE XLS Generation (cf. RESSOURCES DIR)
%-------------------------------------------------------
% DO NOT MODIFY DIRECTLY, USE XLS File (BEGIN)

\newacronym{a2FA}{2FA}{\textit{Two-Factor Authentication} : Authentification à deux facteurs}
\newacronym{aAAA}{AAA}{\textit{Authentication, Authorization, Accounting} : Authentification, autorisation, comptabilisation}
\newacronym{aAC}{AC}{\textit{Access Control} : Accès conditionnel}
\newacronym{aAES}{AES}{\textit{Advanced Encryption Standard} : Algorithme de chiffrement avancé}
\newacronym{aAPT}{APT}{\textit{Advanced Persistent Threat Group} : Groupe de cyber-espionnage avancé}
\newacronym{aASLR}{ASLR}{\textit{Address Space Layout Randomization} : Randomisation de l'espace d'adressage}
\newacronym{aAV}{AV}{\textit{Antivirus Software} : Logiciel antivirus}
\newacronym{aBOT}{BOT}{\textit{Robot (Malicious)} : Robot (malveillant)}
\newacronym{aCC}{C\&C}{\textit{Command And Control} : Commande et contrôle}
\newacronym{aCSRF}{CSRF}{\textit{Cross-Site Request Forgery Attack} : Attaque de confiance de site à site}
\newacronym{aDDOS}{DDOS}{\textit{Distributed Denial Of Service Attack} : Attaque de déni de service distribué}
\newacronym{aDES}{DES}{\textit{Data Encryption Standard} : Algorithme de chiffrement à clé de 64 bits}
\newacronym{aDLP}{DLP}{\textit{Data Loss Prevention} : Gestion de la protection des données sensibles}
\newacronym{aEDR}{EDR}{\textit{Endpoint Detection And Response} : Détection et réponse aux comportements anormaux d'un terminal}
\newacronym{aEK}{EK}{\textit{Exploit Kit} : Kit d'exploit}
\newacronym{aFTP}{FTP}{\textit{File Transfer Protocol} : Protocole de transfert de fichier}
\newacronym{aHTTPS}{HTTPS}{\textit{Protocol For Securing Communications On The Internet} : Protocole de sécurisation HTTP}
\newacronym{aIAM}{IAM}{\textit{Identity And Access Management} : Identité et Gestion des Accés}
\newacronym{aIDS}{IDS}{\textit{Intrusion Detection System} : Système de détection d'intrusion}
\newacronym{aIIS}{IIS}{\textit{Internet Information Services Server}}
\newacronym{aIPS}{IPS}{\textit{Intrusion Prevention System} : Système de prévention d'incident}
\newacronym{aISA}{ISA}{\textit{Internet Security And Acceleration Server} : Serveur de sécurité et d'accélération Internet (Microsoft)}
\newacronym{aLDAP}{LDAP}{\textit{Lightweight Directory Access Protocol} : Protocole de gestion d'annuaire léger}
\newacronym{aMD5}{MD5}{\textit{Message-Digest Algorithm 5}}
\newacronym{aMFA}{MFA}{\textit{Multi-Factor Authentication} : Authentification multi-facteurs}
\newacronym{aNAC}{NAC}{\textit{Network Access Control}}
\newacronym{aOAuth}{OAuth}{\textit{Open Authorization Protocol}}
\newacronym{aPAM}{PAM}{\textit{Privilege Access Management}}
\newacronym{aPKI}{PKI}{\textit{Public Key Infrastructure} : Infrastructure de clefs publiques}
\newacronym{aRSA}{RSA}{\textit{Rivest-Shamir-Adleman Encryption Algorithm}}
\newacronym{aSAML}{SAML}{\textit{Security Assertion Markup Language}}
\newacronym{aSMS}{SMS}{\textit{Short Message Service}}
\newacronym{aSQLI}{SQLI}{\textit{Sql Injection} : Injection de script de commande SQL}
\newacronym{aSSH}{SSH}{\textit{Secure Shell Protocol}}
\newacronym{aSSL}{SSL}{\textit{Protocol For Securing Communications On The Internet} : Protocole de sécurisation des communications sur Internet}
\newacronym{aTKIP}{TKIP}{\textit{Temporal Key Integrity Protocol}}
\newacronym{aTLS}{TLS}{\textit{Transport Layer Security}}
\newacronym{aURL}{URL}{\textit{Uniform Resource Locator}}
\newacronym{aVPN}{VPN}{\textit{Virtual Private Network} : Réseau privé virtuel}
\newacronym{aWAF}{WAF}{\textit{Web Application Firewall}}
\newacronym{aWPA}{WPA}{\textit{Wifi Protected Access}}
\newacronym{aXSS}{XSS}{\textit{Cross-Site Scripting Attack}}

% DO NOT MODIFY (END)
%-------------------------------------------------------





%-------------------------------------------------
\newpage
\ubody % MAIN BODY defined external
%-------------------------------------------------

	\printbibliography
 \printglossary[type=\acronymtype]
%	\printglossary
%-------------------------------------------------	
\newpage
\Ucontribute % Contribution TEXT
%-------------------------------------------------


%-------------------------------------------------
	\newpage
	\tableofcontents  % TOC
	\iftotalfigures
  			\listoffigures
	\fi
%-------------------------------------------------

\end{document}
}
	\else
	\providecommand{\umainload}{%-------------------------------------------------------------
%               FR CYBERDEF SECOPS COURSE
%                          ARTICLE MAIN FILE
%                             2020 eduf@ction
%-------------------------------------------------------------

\documentclass[10pt,fleqn,twoside]{\utemplatepath/LessonModel/Article/edxstyle} %edxstyle

\usepackage{beamerarticle}

\usepackage{ae}
\usepackage[french]{babel}
\usepackage[utf8]{inputenc}
\usepackage[T1]{fontenc}
\usepackage{lmodern}




\usepackage{titletoc} % Required for manipulating the table of contents

\input{\utemplatepath/commons.inc/CommonsPictures/\utemplatetheme.theme/colors.inc}

%-------------------------------------------------------------
%               FR CYBERDEF SECOPS COURSE
%                                  EDX PACKAGES
%.    Commons packages for BOOK, NOTES (.doc) 
%                            and BEAMER (.prz)
%                              2020 eduf@ction
%-------------------------------------------------------------

%\usepackage[hyphens,spaces,obeyspaces]{url}
%\usepackage[hyphens,spaces,obeyspaces]{url}



\PassOptionsToPackage{hyphens,spaces,obeyspaces}{url}

\usepackage[utf8]{inputenc}
\usepackage[french]{babel}
\usepackage[T1]{fontenc}
%\usepackage{markdown}



\usepackage{ccicons}

\usepackage{tikz} % Required for drawing custom shapes
\usetikzlibrary{tikzmark} 
\usetikzlibrary{mindmap,trees, backgrounds}


\usepackage{verbatim}
\usepackage{wrapfig}
\usepackage{xpatch}

\usepackage{hologo}

\usepackage[figure]{totalcount}

\usepackage{fontawesome}

\usepackage[skins]{tcolorbox} % don't forget SKINS option

\usepackage{listings}
\usepackage{upquote}

\lstset{
upquote=true,columns=flexible, basicstyle=\ttfamily,
language=HTML, 
frameround=tttt,
commentstyle=\color{gray},
identifierstyle=\color{blue},
keywordstyle=\color{uMainColor}\bfseries,
xleftmargin=2em,
xrightmargin=2em,
aboveskip=\topsep,
belowskip=\topsep, 
frame=single,
rulecolor=\color{uMainColor},
backgroundcolor=\color{uMainColor!5},
breaklines,
breakindent=1.5em,
showspaces=false,
showstringspaces=false,
showtabs=false,
xleftmargin=2em,
xrightmargin=2em,
aboveskip=1em,
belowskip=1em,
}



\usepackage{graphicx} % Required for including pictures
%\graphicspath{{Pictures/}} % Specifies the directory where pictures are stored

\usepackage{booktabs} % Required for nicer horizontal rules in tables

\usepackage{setspace}
\setstretch{1,2}

\usepackage{csquotes}

\usepackage{fancyhdr}

\usepackage{makeidx} 

\usepackage{xspace} 

\usepackage{floatflt} 

\usepackage{hyperref}
\usepackage{cleveref}

%pdfdisplaydoctitle={\utitle}, Attention le titre ne doit pas contenir de saut de ligne 

\hypersetup{
pdftitle={\uCoursetittle},
pdfauthor={\uauthor},
hidelinks,
backref=true,
pagebackref=true,
hyperindex=true,
pdfcenterwindow = true,
colorlinks=true,
breaklinks=true,
bookmarks=true,
allcolors=uMainColor,
pdfsubject = {\udescription},
pdfcreator = {MACTEX on TEXPAD},
bookmarksopen=true}


% urlcolor=cnam,
% runcolor=uMainColor,
%\pdfstringdefDisableCommands{%
%  \def\\{}%
%  \def\texttt#1{<#1>}%
%}

\usepackage[style=numeric,citestyle=numeric,sorting=nyt,sortcites=true,autopunct=true,hyperref=true,abbreviate=false,backref=true,backend=biber]{biblatex}

%\usepackage{cleveref} old position

\usepackage{multicol}

%--------------------------------------------------------
% pour la génération du QRCODE en page de garde d'article
\usepackage{qrcode} 

%--------------------------------------------------------
% POur la date d'impression en premiere page d'article
\usepackage[useregional]{datetime2}

%--------------------------------------------------------
% Gestion des glossaires et Sources/EDU/acronymes
% Un fichier acroymes.tex centralisé
\usepackage[acronym]{glossaries}

%\makeglossaries
\makenoidxglossaries 

\loadglsentries{Latex/Sources/commons/acronymes.tex}
\loadglsentries{Latex/Sources/commons/glossaire.tex}


%\renewcommand{\bibname}{R\'ef\'erences}

%--------------------------------------------------------
% Pour exclure des "environnement", ici techworkbox par exemple, travaux pour étudiants, non affiché dans le BOOK.
\usepackage{comment}


%--------------------------------------------------------
% pour filigramme sur les pages
\usepackage{eso-pic}

%--------------------------------------------------------
% for \gappto{\UrlBreaks}{\UrlOrds}
% rather than 
%\makeatletter
% \g@addto@macro{\UrlBreaks}{\UrlOrds}
% \makeatother

\usepackage{etoolbox}

\usepackage{breakurl}


\newcommand{\ucite}[1]{\cite{#1}}
%\newcommand{\ucite}[1]{REF BIB #1}
\newcommand{\ugls}[1]{\gls{#1}}
%\newcommand{\ugls}[1]{GLOSSAIRE (#1)}
%\usepackage[Rejne]{fncychap}

%\usepackage{pythontex} % need shell escape

%
%\makeatletter
%\DeclareRobustCommand{\afilename}[1]{%
% \begingroup
%  % \lstname seems to change hyphens into \textendash
%  \def\textendash{-}%
%  \afilename@parse{#1}%
%  \edef\afilename@base{\detokenize\expandafter{\afilename@base}}%
%  \texttt{\afilename@base.\afilename@ext}%
% \endgroup
%}
%\makeatother

%-------------------------------------------------------------
%               FR CYBERDEF SECOPS COURSE
%                                  EDX PACKAGES
%.    Commons COMMANDS for BOOK, NOTES (.doc) 
%                            and BEAMER (.prz)
%                              2020 eduf@ction
%-------------------------------------------------------------

\newcommand{\ucite}[1]{\cite{#1}}
%\newcommand{\ucite}[1]{REF BIB #1}
\newcommand{\ugls}[1]{\gls{#1}}
%\newcommand{\ugls}[1]{GLOSSAIRE (#1)}
\newcommand{\uac}[1]{\acrshort{#1}}



% \newcommand{\ExtractBeforeLastSlash}[1]{\directlua{extract_before_last_slash("#1")}}
% \newcommand{\ExtractAfterLastSlash}[1]{\directlua{extract_after_last_slash("#1")}}
% \newcommand{\RemoveExtension}[1]{\directlua{remove_extension("#1")}}
% \newcommand{\ExtractExtension}[1]{\directlua{extract_extension("#1")}}


\newcommand{\includepicturesubpath}[1]
{#1}

%{\ExtractBeforeLastSlash{#1}/\upicturethemepath\RemoveExtension{\ExtractAfterLastSlash{#1}}.\ExtractExtension{#1}}
    

%---------------------------------------------------
%										 	AUDIO NOTE 
%---------------------------------------------------

\newcommand\audionote[2]{%
  \note{#2}
  \immediate\write\audioexport{==AudioSlide=============}
  \immediate\write\audioexport{#1}
  \immediate\write\audioexport{\insertframenumber}
  \immediate\write\audioexport{---------------------------------}
  \immediate\write\audioexport{#2}
  \immediate\write\audioexport{---------------------------------}
}%

%---------------------------------------------------
%										 	TEXT STYLE SHORTCUT 
%---------------------------------------------------

\newcommand{\tcb}[1]{\textbf{\textcolor{uMainColor}{#1}}}

\newcommand{\voir}[1]{(Schéma~\ref{#1},~page~\pageref{#1}}

\newcommand{\graphicxextension}{pdf}

\newcommand\myemptypage{
   
    \thispagestyle{empty}
    \addtocounter{page}{-1}
    \newpage
  \null
    %CECI EST UNE PAGE VIDE
       \vfill
    }


\newcommand{\uNA}[2]
{
\newacronym{a#1}{#1}{#2}
}

\newcommand{\coauthor}[4]
{
\noindent (#2) \textbf{#1} - #4 : \textit{#3}
}


\newcommand\UKword[1]{\emph{#1}\xspace}
\newcommand\textUK[1]{\textit{\textbf{#1}}\xspace}
\newcommand\Ucenter[1]{\begin{center}{#1}\end{center}}

\newcommand\head[1]{\textbf{#1}}
\newcommand\tb[1]{\textbf{#1}}
\newcommand{\g}[1]{\og #1 \fg{}}
\newcommand{\uload}[1] {\input{../Latex/#1}}
\newcommand{\ucurcolor}{black}

\newcommand{\mak}[1]{\faicon{\aArrowCircleORight} #1}


%-----------------------------------------------------
%										 	\INFOGitHub 
%---------------------------------------------------


\newcommand{\GITfilename}{https://github.com/edufaction/CYBERDEF101/raw/main/Release/Builder/\jobname.pdf}

%\newcommand{\safeqrcode}[2][]{%
%\qrcode[#1]{\detokenize{#2}} }

\newcommand{\safeqrcode}[2][]{\qrcode[#1]{#2}} 

\newcommand{\INFOGitHub}
{
\begin{center} 
\setstretch{1.5}
 {Vérifiez la disponibilité d'une version plus récente de} \\
 {\ulink{\GITfilename}{\textbf{\jobname.pdf }sur GITHUB CYBERDEF} }   \\ 
{{\huge\ccbyncndeu}}  \\  
{Publication en \href{https://creativecommons.org/licenses/by-nc-nd/2.0/fr/}{Creative Common BY-NC-ND} by eduf@ction}    \\  % Copyright notice
{\safeqrcode[height=4 cm,padding]{\GITfilename}}  
\end{center} 
}

\newcommand{\INFOGitHubS}
{
\noindent{\safeqrcode[height=3 cm]{\GITfilename}} \\
\noindent {\tiny Téléchargez une version à jour} 
}



%-----------------------------------------------------
%										 	\upspicture 
%---------------------------------------------------
\newcommand{\upspicture}[3]
{
\renewcommand{\ucurcolor}{#3}
\resizebox{#2\textwidth}{!}{\input{../#1}}
}


%-----------------------------------------------------
%										 	\uindex 
%---------------------------------------------------


\newcommand{\uindex}[1]
{
#1\index{#1}\xspace
}

\newcommand{\ulindex}[1]{\index{#1}}

%-----------------------------------------------------
%										 	\includer 
%---------------------------------------------------

\newcommand{\includer}[1]
{
\input{../Latex/Sources/EDU/SRC1/Chapters/#1}
}

%-----------------------------------------------------
%										 	\edxdico 
%---------------------------------------------------

\newcommand{\edxdico}[2] {\textbf{#1}:\textit{#2} \index{#1}}

%-----------------------------------------------------
%  command DEF : InTextx
%-----------------------------------------------------
\newcommand\InTexta{no text}
\newcommand\InTextb{no text}
\newcommand\InTextc{no text}
\newcommand\rnc{\renewcommand}


%-----------------------------------------------------
%										 	ENV warningbox
%---------------------------------------------------

\newtcolorbox{warningbox}[2][]
{
  colframe = uMainColor!25,
  colback  = black!05,
  coltitle = uMainColor!20!black,
  fonttitle = \bfseries, 
  title    = #2,
  #1,
}


\newtcolorbox{techworkbox}[2][]{sharp corners,colframe = uMainColor!25,colback  = black!05,coltitle = uMainColor!20!black,fonttitle = \bfseries, title={\textcolor{uMainColor}{\faFile~\small {Sujet de mémoire SECOPS}}} : #2,#1,}


\newtcolorbox{notebox2}[2][]
{
sharp corners, 
colback = uMainColor!5!white, 
colframe = uMainColor!75!black,
fonttitle = \bfseries, 
colbacktitle= uMainColor!85!black,
title=#2,#1
}

\newtcolorbox{notebox}[2][]{colback=uMainColor!5!white,
colframe=uMainColor!75!black,fonttitle=\bfseries,
colbacktitle=uMainColor!85!black,enhanced,
attach boxed title to top right={yshift=+1mm},
title=#2,#1}


\newtcolorbox{toolsboxin}[2][]{sharp corners, colback=uMainColor!0!white,
colframe= black!25,fonttitle=\bfseries,
colbacktitle=black!60,enhanced,
attach boxed title to top left={yshift=+1mm},
title=#2,#1}


\newtcolorbox{techworkin}[2][]{sharp corners, colback=uMainColor!5!white,
colframe= uMainColor!100, fonttitle=\bfseries,
colbacktitle=black!5, enhanced,
 coltitle=black!50,
attach boxed title to top right ={yshift=+1mm},
title=#2,#1}



%-----------------------------------------------------
%										 	ENV ToolsBox
%---------------------------------------------------


\newcommand\tooleditor{No EDITOR in Datas.tex}
\newcommand\toolurl{No URL INSERTED in Datas.tex}
\newcommand\toolclass{No EDITOR in Datas.tex}
\newcommand\toolname{No NAME in Datas.tex}
\newcommand\toolanalyst{No ANALYST in Datas.tex}

\newcommand{\toolsbox}[2]
{
\input{../Latex/Sources/EDU/SRC1/SecComponents/#2.#1/datas.tex} 
\begin{toolsboxin}[sidebyside,righthand width=.7\textwidth]{#1}
\includegraphics[width=0.9\textwidth]{../Latex/Sources/EDU/SRC1/SecComponents/#2.#1/icon.png} \\
\tcblower
\input{../Latex/Sources/EDU/SRC1/SecComponents/#2.#1/description.tex}
\end{toolsboxin}
\noindent \small {\faCogs~Classe : \uppercase{\textbf{#2}}, Site de référence : \textbf{\ulink{\toolurl}{#1}}}
\newline
\small {\faGlobe~Editeur : \textbf{\tooleditor}}
\small {\faEye~Analyste : \textbf{\toolanalyst}}
}

%\renewcommand\toolsclass{SIEM}
%\renewcommand\toolsname{RSA Netwitness}
%\renewcommand\toolseditor{RSA}
%\renewcommand\toolsurl{https://www.rsa.com/fr-fr/products/threat-detection-response}

%-----------------------------------------------------
%										 	ENV  TechWork
%---------------------------------------------------


\newcommand{\techwork}[2]
{
\begin{techworkin} [righthand width=.8\textwidth]{\faCogs~\small {Fiche TECHNO} : #1}#2
\end{techworkin}
}

%-----------------------------------------------------
%										 	\uchap
%---------------------------------------------------
\ifdefined\INTERTITLE
\providecommand{\uchap}[1]
	{
	\begin{LARGE}
	\textbf{\Ucenter{#1 \\ -oOo-}}	
	\end{LARGE}
	}
\else
	\providecommand{\uchap}[1]{}
\fi

%-----------------------------------------------------
%											\uexpand
%---------------------------------------------------
\newcommand {\uexpand}[1]{
\mode<all>\input{../Latex/Sources/EDU/SRC1/Chapters/#1}
}


%-----------------------------------------------------
%											\spicture 
% direct include graphic with Scrivener
%-----------------------------------------------------

\newcommand{\spicture}[3]{
\begin{figure}[!hbtp] %h !hbtp
  \begin{center}
	 \includegraphics[width=0.8\textwidth]{#1}
  \end{center}
\caption{#3 - (slbl-#2)\label{slbl-#2}}
\end{figure}
}

%\StrBefore{#2}{.pdf}

%-----------------------------------------------------
%											\upicture 
%---------------------------------------------------

\ifdefined\PRZMODE

\providecommand{\upicture}[4]{
\framesubtitle{#2}
%\begin{figure}[!h] %h
  \begin{center}
	 \includegraphics[height=#3\textheight, width=#3\textwidth, keepaspectratio ]{\includepicturesubpath{#1.pdf}}
  \end{center}
%  \end{figure}
}

\else

\providecommand{\upicture}[4]{
\begin{figure}[!hbtp] %h !hbtp
  \begin{center}
	 \includegraphics[width=#3\textwidth]{\includepicturesubpath{#1.pdf}}
  \end{center}
\caption{#2 \label{#4}}
\end{figure}
}
	
\fi


\ifdefined\PRZMODE

\providecommand{\uwpicture}[4]{
\framesubtitle{#2}
%\begin{figure}[!h] %h
  \begin{center}
	 \includegraphics[height=0.8\textheight, width=0.8\textwidth, keepaspectratio ]{\includepicturesubpath{#1.pdf}}
  \end{center}
%  \end{figure}
}
\else
\providecommand{\uwpicture}[4]
{
\begin{wrapfigure}{O}{#3\textwidth} % inside / Outside
  \begin{center}
	 \includegraphics[width=#3\textwidth]{\includepicturesubpath{#1.pdf}}
\end{center}
\caption{#2 \label{#4}}
\end{wrapfigure}
}
\fi

%-----------------------------------------------------
%	\picframe for BEAMER (Image included in Beamer Slide)
%---------------------------------------------------

\ifdefined\PRZMODE

\providecommand{\picframe}[4]
{
\begin{frame}
\frametitle<presentation>{#2}
  \begin{center}

   \begin{tikzpicture}[remember picture, overlay]
    \node at (current page.center) {\includegraphics[width=#3\paperwidth, height=#3\paperheight,keepaspectratio]{\includepicturesubpath{#1.pdf}}};
    \end{tikzpicture}

  \end{center}
\end{frame}

}
\else
\providecommand{\picframe}[4]
{
\begin{figure}[!hbtp] % !hbtp ou !h 
\centering 
%{\includegraphics[width=#3\paperwidth, height=#3\paperheight,keepaspectratio]{\includepicturesubpath{#1.pdf}}}
\includegraphics[width=#3\textwidth]{\includepicturesubpath{#1.pdf}}
\caption{#2 \label{#4}}
\end{figure}
}
\fi

%-----------------------------------------------------
%.  for BEAMER (Image included in Beamer Slide)
%	\picframeS Simplified with Name and Label from One Param
%---------------------------------------------------


\ifdefined\PRZMODE

\providecommand{\picframeS}[2]
{
\begin{frame}
\frametitle<presentation>{#2}
  \begin{center}
	 \includegraphics[height=0.95\textheight, width=0.95\textwidth, keepaspectratio ]{\picpath/img-#1.\graphicxextension}
  \end{center}
\end{frame}

}
\else
\providecommand{\picframeS}[2]
{
\begin{figure}[!hbtp] % !hbtp ou !h 
\centering \includegraphics[width=0.9\textwidth]{\picpath/img-#1.\graphicxextension}
\caption{(s) #2 \label{lbl:#1}}
\end{figure}
}

\fi


%{figure}[!h]
%{wrapfigure}{O}{#3\textwidth}
%-----------------------------------------------------
%	FRAME			\wpicframe 
%---------------------------------------------------

\ifdefined\PRZMODE

\providecommand{\wpicframe}[4]
{
\begin{frame}
\frametitle<presentation>{#2}
  \begin{center}
	 \includegraphics[height=0.95\textheight, width=0.95\textwidth, keepaspectratio ]{#1.\graphicxextension}
  \end{center}
\end{frame}

}
\else
\providecommand{\wpicframe}[4]
{
\begin{wrapfigure}{O}{#3\textwidth}
\centering \includegraphics[width=#3\textwidth]{#1.\graphicxextension}
\caption{#2 \label{#4}}
\end{wrapfigure}
}
\fi

%-----------------------------------------------------
%	FRAME										\texframe 
%---------------------------------------------------
\ifdefined\PRZMODE

\providecommand{\texframe}[3]
{
\begin{frame}
\frametitle<presentation>{#1}
\framesubtitle<presentation>{#2}
#3
\end{frame}
}
\else
\providecommand{\texframe}[3]
{
#3
}

\fi

%-----------------------------------------------------
%											\uref 
%---------------------------------------------------
%\newcommand{\uref}[2]{(Voir~#1~\cref{#2} page~\pageref{#2})}

\newcommand{\uref}[2]{(Voir~#1~\ref{#2} page~\pageref{#2})}
%-----------------------------------------------------
%											\ubg
%---------------------------------------------------
\newcommand{\ubg}[1]{
\textbf{\g{\uppercase{#1}}}
}
%-----------------------------------------------------
%											\updfimage 
%---------------------------------------------------

\newcommand{\updfimage}[3]
{
\begin{wrapfigure}{R}{#3\textwidth}
  \begin{center}
	 \resizebox{#3\textwidth}{!}{\includegraphics{\includepicturesubpath{#1.pdf}}}
  \end{center}
\caption{#2}
\end{wrapfigure}
}


%************************PACKAGE***

%\usepackage{tikz}% Required for drawing custom shapes
\usetikzlibrary{shadows}


%-----------------------------------------------------
%											 		\rem 
% insertion d'une remarque avec R
%---------------------------------------------------
\newcommand\rem[1]{%
   \marginpar{%
   \tikzpicture[baseline={(title.base)}]
      \node[inner sep=5pt,text width=4cm,drop shadow={shadow yshift=-5pt,shadow xshift=5pt,uMainColor},fill=white] (box) {\vskip5pt \nointerlineskip #1};
      \node[right=10pt,inner sep=0pt,outer sep=10pt] (title) at (box.north west) {\bfseries\color{uMainColor}Remarque};
      \draw[draw=uMainColor,very thick](title.west)--(box.north west)--(box.south west)--(box.south east)--(box.north east)--(title.east);
      \fill[uMainColor]([yshift=-10pt]box.north west)--+(-5pt,-5pt)--+(0pt,-10pt);
   \endtikzpicture}%
}
%-----------------------------------------------------
%											 \utikzimage 
% insertion d'une image sur fichier tkz.tex
%---------------------------------------------------

\newcommand{\utikzimage}[3]
{
\begin{wrapfigure}{R}{#3\textwidth}
  \begin{center}
	 \resizebox{#3\textwidth}{!}{\input{../#1.tkz.tex}}
  \end{center}
\caption{#2}
\end{wrapfigure}
}

%-----------------------------------------------------
%														\link 
% Url et footnote du lien
%-----------------------------------------------------

% TT every character and hyphenate after it
\def\hyphenateAndTtWholeString #1{\xHyphenate#1$\wholeString\unskip}

\def\xHyphenate#1#2\wholeString {\if#1$%
    \else\transform{#1}%
    \takeTheRest#2\ofTheString\fi}

\def\takeTheRest#1\ofTheString\fi
{\fi \xHyphenate#1\wholeString}

\def\transform#1{\url{#1}\hskip 0pt plus 1pt}

% Define the \urlx command which works like \url, but with line brakes
\def\urlx #1{\href{#1}{\hyphenateAndTtWholeString{#1}}}
  

 %-----------------------------------------------------
%														\link 
%-----------------------------------------------------  
  
  \newcommand{\ulink}[2]
  {\href{#1}
  {#2~\raisebox{-0.2ex}{\faExternalLink}\footnote{\url{#1}}}\xspace}
  
  \newcommand{\wikipedia}[1]
  { \href{https://fr.wikipedia.org/wiki/#1}{\textcolor{grey}{#1}}\raisebox{0.5ex}{\tiny{\textcolor{uMainColor}{\faWikipediaW}}}}
  


 %-----------------------------------------------------
%														\utodo 
%-----------------------------------------------------

  \newcommand{\utodo}{\begin{center}{\noindent \textcolor{cyan}{\faGears~{En cours de rédaction, DRAFT non publiable} \faGears}}\end{center}}
  
    \newcommand{\utocomplete}{\begin{center}{\noindent \textcolor{blue}{\faClipboard~{à compléter, les éléments ne donnent qu'une vue trop réduite ou trop parcellaire du sujet}}}\end{center}}
  

 




%-------------------------------------------------------------
%               FR CYBERDEF SECOPS COURSE
%
%                                    EDX STYLE
%
%                              2020 eduf@ction
%-------------------------------------------------------------
\usepackage{enumitem}

\usepackage{parskip}
\setlength{\parindent}{0em}
\setlength{\parskip}{5pt} % 1ex plus 0.5ex minus 0.2ex}

\def\uCoverColor{}

%---------------------------------------------------------------
%	MARGINS
%---------------------------------------------------------------

\usepackage{geometry}

\geometry{
	paper=a4paper, 
%	top=1cm, 
%	bottom=1cm, 
%	left=1cm, 
%	right=1cm, 
%	headheight=20pt, 
%	footskip=1.4cm, 
%	headsep=10pt, 
%	showframe, 
	showcrop 
}

%---------------------------------------------------------------
%	FOOTNOTE (format)
%---------------------------------------------------------------

\makeatletter
\long\def\@makefntextFB#1{%
    \ifx\thefootnote\ftnISsymbol
        \@makefntextORI{#1}%
    \else
        \rule\z@\footnotesep
        \setbox\@tempboxa\hbox{\@thefnmark}%
            \ifdim\wd\@tempboxa>\z@
                \kern2em\llap{\@thefnmark.\kern0.5em}%
            \fi
        \hangindent2em\hangafter\@ne#1
    \fi}
\makeatother

%-----------------------------------------------------
%	COLOR
%-----------------------------------------------------
%\usepackage{xcolor} % Required for specifying colors by name
%\definecolor{uMainColor}{RGB}{160,0,0} 
\colorlet{edxcolorcover}{uCoverColor}
%\definecolor{edxcolorcover}{RGB}{160,0,0} 

%\definecolor{cnam}{RGB}{128, 0, 32} 
                 

  \definecolor{comments}{rgb}{0.7,0,0}    % rouge foncé
  \definecolor{link}{rgb}{0,0.4,0.6}      % ~RoyalBlue de dvips
  \definecolor{url}{rgb}{0.6,0,0}         % rouge-brun
  \definecolor{citation}{rgb}{0,0.5,0}    % vert foncé
  \definecolor{ULlinkcolor}{rgb}{0,0,0.3} % de ulthese.cls
  \definecolor{rouge}{rgb}{0.85,0,0.07}   % rouge bandeau identitaire
  \definecolor{or}{rgb}{1,0.8,0}          % or bandeau identitaire

%-----------------------------------------------------
%	FONTS
%-----------------------------------------------------

%%\usepackage{avant} % Use the Avantgarde font for headings
%\usepackage{times} % Use the Times font for headings
%\usepackage{mathptmx} % Use the Adobe Times Rgman as the default text font together with math symbols from the Sym­bol, Chancery and Com­puter Modern fonts

%%\usepackage{microtype} % Slightly tweak font spacing for aesthetics
%%\usepackage[utf8]{inputenc} % Required for including letters with accents
%\usepackage[T1]{fontenc} % Use 8-bit encoding that has 256 glyphs

% Using AVANT GARDE Family

%\renewcommand{\familydefault}{\sfdefault}

%\fontfamily{pag}\selectfont
%\renewcommand{\familydefault}{pag} % sfdefault pag lmss

%bch         Charter
%lmr         Latin Modern Roman
%lmss        Latin Modern Sans Serif
%lmssq       Latin Modern Sans Serif extended
%lmtt        Latin Modern Typewriter
%lmvtt       Latin Modern Typewriter proportional
%pag         Avant Garde
%pbk         Bookman
%pcr         Courier
%phv         Helvetica
%pnc         New Century Schoolbook
%ppl         Palatino
%ptm         Times
%put         Utopia
 
%\renewcommand{\rmdefault}{pag} % text 
%\renewcommand{\sfdefault}{pag} % titre


%\renewcommand{\rmdefault}{pag} % text 
%\renewcommand{\sfdefault}{phv} % titre


%-----------------------------------------------------
%	BIBLIOGRAPHY AND INDEX
%-----------------------------------------------------


\usepackage{calc} % For simpler calculation - used for spacing the index letter headings correctly

\usepackage{imakeidx}
%\makeindex % Tells LaTeX to create the files required for indexing

\makeindex[columns=2,intoc=true,options={-s \upath/indexstyle.ist}]

%-----------------------------------------------------
%	HEADERS AND FOOTERS
%-----------------------------------------------------


\usepackage{fancyhdr} % Required for header and footer configuration
\pagestyle{fancy} % Enable the custom headers and footers

\ifcsname chapter \endcsname   % verify if command name exists in the CLASS (Book, report, article ...)
% BEGIN IF
\renewcommand{\chaptermark}[1]{\markboth{\sffamily\normalsize\bfseries\chaptername\ \thechapter.\ #1}{}} % Styling for the current chapter in the header
  \else
%nothing
\fi
% END IF

\renewcommand{\sectionmark}[1]{\markright{\sffamily\normalsize\thesection\hspace{5pt}#1}{}} % Styling for the current section in the header

\fancyhf{} % Clear default headers and footers
\fancyhead[LE,RO]{\sffamily\normalsize\thepage} % Styling for the page number in the header
\fancyhead[LO]{\sffamily\normalsize\rightmark} % Print the nearest section name on the left side of odd pages
\fancyhead[RE]{\sffamily\normalsize\leftmark} % Print the current chapter name on the right side of even pages
\fancyfoot[RE,LO]{\sffamily\normalsize\uCoursetittle} % Uncomment to include a footer
\fancyfoot[LE,RO]{\hyperlink{toc}{\sffamily\normalsize\uMinilogo}} % Uncomment to include a footer
% \hyperlink{toc}
\renewcommand{\headrulewidth}{1.0pt} % Thickness of the rule under the header

\fancypagestyle{plain}{% Style for when a plain pagestyle is specified
	\fancyhead{}\renewcommand{\headrulewidth}{0pt}} 

% Removes the header from odd empty pages at the end of chapters
\makeatletter
\renewcommand{\cleardoublepage}{
\clearpage\ifodd\c@page\else
\hbox{}
\vspace*{\fill}
\thispagestyle{empty}
\newpage
\fi}


%----------------------------------------------------------------------------------------
%	MAIN TABLE OF CONTENTS
%----------------------------------------------------------------------------------------
%\usepackage{titletoc} % Required for manipulating the table of contents

\contentsmargin{0cm} % Removes the default margin

% Part text styling (this is mostly taken care of in the PART HEADINGS section of this file)
\titlecontents{part}
	[0cm] % Left indentation
	{\addvspace{20pt}\bfseries} % Spacing and font options for parts
	{}
	{}
	{}

% Chapter text styling
\titlecontents{chapter}
	[1.25cm] % Left indentation
	{\addvspace{12pt}\large\sffamily\bfseries} % Spacing and font options for chapters
	{\color{uMainColor!60}\contentslabel[\Large\thecontentslabel]{1.25cm}\color{uMainColor}} % Formatting of numbered sections of this type
	{\color{uMainColor}} % Formatting of numberless sections of this type
	{\color{uMainColor!60}\normalsize\;\titlerule*[.5pc]{.}\;\thecontentspage} % Formatting of the filler to the right of the heading and the page number

% Section text styling
\titlecontents{section}
	[1.25cm] % Left indentation
	{\addvspace{3pt}\sffamily\bfseries} % Spacing and font options for sections
	{\contentslabel[\thecontentslabel]{1.25cm}} % Formatting of numbered sections of this type
	{} % Formatting of numberless sections of this type
	{\hfill\color{black}\thecontentspage} % Formatting of the filler to the right of the heading and the page number

% Subsection text styling
\titlecontents{subsection}
	[1.25cm] % Left indentation
	{\addvspace{1pt}\sffamily\small} % Spacing and font options for subsections
	{\contentslabel[\thecontentslabel]{1.25cm}} % Formatting of numbered sections of this type
	{} % Formatting of numberless sections of this type
	{\ \titlerule*[.5pc]{.}\;\thecontentspage} % Formatting of the filler to the right of the heading and the page number

%\titlecontents{subsubsection}
%	[1.25cm] % Left indentation
%	{\addvspace{1pt}\sffamily\tiny} % Spacing and font options for subsections
%	{\contentslabel[\thecontentslabel]{1.25cm}} % Formatting of numbered sections of this type
%	{} % Formatting of numberless sections of this type
%	{\ \titlerule*[.5pc]{.}\;\thecontentspage} % Formatting of the filler to the right of the heading and the page number

\titlecontents{subsubsection}
[1.25cm]
  {\addvspace{1pt}\sffamily\tiny}
 {\contentslabel[\thecontentslabel]{1.25cm}}
  {}
  {}

% Figure text styling
\titlecontents{figure}
	[1.25cm] % Left indentation
	{\addvspace{1pt}\sffamily\small} % Spacing and font options for figures
	{\thecontentslabel\hspace*{1em}} % Formatting of numbered sections of this type
	{} % Formatting of numberless sections of this type
	{\ \titlerule*[.5pc]{.}\;\thecontentspage} % Formatting of the filler to the right of the heading and the page number

% Table text styling
\titlecontents{table}
	[1.25cm] % Left indentation
	{\addvspace{1pt}\sffamily\small} % Spacing and font options for tables
	{\thecontentslabel\hspace*{1em}} % Formatting of numbered sections of this type
	{} % Formatting of numberless sections of this type
	{\ \titlerule*[.5pc]{.}\;\thecontentspage} % Formatting of the filler to the right of the heading and the page number

%----------------------------------------------------------------------------------------
%	MINI TABLE OF CONTENTS IN PART HEADS
%----------------------------------------------------------------------------------------

% Chapter text styling
\titlecontents{lchapter}
	[0em] % Left indentation
	{\addvspace{15pt}\large\sffamily\bfseries} % Spacing and font options for chapters
	{\color{uMainColor}\contentslabel[\Large\thecontentslabel]{1.25cm}\color{uMainColor}} % Chapter number
	{}  
	{\color{uMainColor}\normalsize\sffamily\bfseries\;\titlerule*[.5pc]{.}\;\thecontentspage} % Page number

% Section text styling
\titlecontents{lsection}
	[0em] % Left indentation
	{\sffamily\small} % Spacing and font options for sections
	{\contentslabel[\thecontentslabel]{1.25cm}} % Section number
	{}
	{}

% Subsection text styling (note these aren't shown by default, display them by searchings this file for tgcdepth and reading the commented text)
%\titlecontents{lsubsection}
%	[.5em] % Left indentation
%	{\sffamily\footnotesize} % Spacing and font options for subsections
%	{\contentslabel[\thecontentslabel]{1.25cm}}
%	{}
%	{}


\usepackage[]{ccicons}


%----------------------------------------------------------------------------------------
%	THEOREM STYLES
%----------------------------------------------------------------------------------------

\usepackage{amsmath,amsfonts,amssymb,amsthm} % For math equations, theorems, symbols, etc

\newcommand{\intoo}[2]{\mathopen{]}#1\,;#2\mathclose{[}}
\newcommand{\ud}{\mathop{\mathrm{{}d}}\mathopen{}}
\newcommand{\intff}[2]{\mathopen{[}#1\,;#2\mathclose{]}}
\renewcommand{\qedsymbol}{$\blacksquare$}

\ifcsname chapter \endcsname   % verify if command name exists in the CLASS (Book, report, article ...)
	\newtheorem{notation}{Notation}[chapter]
 \else
	\newtheorem{notation}{Notation}[section]
\fi


% Boxed/framed environments
\newtheoremstyle{ocrenumbox}% Theorem style name
{0pt}% Space above
{0pt}% Space below
{\normalfont}% Body font
{}% Indent amount
{\small\bf\sffamily\color{uMainColor}}% Theorem head font
{\;}% Punctuation after theorem head
{0.25em}% Space after theorem head
{\small\sffamily\color{uMainColor}\thmname{#1}\nobreakspace\thmnumber{\@ifnotempty{#1}{}\@upn{#2}}% Theorem text (e.g. Theorem 2.1)
\thmnote{\nobreakspace\the\thm@notefont\sffamily\bfseries\color{black}---\nobreakspace#3.}} % Optional theorem note

\newtheoremstyle{ocrenumbox}% Theorem style name
{0pt}% Space above
{0pt}% Space below
{\normalfont}% Body font
{}% Indent amount
{\small\bf\sffamily\color{uMainColor}}% Theorem head font
{\;}% Punctuation after theorem head
{0.25em}% Space after theorem head
{\small\sffamily\color{uMainColor}\thmname{#1}\nobreakspace\thmnumber{\@ifnotempty{#1}{}\@upn{#2}}% Theorem text (e.g. Theorem 2.1)
\thmnote{\nobreakspace\the\thm@notefont\sffamily\bfseries\color{black}---\nobreakspace#3.}} % Optional theorem note

\newtheoremstyle{blacknumex}% Theorem style name
{5pt}% Space above
{5pt}% Space below
{\normalfont}% Body font
{} % Indent amount
{\small\bf\sffamily}% Theorem head font
{\;}% Punctuation after theorem head
{0.25em}% Space after theorem head
{\small\sffamily{\tiny\ensuremath{\blacksquare}}\nobreakspace\thmname{#1}\nobreakspace\thmnumber{\@ifnotempty{#1}{}\@upn{#2}}% Theorem text (e.g. Theorem 2.1)
\thmnote{\nobreakspace\the\thm@notefont\sffamily\bfseries---\nobreakspace#3.}}% Optional theorem note


\newtheoremstyle{blacknumboxS}% name of the style to be used
{5pt}% measure of space to leave above the theorem. E.g.: 3pt
{5pt}% measure of space to leave below the theorem. E.g.: 3pt
{\normalfont}% name of font to use in the body of the theorem
{}% measure of space to indent
{\small\bf\sffamily}% name of head font
{}% punctuation between head and body
{ }% space after theorem head; " " = normal interword space
{\color{black}\faEye~  \color{uMainColor}\small\sffamily\thmnote{#3 : }}

% Non-boxed/non-framed environments
\newtheoremstyle{ocrenum}% Theorem style name
{5pt}% Space above
{5pt}% Space below
{\normalfont}% Body font
{}% Indent amount
{\small\bf\sffamily\color{uMainColor}}% Theorem head font
{\;}% Punctuation after theorem head
{0.25em}% Space after theorem head
{\small\sffamily\color{uMainColor}\thmname{#1}\nobreakspace\thmnumber{\@ifnotempty{#1}{}\@upn{#2}}% Theorem text (e.g. Theorem 2.1)
\thmnote{\nobreakspace\the\thm@notefont\sffamily\bfseries\color{black}---\nobreakspace#3.}} % Optional theorem note
\makeatother

% Defines the theorem text style for each type of theorem to one of the three styles above



\ifcsname chapter \endcsname

\newcounter{dummy} 

\numberwithin{dummy}{section}
	\theoremstyle{ocrenumbox}
		\newtheorem{theoremeT}[dummy]{Proposition}
		\newtheorem{problem}{Problem}[chapter]
		\newtheorem{exerciseT}{Outillage}[chapter]
	\theoremstyle{blacknumex}
		\newtheorem{exampleT}{Example}[chapter]
	\theoremstyle{blacknumbox}
		\newtheorem{vocabulary}{Vocabulary}[chapter]
		\newtheorem{definitionT}{Definition}[section]
		\newtheorem{corollaryT}[dummy]{Concept}
	\theoremstyle{blacknumboxS}
		\newtheorem{remarqueST}[]{}
		\newtheorem{remarqueT}[dummy]{Remarque}
	\theoremstyle{ocrenum}
	\newtheorem{proposition}[dummy]{Proposition}
  
 \else
 
\newcounter{dummy} 
\numberwithin{dummy}{section}

	\theoremstyle{ocrenumbox}
		\newtheorem{theoremeT}[dummy]{Proposition}
		%\newtheorem{problem}{Problem}[section]
		\newtheorem{exerciseT}{Outillage}[section]
	\theoremstyle{blacknumex}
		\newtheorem{exampleT}{Example}[section]
	\theoremstyle{blacknumbox}
		\newtheorem{vocabulary}{Vocabulary}[section]
		\newtheorem{definitionT}{Definition}[section]
		\newtheorem{corollaryT}[dummy]{concept}
	\theoremstyle{blacknumboxS}
		\newtheorem{remarqueST}[]{}
		\newtheorem{remarqueT}[dummy]{Remarque}
	\theoremstyle{ocrenum}
		\newtheorem{proposition}[dummy]{Proposition}
\fi



%----------------------------------------------------------------------------------------
%	DEFINITION OF COLORED BOXES
%----------------------------------------------------------------------------------------

\RequirePackage[framemethod=default]{mdframed} % Required for creating the theorem, definition, exercise and corollary boxes

% Theorem box
\newmdenv[skipabove=7pt,
skipbelow=7pt,
backgroundcolor=black!5,
linecolor=uMainColor,
innerleftmargin=5pt,
innerrightmargin=5pt,
innertopmargin=5pt,
leftmargin=0cm,
rightmargin=0cm,
innerbottommargin=5pt]{tBox}

% Exercise box	  
\newmdenv[skipabove=7pt,
skipbelow=7pt,
rightline=false,
leftline=true,
topline=false,
bottomline=false,
backgroundcolor=uMainColor!10,
linecolor=uMainColor,
innerleftmargin=5pt,
innerrightmargin=5pt,
innertopmargin=5pt,
innerbottommargin=5pt,
leftmargin=0cm,
rightmargin=0cm,
linewidth=4pt]{eBox}	

% Definition box
\newmdenv[skipabove=7pt,
skipbelow=7pt,
rightline=false,
leftline=true,
topline=false,
bottomline=false,
linecolor=uMainColor,
innerleftmargin=5pt,
innerrightmargin=5pt,
innertopmargin=0pt,
leftmargin=0cm,
rightmargin=0cm,
linewidth=4pt,
innerbottommargin=0pt]{dBox}	

% Corollary box
\newmdenv[skipabove=7pt,
skipbelow=7pt,
rightline=false,
leftline=true,
topline=false,
bottomline=false,
linecolor=gray,
backgroundcolor=black!5,
innerleftmargin=5pt,
innerrightmargin=5pt,
innertopmargin=5pt,
leftmargin=0cm,
rightmargin=0cm,
linewidth=4pt,
innerbottommargin=5pt]{cBox}

% Creates an environment for each type of theorem and assigns it a theorem text style from the "Theorem Styles" section above and a colored box from above
%\newenvironment{theorem}{\begin{tBox}\begin{theoremeT}}{\end{theoremeT}\end{tBox}}
\newenvironment{exercise}{\begin{eBox}\begin{exerciseT}}{\hfill{\color{uMainColor}\tiny\ensuremath{\blacksquare}}\end{exerciseT}\end{eBox}}				  
%\newenvironment{definition}{\begin{dBox}\begin{definitionT}}{\end{definitionT}\end{dBox}}	
%\newenvironment{example}{\begin{exampleT}}{\hfill{\tiny\ensuremath{\blacksquare}}\end{exampleT}}		
%\newenvironment{corollary}{\begin{cBox}\begin{corollaryT}}{\end{corollaryT}\end{cBox}}	
\newenvironment{nota}{\begin{cBox}\begin{remarqueT}}{\end{remarqueT}\end{cBox}}	

%----------------------------------------------------------------------------------------
%	REMARK ENVIRONMENT
%----------------------------------------------------------------------------------------

\newenvironment{remark}{\par\vspace{10pt}\small % Vertical white space above the remark and smaller font size
\begin{list}{}{
\leftmargin=35pt % Indentation on the left
\rightmargin=25pt}\item\ignorespaces % Indentation on the right
\makebox[-2.5pt]{\begin{tikzpicture}[overlay]
\node[draw=uMainColor!60,line width=1pt,circle,fill=uMainColor!25,font=\sffamily\bfseries,inner sep=2pt,outer sep=0pt] at (-15pt,0pt){\textcolor{uMainColor}{i}};\end{tikzpicture}} % Orange R in a circle
\advance\baselineskip -1pt}{\end{list}\vskip5pt} % Tighter line spacing and white space after remark

%----------------------------------------------------------------------------------------
%	SECTION NUMBERING IN THE MARGIN
%----------------------------------------------------------------------------------------

\makeatletter
% \renewcommand{\@seccntformat}[1]{\textcolor{uMainColor}{\csname the#1\endcsname}\quad}   

%\makeatother

%----------------------------------------------------------------------------------------
%	PART HEADINGS
%----------------------------------------------------------------------------------------

% Numbered part in the table of contents
\newcommand{\@mypartnumtocformat}[2]{%
	\setlength\fboxsep{0pt}%
	\noindent\colorbox{uMainColor!20}{\strut\parbox[c][.7cm]{\ecart}{\color{uMainColor!70}\Large\sffamily\bfseries\centering#1}}\hskip\esp\colorbox{uMainColor!40}{\strut\parbox[c][.7cm]{\linewidth-\ecart-\esp}{\Large\sffamily\centering#2}}%
}

% Unnumbered part in the table of contents
\newcommand{\@myparttocformat}[1]{%
	\setlength\fboxsep{0pt}%
	\noindent\colorbox{uMainColor!40}{\strut\parbox[c][.7cm]{\linewidth}{\Large\sffamily\centering#1}}%
}

\newlength\esp
\setlength\esp{4pt}
\newlength\ecart
\setlength\ecart{1.2cm-\esp}
\newcommand{\thepartimage}{}%
\newcommand{\partimage}[1]{\renewcommand{\thepartimage}{#1}}%
\def\@part[#1]#2{%
\ifnum \c@secnumdepth >-2\relax%
\refstepcounter{part}%
\addcontentsline{toc}{part}{\texorpdfstring{\protect\@mypartnumtocformat{\thepart}{#1}}{\partname~\thepart\ ---\ #1}}
\else%
\addcontentsline{toc}{part}{\texorpdfstring{\protect\@myparttocformat{#1}}{#1}}%
\fi%
\startcontents%
\markboth{}{}%
{\thispagestyle{empty}%
\begin{tikzpicture}[remember picture,overlay]%
\node at (current page.north west){\begin{tikzpicture}[remember picture,overlay]%	
\fill[uMainColor!20](0cm,0cm) rectangle (\paperwidth,-\paperheight);
\node[anchor=north] at (4cm,-3.25cm){\color{uMainColor!40}\fontsize{220}{100}\sffamily\bfseries\thepart}; 
\node[anchor=south east] at (\paperwidth-1cm,-\paperheight+1cm){\parbox[t][][t]{8.5cm}{
\printcontents{l}{0}{\setcounter{tocdepth}{1}}% The depth to which the Part mini table of contents displays headings; 0 for chapters only, 1 for chapters and sections and 2 for chapters, sections and subsections
}};
\node[anchor=north east] at (\paperwidth-1.5cm,-3.25cm){\parbox[t][][t]{15cm}{\strut\raggedleft\color{white}\fontsize{30}{30}\sffamily\bfseries#2}};
\end{tikzpicture}};
\end{tikzpicture}}%
\@endpart}
\def\@spart#1{%
\startcontents%
\phantomsection
{\thispagestyle{empty}%
\begin{tikzpicture}[remember picture,overlay]%
\node at (current page.north west){\begin{tikzpicture}[remember picture,overlay]%	
\fill[uMainColor!20](0cm,0cm) rectangle (\paperwidth,-\paperheight);
\node[anchor=north east] at (\paperwidth-1.5cm,-3.25cm){\parbox[t][][t]{15cm}{\strut\raggedleft\color{white}\fontsize{30}{30}\sffamily\bfseries#1}};
\end{tikzpicture}};
\end{tikzpicture}}
\addcontentsline{toc}{part}{\texorpdfstring{%
\setlength\fboxsep{0pt}%
\noindent\protect\colorbox{uMainColor!40}{\strut\protect\parbox[c][.7cm]{\linewidth}{\Large\sffamily\protect\centering #1\quad\mbox{}}}}{#1}}%
\@endpart}
\def\@endpart{\vfil\newpage
\if@twoside
\if@openright
\null
\thispagestyle{empty}%
\newpage
\fi
\fi
\if@tempswa
\twocolumn
\fi}

%----------------------------------------------------------------------------------------
%	CHAPTER HEADINGS
%-------------------------------------------------%--------------------------------------

%
%
%% A switch to conditionally include a picture, implemented by Christian Hupfer
\newif\ifusechapterimage
\usechapterimagetrue
\newcommand{\thechapterimage}{}%
\newcommand{\chapterimage}[1]{\ifusechapterimage\renewcommand{\thechapterimage}{#1}\fi}%
\newcommand{\autodot}{.}

\def\@makechapterhead#1{%
{\parindent \z@ \raggedright \normalfont
\ifnum \c@secnumdepth >\m@ne
\if@mainmatter
\begin{tikzpicture}[remember picture,overlay]
	
\node at (current page.north west)
			{\begin{tikzpicture}[remember picture,overlay]
			\node[anchor=north west,inner sep=0pt] at (0,0) {\ifusechapterimage\includegraphics[width=\paperwidth]{\thechapterimage}\fi};
			\draw[anchor=west] (\Gm@lmargin,-9cm) node [line width=2pt,rounded corners=15pt,draw=uMainColor,fill=white,fill opacity=1,inner sep=15pt]{\strut\makebox[22cm]{}};
			
			\draw[anchor=west] (\Gm@lmargin+.3cm,-9cm) node {\fontsize{20}{30}\selectfont\sffamily\bfseries\color{uMainColor}~#1\strut};
			%{\Huge\sffamily\bfseries\color{uMainColor}\thechapter\autodot~#1\strut};
			\node at (current page.north east)  [xshift=-3.95cm, yshift=-4.05cm, text opacity=1]  {{\color{white}\centering\fontsize{200}{30}\selectfont \bfseries\sffamily\thechapter\strut} };
			\node at (current page.north east)  [xshift=-4cm, yshift=-4cm, text opacity=1]  {{\color{uMainColor}\centering\fontsize{200}{30}\selectfont \bfseries\sffamily\thechapter\strut} };
			
			\end{tikzpicture}};
	\end{tikzpicture}

\else

\begin{tikzpicture}[remember picture,overlay]
\node at (current page.north west)
{\begin{tikzpicture}[remember picture,overlay]
\node[anchor=north west,inner sep=0pt] at (0,0) {\ifusechapterimage\includegraphics[width=\paperwidth]{\thechapterimage}\fi};
\draw[anchor=west] (\Gm@lmargin,-9cm) node [line width=2pt,rounded corners=15pt,draw=uMainColor,fill=white,fill opacity=1,inner sep=15pt]{\strut\makebox[22cm]{}};

\draw[anchor=east] (\Gm@lmargin,-9cm) node [line width=2pt,rounded corners=15pt,draw=uMainColor,fill=white,fill opacity=1,inner sep=15pt]{\strut\makebox[22cm]{}};
\draw[anchor=west] (\Gm@lmargin+.3cm,-9cm) node {\huge\sffamily\bfseries\color{uMainColor}#1\strut};

\end{tikzpicture}};
 
\end{tikzpicture}
\fi\fi\par\vspace*{270\p@}}}

\makeatother
%-------------------------------------------

%\def\@makeschapterhead#1{%
%\begin{tikzpicture}[remember picture,overlay]
%\node at (current page.north west)
%{\begin{tikzpicture}[remember picture,overlay]
%\node[anchor=north west,inner sep=0pt] at (0,0) {\ifusechapterimage\includegraphics[width=\paperwidth]{\thechapterimage}\fi};
%\draw[anchor=west] (\Gm@lmargin,-9cm) node [line width=2pt,rounded corners=15pt,draw=uMainColor,fill=white, fill opacity=0.8,inner sep=15pt]{\strut\makebox[22cm]{}};
%\draw[anchor=west] (\Gm@lmargin+.3cm,-9cm) node {\huge\sffamily\bfseries\color{uMainColor}#1\strut};
%\end{tikzpicture}};
%\end{tikzpicture}
%\par\vspace*{270\p@}}




%----------------------------------------------------------------------------------------
%	LINKS
%----------------------------------------------------------------------------------------


\usepackage{bookmark}
\bookmarksetup{
open,
numbered,
addtohook={%
\ifnum\bookmarkget{level}=0 % chapter
\bookmarksetup{bold}%
\fi
\ifnum\bookmarkget{level}=-1 % part
\bookmarksetup{color=uMainColor,bold}%
\fi
}
}

%----------------------------------------------------------------------------------------
%	FANCY CHAPTER Header
%----------------------------------------------------------------------------------------





%\definecolor{uMainColor}{RGB}{255,128,0}   



%-------------------------------------------------
%	BIBLIOGRAPHY
%-------------------------------------------------

\addbibresource{../Latex/Sources/commons/bibliography.bib} % BibTeX bibliography file
\defbibheading{bibempty}{}

%-------------------------------------------------
%	COLORS & BORDERS
%-------------------------------------------------
\setlength{\fboxrule}{0.75pt} % Width of the border around the abstract
%\definecolor{uMainColor}{RGB}{255,128,0} % Color of the article title and section
\definecolor{color2}{RGB}{220,220,220} % Color of the boxes behind the abstract and headings

%-------------------------------------------------
% ITEMS DEFINITION
%-------------------------------------------------

\setlist [itemize,1]{label=\color{uMainColor}\faCaretRight }


\usepackage{environ}
\newwrite\audioexport



\begin{document}





  %-------------------------------------------------------
  % AudioNote Open File
  %-------------------------------------------------------
\mode<all>{\immediate\openout\audioexport=\jobname.txt}


\addtocontents{toc}{\protect\hypertarget{toc}{}}

%-------------------------------------------------
%	ARTICLE INFORMATION
%-------------------------------------------------

\Abstract {\uabstract%===========================
% COURS "INTRO CYBERDEF"
% Abstract général des articles
%===========================


Il fait partie du cours introductif aux fondamentaux de la sécurité des systèmes d'information vue sous deux prismes quelques fois opposés dans la littérature : la gouvernance et la gestion opérationnelle de la sécurité.
Le cours est constitué d'un ensemble de notes de synthèse indépendantes compilées en un document unique, mais édité par chapitre dans le cadre de ce cours.\\
C\edoc ne constitue pas à lui seul le référentiel du cours CYBERDEF101 (SEC101 du Cnam). Il compile des notes de cours mises à disposition de l'auditeur comme support pédagogique partiel à ce cours introductif à la cyberdéfense d'entreprise. }
\JournalInfo{\uJournalInfo} % Journal information
\Archive{Notes de cours éditées le  \DTMnow} % Additional notes (e.g. copyright, DOI, review/research article)
\PaperTitle{\utitle} % Article title
\Authors{\uauthor\textsuperscript{1,}\textsuperscript{2}*} % Authors
\affiliation{\textsuperscript{1}\textit{\uproa}} % Author affiliation
\affiliation{\textsuperscript{2}\textit{\uprob}} % Author affiliation
\affiliation{*\textbf{email}: \umaila\ -- \umailb} % Corresponding author
\Keywords{\ukeywords} % Keywords 
\newcommand{\keywordname}{Mots clefs} % Defines the keywords heading name

%-------------------------------------------------
%	MAKE TITLE
%-------------------------------------------------

\maketitle

\begin{tikzpicture}[remember picture,overlay]\node at (current page.north west)[xshift=3cm, yshift=-3cm, text opacity=1.0]{\includegraphics[width=0.1\paperwidth]{\utemplatepath/commons.inc/CommonsPictures/\utemplatetheme.theme/shield-20.pdf}}; 
\end{tikzpicture} 


\AddToShipoutPicture{%
\AtPageCenter{%
\begin{tikzpicture}[remember picture,overlay]
\node [rotate=60,scale=5,text opacity=0.1] at (current page.center) { \textit{\color{uAlertColor}\INFODistrib} };
\end{tikzpicture} 
%\begin{tikzpicture}[remember picture,overlay]
%\node [rotate=60,scale=10,text opacity=0.3]
%at (current page.center) {moi};
%\end{tikzpicture}
}%
}


%-------------------------------------------------	
\begin{center}
	\textcolor{uMainColor}{\textbf{\INFODistrib}} % Distribution information
\end{center}

\INFOGitHub % Link to download file and copyright
%-------------------------------------------------

%-------------------------------------------------
	\newpage
	\tableofcontents  % TOC
	\iftotalfigures
  			\listoffigures
	\fi
%-------------------------------------------------

%-------------------------------------------------
\newpage
\ubody % MAIN BODY defined external
%-------------------------------------------------

	\printbibliography
 \printglossary[type=\acronymtype]
%	\printglossary
%-------------------------------------------------	
\newpage
\Ucontribute % Contribution TEXT
%-------------------------------------------------

  %-------------------------------------------------------
  % AudioNote Open File
  %-------------------------------------------------------
  \mode<all>{\immediate\closeout\audioexport}

\end{document}
}
	\fi
\fi


 


%************************************************************
% Chargement des variables dédiées à l'article
%************************************************************
\lessonID{7.1}
\setauthor{EDU}

%========================================
\newcommand {\ukeywords}
%========================================
{%---------------------------------------------------------------------
Vulnérabilités, Pentest, Bugbounty, Audit sécurité
}%

%========================================
\newcommand {\utitle}
{%---------------------------------------------------------------------
Vulnerability Management
}%---------------------------------------------------------------------

%========================================
\newcommand {\uabstract}
{%---------------------------------------------------------------------
C\edoc introduit la gestion de fragilités du numérique de la partie cyberdéfense de la sécurité opérationnelle : Anticiper, Détecter, Réagir  et ceci sur les trois grands invariants des risques numériques : les vulnérabilités, les menaces et l'impact.
}
%---------------------------------------------------------------------

%************************************************************
%  variable définissant  le corps de l'article
%************************************************************

%========================================
\newcommand {\ubody}
{%---------------------------------------------------------------------
%$BEGIN_UBODY


% VUL MAN

\mode<all>{%-------------------------------------------------------------
%               FR CYBERDEF SECOPS COURSE
%                                        SECOPS
%                    Vulnerability/Threat/Incident
%.           $Chapitre : Vulnerabilty Management
%                                $theme : Intro
%                    $File : chap-Vulman-into.tex
%                             2020 eduf@ction
%-------------------------------------------------------------
\uchap{Introduction sur les vulnérabilités, chapitre 3.1}
%-------------------------------------------------------------

\section{Fragilités numériques}

Nous avons vu dans l'équation d'évaluation simple du risque, que ce dernier dépendait directement des fragilités de l'entreprise. C'est par l'exploitation de ces failles que l'attaquant va pouvoir déployer toutes ses ambitions.

% Begin PRZ ===========================
\begin{frame}
\frametitle<presentation>{Fragilités et Vulnérabilités}
% end header PRZ =======================
La notion de fragilité numérique ou digitale de l'entreprise est à prendre au sens large. Elle comprend les fragilités \textbf{humaines}, \textbf{organisationnelles} et \textbf{techniques} mais aussi la sensibilité à des scénarios d'attaques. C'est en effet la susceptibilité d’une organisation à subir des défaillances dans le temps que l'on nomme vulnérabilités.
\end{frame}
% end PRZ ===========================

\subsection{Détecter les fragilités de l’entreprise}

La première tâche de fond en cybersécurité pour une équipe dédiée est d'identifier les fragilités de l'ensemble de l'environnement numérique \footnote{systèmes d'information de l'entreprise, services dans les cloud, réseaux sociaux...} de l'entreprise. Elle s'inscrit dans la dynamique de l'\textbf{anticipation} avec la recherche de fragilités ou de risques cyber dans l'entreprise et leur correction. Généralement les tâches associées à la couverture des vulnérabilités se déploient avant la \textbf{détection} d'évènement à risque, d'attaques, de déviance dans l'environnement mais aussi à l'extérieur du périmètre de l'entreprise.
Elle se positionne néanmoins comme une activité qui peut déclencher des mécanismes de la \textbf{réaction} aux incidents, de la gestion de crise, par la nécessaire remédiation en cas de vulnérabilité critique.
De multiples notions sont liées à la gestion des vulnérabilités, telles que le « scan » de vulnérabilités, l’évaluation de vulnérabilités (vulnerability assessment en anglais) ou l’application de correctifs (vulnerability patching ou patch management en anglais). On trouvera ces concepts bien décrits dans le livre blanc \ulink{https://www.sans.org/reading-room/whitepapers/threats/implementing-vulnerability-management-process-34180}{du SANS Institute Implementing a Vulnerability Management Process}.


% Begin PRZ ===========================
\begin{frame}
\frametitle<presentation>{Identifier ses vulnérabilités}
% end header PRZ =======================
On peut distinguer deux grandes typologies d'actions pour identifier ces fragilités :
\begin{itemize}
	\item l'audit de sécurité, qui permet de détecter des fragilités exploitables. Ce type d'audit peut se dérouler sous la forme de scénario exécuté par des équipes de \g{tests d'intrusion} soit sous la forme de campagne exécutée avec des scanners de vulnérabilités.
	\item la veille en vulnérabilités associée à la cartographie de l'environnement technique permettent de déclencher une alerte de sécurité si une vulnérabilité apparaissait sur un des produits, services ou logiciel surveillés.
\end{itemize}
\end{frame}
% end PRZ ===========================

La difficulté principale de ces activités est de bien définir les périmètres techniques et de responsabilité sur lesquelles elles portent.

Si l'audit de sécurité permet d'évaluer les fragilités des éléments (composants) de l'entreprise en se mettant dans la peau de l'attaquant, afin de découvrir les scénarios potentiellement actifs sur l'environnement digital de l'entreprise, il n'en demeure pas moins important de mettre en place des mécanismes complémentaires et continus pour la veille, la recherche, la détection, la correction de ces vulnérabilités.


\subsection{Anticiper et surveiller les menaces }

Comme nous l'avons vu, une grande partie des attaques sur l'entreprise est liée à l'exploitation de fragilités de celle-ci, ces fragilités étant dans la plupart des cas connues.  

% Begin PRZ ===========================
\begin{frame}
\frametitle<presentation>{Exploitation des vulnérabilités}
L'exploitation de ces fragilités, sont de deux grandes natures :
% end header PRZ =======================
\begin{itemize}
	\item attaques exploitant de manière \textbf{opportuniste} des fragilités non cataloguées  avec ou sans ciblage particulier de l'attaqué;
	\item attaques \textbf{ciblées} exploitant de manière spécifique des fragilités connues mais pas corrigées ou des fragilités non encore connues par les défenseurs.
\end{itemize}
\end{frame}
% end PRZ ===========================

On trouvera dans le chapitre \ref{chap:GESTVUL}, une description plus précise de ces notions de vulnérabilités connues et non connues. Les menaces sont généralement des scénarios, des codes malveillants, des mécanismes d'agression ...
Le principe de gestion de la menace relève de la même dynamique de gestion que celle liée aux vulnérabilités. 


}
\mode<all>{%-------------------------------------------------------------
%               FR CYBERDEF SECOPS COURSE
%             $Chapitre : Vulnerabilty Management
%                               $theme : Intro
%                    $File : chap-Vulman-def.tex
%                             2020 eduf@ction
%-------------------------------------------------------------

%----------------------------------------------------------
% Les bases sur vulnérabilités
%----------------------------------------------------------
\subsection{Les basics sur les vulnérabilités HOT}

Quand nous parlons de vulnérabilités, nous parlons globalement des fragilités dans l’environnement du numérique de l’entreprise. Nous pouvons distinguer trois grandes classes de fragilités :

% Begin PRZ ===========================
\begin{frame}
\frametitle<presentation>{Fragilités HOT}
% end header PRZ =======================
Les fragilités de base dans une entreprise
\upicture{../Latex/Sources/EDU/SRC1/Pictures/img-hot}{les types de vulnérabilités}{0.7}{lblvulhot}
\end{frame}
% end PRZ ===========================

% Begin PRZ ===========================
\begin{frame}
\frametitle<presentation>{Fragilités HOT}
% end header PRZ =====================
\begin{itemize}
\item Fragilités techniques : généralement dénommées vulnérabilités au sens où ces fragilités rendent vulnérable tout ou partie d’un système. Pour rechercher ces vulnérabilités, on utilisera des techniques d’audit, de scan de fuzzing ... Ce sont ces vulnérabilités informatiques et réseaux que nous présenterons plus en détails;
\item Fragilités humaines : généralement des déviances comportementales, détournement d’usage légitime, sensibilité à l’ingénierie sociale, vulnérabilités sociales ou physiologiques que l’attaquant pour utiliser. Ces fragilités sont détectables avec des audits (exemple tests mail phishing). Elles sont réduites par des mécanismes de formations et de sensibilisation, ainsi que dans certains cas des processus d’habilitation ;
\item Fragilités organisationnelles : un attaquant peut utiliser des déficiences organisationnelles pour obtenir des éléments pour conduire son attaque (exemple : pas de processus de vérification d’identité lors de demande sensible par téléphone).
\end{itemize}
\end{frame}
% end PRZ ===========================



Un attaquant utilisera bien entendu l’ensemble de ces fragilités pour conduire sa mission.

Dans le domaine technique de cette sécurité numérique, une vulnérabilité ou faille est une faiblesse dans un système, permettant à un attaquant de porter atteinte à la fonction de ce système, c'est-à-dire à son fonctionnement normal, à la confidentialité et l'intégrité des données qu'il contient.

Ces vulnérabilités sont la conséquence de faiblesses dans la conception, le développement, le déploiement, la mise en œuvre ou l'utilisation d'un composant matériel ou logiciel du système. 

Il y a trois grandes classes de faiblesses ou vulnérabilités numériques :

% Begin PRZ ===========================
\begin{frame}
\frametitle<presentation>{Zoom Fragilités TECHNIQUES}
% end header PRZ =======================
\begin{itemize}
\item \textbf{Failles de configuration} ou de défaut d’usage (utilisation d’un système en dehors de ses zones de fonctionnement stable et maitrisé)
\item \textbf{Failles Logicielles } : failles de développement, de programmation qui conduisent généralement de l'exploitation de bugs logiciels. Il faut distinguer les logiciels développés de manière dédiée, et les logiciels dits sur étagère.
Les dysfonctionnements des logiciels sur étagère (éditeurs logiciels) sont en général corrigés à mesure de leurs découvertes, mais il y a un délai entre le moment de la découverte et la correction,
\item \textbf{Failles de conception} : failles issues de défaut de conception. Ces failles sont souvent liées à des failles protocolaires issues de faille de conception d'un protocole de communication, ou de format de données.
\end{itemize}
\end{frame}
% end PRZ ===========================

%%%%%%%%

% Begin PRZ ===========================
\begin{frame}
\frametitle<presentation>{Groupes de failles}
% end header PRZ =======================
Nous pouvons décomposer les failles dites logicielles, en deux groupes 
\begin{itemize}
\item Les failles des logiciels ou \textbf{codes sur mesure}, développés dans l'entreprise ou par un tiers mais non édités en tant que logiciel indépendant. Nous pouvons y inclure tous les codes logiciels développés en interne.
\item Les failles logicielles de produits ou codes connus, reconnus souvent dénommées \textbf{progiciels} (produits logiciels). On peut aussi y distinguer deux sous classes les logiciels où les sources sont accessibles, et les codes dits fermés ou l'utilisateur ne dispose que du code binaire exécutable. Nous verrons que les démarches de recherche de failles dans ces deux types de code sont un peu différentes.
\end{itemize}

\end{frame}
% end PRZ ===========================


% Begin PRZ ===========================
\begin{frame}
\frametitle<presentation>{Typologie de failles}
% end header PRZ =======================
\upicture{../Latex/Sources/EDU/SRC1/Pictures/img-vul-types}{Les types de vulnérabilités}{1}{lblvultypes}
\end{frame}
% end PRZ ===========================

% Begin PRZ ===========================
\begin{frame}
\frametitle<presentation>{AllowAll vs DenyAll}
% end header PRZ =======================
Quand on parle de fragilités, il n'y pas que les failles de conception ou de développement. Les failles de configuration des systèmes d'information représentent encore une grande partie des fragilités utilisées par les attaques.
On trouve encore des administrateurs système qui utilisent dans les outils de filtrage la règle :

\begin{notebox}{AllowAll vs DenyAll}
Tout est autorisé sauf ce qui est interdit (\textbf{Allow All })
plutôt que de respecter le concept de base de la sécurité 
tout est interdit (\textbf{Deny All}) sauf ce qui est autorisé.
\end{notebox}

\end{frame}
% end PRZ ===========================


La recherche et la découverte de ces vulnérabilités utilisent donc des outillages un peu différents. On distinguera donc : 

\begin{itemize}
  \item les failles systèmes et de configuration,
  \item les failles dans le développement, dites failles applicatives.
 \end{itemize}

\subsection{Exemples de vulnérabilités}

A titre d'exemple et d'illustration je vous propose d'examiner rapidement des vulnérabilités techniques : deux failles de conception  et  une de programmation.
Nous ne rentrerons pas dans les détails des vulnérabilités. Ce chapitre a pour objectif de présenter concrètement ce qu'est une vulnérabilité, sur la base d'exemples simples.

\begin{techworkbox}{Exemples et cas typiques}
Les exemples de vulnérabilités tant en développement qu'en configuration et les moyens de les couvrir sont de bons sujets pour \fichetech  dans le cadre des travaux demandés.
\end{techworkbox}

% voir : https://beta.hackndo.com/la-faille-xss/

La  majorité des failles informatiques du domaine du Web et des applications sur mesure est due à une utilisation non prévue de l'applicatif. Un utilisateur peut envoyer une information plus longue que prévue (buffer overflow), ou une valeur non gérée (négative, quand le logiciel attend une valeur positive), ou quand il ajoute des symboles non attendus (des guillemets, caractères spéciaux alors qu'il était prévu seulement des lettres), si les vérifications des données ne sont pas faites correctement, alors le logiciel, le programme ou l’application peut se mettre dans un état qui, dans certains cas, peut être détourné.

\subsubsection{Faille type XSS}

Si nous prenons par exemple, un code qui affiche une image avec un titre, que ce titre d'image soit saisi par un utilisateur et qu'aucun contrôle ne soit fait. Dans l'application, l'affichage se fait par un code PHP du style : \\


% Begin PRZ ===========================
\begin{frame}[containsverbatim]
\frametitle<presentation>{Faille, type XSS}
\framesubtitle<presentation>{Exemple}
% end header PRZ =======================
%---------------------------------------------------
\begin{lstlisting} [language=php,morekeywords={readtitle,readimage}]
<?php ...
	$image = readimage()."png";
	$title = readtitle();
...
	print '<img src="$image" title="$title" />';
...?>
\end{lstlisting}
%---------------------------------------------------

et permet de générer le code HTML suivant : \\

%---------------------------------------------------
\begin{lstlisting} [language=html]
<html>...
	<img src="../path/monimage.png" title="un titre de mon image" />
...</html>
\end{lstlisting}
%---------------------------------------------------
\end{frame}
% end PRZ ===========================




% Begin PRZ ===========================
\begin{frame}[containsverbatim]
\frametitle<presentation>{Faille, type XSS}
\framesubtitle<presentation>{Exemple}
% end header PRZ =======================
Un utilisateur malveillant pourrait avoir saisi autre chose qu'un simple titre, et faire en sorte que la variable \textbf{\$tittle} puisse contenir une chaîne de caractère un peu particulière. Le pirate aura entré, par exemple, comme titre de sa photo sur ce site un peu faible, une chaine comme  : \\
\g{\verb!un titre de mon image/"><script>...script malveillant...;</script>!}
\\
%---------------------------------------------------
\begin{lstlisting} [language=html]
<html>...
	<img src="../path/monimage.png" title="un titre de mon image/"><script>...scriptmalveillant...;</script>" 
...
</html>
\end{lstlisting}
%---------------------------------------------------

\end{frame}
% end PRZ ===========================



L'exécution du script javascript malveillant se fera à la lecture de cette page générée. Si cette donnée et stockée sur un serveur, l'action sera effective pour toute les personnes qui consulteront l'image avec son titre piégé par un script malveillant.

\subsubsection{Faille type SQL Injection}

Nous allons rapidement explorer un grand classique des vulnérabilités sur les applications de sites Web sur Internet : l'injection SQL. Le principe est d'injecter dans une requête SQL (langage d'interrogation de base de données), utilisée dans un application PHP par exemple. Supposons que dans l'application, la requête suivante soit utilisée : \\
%---------------------------------------------------
\begin{lstlisting} [language=SQL]
SELECT fieldlist
  FROM table 
 WHERE field = ' $EMAIL ';
\end{lstlisting}
%---------------------------------------------------
Supposons que la saisie de l'utilisateur, saisisse un email avec une chaîne un peu modifiée (ajout d'un simple \g{'} en plus)  :
\\
%---------------------------------------------------
\begin{lstlisting} [language=SQL,escapeinside=|| ]
SELECT fieldlist
  FROM table 
 WHERE field = '  |\tikzmark{starta}|contact@test.com'|\tikzmark{enda}| ';
\end{lstlisting}
%---------------------------------------------------


%---------------------------------------------------
\begin{tikzpicture}[remember picture,overlay] 
\draw[uMainColor,rounded corners] 
    ([shift={(-3pt,2ex)}]pic cs:starta) 
    rectangle 
    ([shift={(3pt,-0.65ex)}]pic cs:enda); 
\end{tikzpicture} 
%---------------------------------------------------
L'exécution de cette requête va générer une erreur, et en fonction de la gestion des erreurs du code PHP, l'utilisateur pourra apercevoir que cette requête a provoqué une erreur d'exécution. Ceci permet à l'utilisateur de rapidement déterminer que le code est sensible à une attaque par injection SQL. Il peut alors à loisir trouver la meilleure manière de l'exploiter, en entrant un email forgé avec une chaîne plus malicieuse.\\

La chaîne  \g{\verb!OU 'x'='x'!} étant toujours VRAI, on pourrait obtenir des informations complètes de certaines tables.
Bien entendu, l'usage de vulnérabilité SQL injection n'est généralement pas trivial, mais avec un peu d'habitude, il est possible de construire des attaques sophistiquées sur des codes vulnérables.\\
%---------------------------------------------------
\begin{lstlisting} [language=SQL,escapeinside=||]
SELECT fieldlist
  FROM table 
 WHERE field = ' |\tikzmark{startb}|somebody' OU 'x' = 'x|\tikzmark{endb}| ';
\end{lstlisting}


\begin{tikzpicture}[remember picture,overlay] 
\draw[uMainColor,rounded corners] 
    ([shift={(-3pt,2ex)}]pic cs:startb) 
    rectangle 
    ([shift={(3pt,-0.65ex)}]pic cs:endb); 
\end{tikzpicture} 
%---------------------------------------------------


\subsubsection{Vulnérabilités WEB}

Vous trouverez sur le site \ulink{https://www.owasp.org}{Open Web Application Security Project}, le top TEN des vulnérabilités découvertes sur les sites WEB
et pour ceux qui souhaitent creuser un peu plus, il existe de nombreux sites présentant en détail des vulnérabilités et des manières de les exploiter (à des fins pédagogiques!). Le site de \ulink{https://beta.hackndo.com}{Pixis (Hackndo)}, par exemple, vous donne quelques partages particuliers d'un Ethical hacker.

\subsection{Failles de programmation}

%SMB etherblue
\subsubsection{Exemple : SMB etherblue}
On peut trouver des vulnérabilités dans des produits et services très connus, et déployés depuis très longtemps. 
Parmi ces vulnérabilités les plus célèbres,~\uref{voir  figure}{lblfaillesmb}, la classe de vulnérabilités du protocole SMB en version 1, est celle qui continue encore à faire des victimes.
Le protocole SMB (Server Message Block) est un protocole permettant le partage de ressources (fichiers et imprimantes) sur des réseaux locaux avec des PC sous Windows. Sa version 1 du protocole SMB, vulnérable à la faille EternalBlue.
\upicture{../Latex/Sources/EDU/SRC1/Pictures/img-smbfaille}{Tempo faille SMB - google}{0.6}{lblfaillesmb}

Dans la base de données du système CVE, on retrouve l’identifiant de cette vulnérabilité : \textbf{CVE-2017-0144}. 

\subsubsection{Exemple : programmation erronée}

Le célèbre débordement de pile, ou BUFFER OVERFLOW, fait partie des exemples d'erreur de programmation, en particulier avec des langages permissifs,  pouvant conduire à des situations exploitables par des codes malveillants.

Si nous prenons l'exemple simple d'une fonction de copie de mémoire par indexation de tableau en C suivante :

\begin{lstlisting} [language=c++]
void *memcpy(char *dest; char *src, size_t *n)
{  /* copie de bloc de memoire */
// size_t : taille du tableau
// *src : pointeur  source
// *dest : pointeur  destination
  for (size_t i=0; i<n; i++)
    dst[i] = src[i];
}
\end{lstlisting} 

Un nombre négatif (et donc incorrect en principe pour copier une zone de mémoire dans une autre) peut être passé en valeur dans 'n', et par la suite est interprété comme un grand nombre à cause du complément à 2 et de la façon dont les nombres négatifs sont codés. Ceci qui provoque un débordement de mémoire et un plantage si ce débordement touche une zone de mémoire virtuelle non liée à une zone de mémoire physique.

%https://blog.httpcs.com/heartbleed-openssl-explication-et-exploit/ (faille openssl) HeartBleed dans l'opensource

\subsection{Vulnérabilités et configuration}

Un des plus grands classiques de vulnérabilité système concerne  les défauts de configuration, en particulier les défauts de configuration des équipements et systèmes dans un environnement réseau.

L'outil le plus classique et accessible est NMAP.  Cet outil est conçu pour détecter les ports ouverts, identifier les services hébergés et obtenir des informations sur le système  de l'ordinateur distant. 
La technique de scan de port est utilisée par les administrateurs des systèmes informatiques pour contrôler la sécurité des serveurs de leurs réseaux. La même technique est aussi utilisée par les acteurs malveillants pour tenter de trouver des failles dans des systèmes informatiques. De nos jours un balayage de ports (port scan ou portscan en anglais) effectué sur un système tiers est généralement considéré comme une agression, car il préface une intrusion. Il est donc recommandé de l'utiliser de manière responsable lorsqu'on utilise ce type d'outils dans son entreprise. En effet --, 
scan de ports est une des activités considérées comme suspectes par un système de détection d'intrusion. 





\subsection{Vulnérabilités et exploits}

Il arrive que la procédure d'exploitation d'une faille d'un logiciel soit  documentée et utilisable soit sous la forme d'un code logiciel et/ou de procédure descriptive détaillée appelée \g{exploit}. Ces exploits ne sont pas systématiquement publiés.

\subsection{Vulnérabilités et divulgation}

 La divulgation publique des vulnérabilités  est soumise à un modèle de divulgation de vulnérabilité dans lequel une vulnérabilité ou un problème est révélé uniquement après une période permettant à la vulnérabilité ou au problème d'être corrigé. Cette période distingue le modèle de la divulgation complète.
 
Tout fournisseur de logiciels de sécurité, de services et de recherches de vulnérabilité, se doit de prendre des précautions vis à vis de vulnérabilités découvertes, en particulier les délais de publication. On parle généralement de \UKword{Vulnerability Disclosure Policy}.

En effet, développeurs de matériel et de logiciels ont souvent besoin de temps et de ressources pour corriger ces vulnérabilités. 

Dans certains cas, lorsque la découverte n'a pas été faite via une recherche commanditée (Audit, Pentest, BugBounty), la communauté sécurité et les scientifiques  estiment qu’il est de leur responsabilité sociale de sensibiliser le public aux vulnérabilités ayant un impact important en les publiant. Cacher ces problèmes pourrait créer un faux sentiment de sécurité. Pour éviter cela, les parties impliquées unissent leurs forces et s’accordent sur un délai pour réparer la vulnérabilité et prévenir tout dommage futur. En fonction de l'impact potentiel de la vulnérabilité, du temps requis pour qu'un correctif d'urgence ou une solution de contournement soit développé et appliqué, ainsi que d'autres facteurs, cette période peut varier de quelques jours à plusieurs mois. 

Par ailleurs, la confidentialité des découvertes est généralement requise lors des audits. Le commanditaire et l'expert signent un accord dénommé \UKword{Vulnerability Non Disclosure Agreement}, qui permet de s'assurer que la publication des vulnérabilités restera à la main du commanditaire.

Dans un mode publique avec mode de divulgation de vulnérabilités ouvertes les experts en sécurité  s'attendent à être indemnisés financièrement, mais avec le risque que signaler ces vulnérabilités au fournisseur avec l'exigence d'une indemnisation soit considéré comme une extorsion. 

Un marché des vulnérabilités s'est développé, on peut citer en particulier la société \ulink{https://zerodium.com}{Zerodium}, ou les sociétés de BugBounty, mais la commercialisation des vulnérabilités reste un sujet très controversé lié au concept de divulgation des vulnérabilités.   C'est normalement dans le rôle d'un CERT d'assurer cette coordination des divulgations.


% Begin PRZ ===========================
\begin{frame}
\frametitle<presentation>{zerodium}
% end header PRZ =======================
\upicture{../Latex/Sources/EDU/SRC1/Pictures/img-zerodium}{Le marché des failles Serveurs avec Zerodium}{0.8}{lblzerodium}
\end{frame}
% end PRZ ===========================


\begin{techworkbox}{Le marché des vulnérabilités}
Le marché des failles de sécurité est un marché particulier dans lequel des hackers de toute nature trouvent le moyen de financer leurs activités de R\&D et de hacking. Les grands éditeurs commerciaux et libres y trouvent leur compte. C'est un sujet intéressant pour \fichetech. 	
\end{techworkbox}





}
\mode<all>{%-------------------------------------
% Chapitre
% Vulnerability Management
% CVE
% File : chap-Vulman-cve.tex
%--------------------------------------
\uchap{\jobname}

\subsection{CVE, CVSS et CWE}

% Begin PRZ ===========================
\begin{frame}
\frametitle<presentation>{Common Vulnerabilities \& Weakness}
% end header PRZ =======================
\upicture{../Latex/Sources/EDU/SRC1/Pictures/img-cvecwe}{Quelques concepts de gestion sur les vulnérabilités}{0.8}{lblcvecwe}
\end{frame}
% end PRZ ===========================

\subsection{Common Vulnerabilities and Exposure (CVE)}

De nombreuses vulnérabilités sont découvertes chaque jour dans des produits et logiciels. Les informations techniques sur ces vulnérabilités permet de les détecter, et de les caractériser. il était important dans le monde des technologies de l'information qu'elles puissent être identifiées et décrites de manière unique, et que ces caractérisations soient accessibles à tous.

L’objectif fondamental de la création du CVE est de constituer un dictionnaire qui recense toutes les failles avec une description succincte de la vulnérabilité, ainsi qu’un ensemble de liens que les utilisateurs peuvent consulter pour plus d’informations. Cette base  est proposée pour consultation et reste maintenue par le Mitre Corporation.  Cet organisme à but non lucratif américaine a pour  objectif de travailler dans des domaines technologiques comme l'ingénierie des systèmes, les technologies de l'information, la sécurité. 



% Begin PRZ ===========================
\begin{frame}
\frametitle<presentation>{MITRE et codification}
% end header PRZ =======================
\UKword{ Common Vulnerabilities and Exposures} ou CVE est une base de données (Dictionnaire) des informations publiques relatives aux vulnérabilités de sécurité. Le dictionnaire est maintenu par l'organisme MITRE.  Les identifiants CVE sont des références de la forme CVE-AAAA-NNNN 
 
 Pour consulter les CVE,  il suffit de se rendre sur \ulink {https://www.cve.mitre.org}{CVE.mitre.org}

\end{frame}
% end PRZ ===========================

Le système CVE permet de recenser toutes les failles et les menaces liées à la sécurité des systèmes d’information avec  un identifiant unique  attribué à chaque faille.

On trouve ainsi dans cette base par exemple :

\begin{itemize}
  \item l’identifiant de l’une des  vulnérabilités qui a permis une attaque massive via le rançonlogiciel \textbf{Wannacry} : CVE-\textbf{2017}-0144,  faille dans le protocole SMB découverte en 2017 et la 144ième faille découverte de l'année.
  \item ou \textbf{Heardbleed} l’une des failles les plus importantes des années 2010 : CVE-\textbf{2014}-0160 présente dans la couche logicielle open source OpenSSL. OpenSSL est une librairie avec deux bibliothèques, libcrypto et libssl,  respectivement une implémentation des algorithmes cryptographiques et du protocole de communication SSL/TLS fortement utilisé par internet.
\end{itemize}


\subsection{Common Vulnerability Scoring System (CVSS)}

Bien entendu, disposer d'un identifiant d'une vulnérabilité est important, mais  un gestionnaire de sécurité dans l'entreprise, doit aussi disposer d'élément pour juger de la gravité de cette vulnérabilité. 


% Begin PRZ ===========================
\begin{frame}
\frametitle<presentation>{CVSS}
% end header PRZ =======================
Le  \UKword{Common Vulnerability Scoring System (CVSS) à sa version 3} issu des travaux du  \ulink{https://www.first.org/cvss/}{FIRST, Forum of Incident Response and Security Teams}, est un cadre méthodologique permettant d'évaluer en particulier la criticité d'une vulnérabilité.
\end{frame}
% end PRZ ===========================

%https://www.advens.fr/fr/ressources/blog/regard-critique-sur-cvss

C'est un système permettant de calculer une note évaluant la criticité d'une vulnérabilité, et de construire une chaîne de caractères (un vecteur) présentant les caractéristiques de cette vulnérabilité, et les critères utilisés pour ce calcul.

% Begin PRZ ===========================
\begin{frame}
\frametitle<presentation>{Cotation CVSS}
\framesubtitle<presentation>{Vecteur}
% end header PRZ =======================
Les notes et vecteurs CVSS sont toujours le résultat de trois groupes de critères d'évaluation (\g{Base}, \g{Temporal} et \g{Environnemental}) ayant chacun leur note ainsi que leur vecteur :
\begin{itemize}
  \item Le groupe des critères de\textbf{ \g{Base}} évalue l'impact maximum théorique de la vulnérabilité.
  \item Le groupe des critères\textbf{ \g{Temporel}  }pondère le groupe \g{Basic} en prenant en compte  l'évolution dans le temps de la menace liée à la vulnérabilité  (par exemple, l'existence d'un programme d’exploitation ou d'un correctif).
  \item Le groupe des critères\textbf{ \g{Environnemental} }pondère le groupe \g{Temporel} en prenant en compte les caractéristiques de la vulnérabilité pour un Système d'Information donné.
\end{itemize}
\end{frame}
% end PRZ ===========================

%https://www.cert-ist.com/public/fr/SO_detail?format=html&code=cvss%20v3
% Begin PRZ ===========================
\begin{frame}
\frametitle<presentation>{Cotation CVSS}
\framesubtitle<presentation>{criticité}
% end header PRZ =======================
La richesse du modèle apporte une complexité dans sa lecture rapide, toutefois globalement, on peut lire un score CVSS en terme de criticité avec la grille de lecture suivante :

\begin{itemize}
  \item Un score de 0 à 3.9 correspond à une criticité basse
  \item Un score de 4 à 6.9 correspond à une criticité moyenne
  \item Un score de 7 à 10 correspond à une criticité haute
\end{itemize}
\end{frame}
% end PRZ ===========================


Un autre exemple que je vous engage à explorer pour bien comprendre le fonctionnement CVE sont les vulnérabilités : 
\begin{itemize}
  \item CVE-2020-1023, CVE-2020-1024, and CVE-2020-1102 - SharePoint Remote Code Execution Vulnerability
  \item          CVE-2020-1067 - Windows OS Remote Code Execution Vulnerability
  \item        CVE-2020-1058 (VBScript), CVE-2020-1060 (VBScript), CVE-2020-1064 (Trident=>I.E.) – Internet Explorer Remote Code Execution Vulnerability (which could be used during Web Browsing)
  \item         CVE-2020-1096 - Microsoft Edge PDF Remote Code Execution Vulnerability
\end{itemize}



\subsection{Common Weakness Enumeration (CWE)}

 L'énumération des faiblesses ordinaires ou classiques c'est ainsi qu'il faudrait traduire CWE publié par le \ulink{http://cwe.mitre.org}{MITRE}  est un site qui liste par ailleurs le top des  erreurs de programmation dangereuses et les plus fréquentes fréquentes. En effet, malheureusement les  développeurs font souvent les mêmes erreurs, et malgré un évolution des formations en développement, les développeurs juniors continuent à reproduire les mêmes erreurs de conception et de programmation qui conduisent à des failles logicielles.
La plupart des vulnérabilités applicatives viennent en particulier de quelques erreurs bien connues, qui reviennent régulièrement et pour lesquelles un grande partie des attaques existantes se basent.
C'est le but de la CWE (Common weakness Enumeration) qui est de recenser les erreurs de programmation commises. 
On y retrouve des grands classiques, comme la validation des champs d'un formulaire, la célèbre injection SQL, les problèmes de gestion du système, les contrôles d'accès mal gérés, les tests réalisés par le client plutôt que par le serveur...
Le but du top 25 est d'attirer l'attention des programmeurs sur leurs propres erreurs les plus courantes, mais également de faire réfléchir les formateurs : trop souvent, ces problèmes courants sont oubliés des cours de programmation et de sécurité. Après une brève présentation de chaque problème, la CWE propose des principes généraux pour l'éviter ; le tout est clarifié autant que possible et devrait être compréhensible avec un peu d'effort par la plupart des développeurs. On explorera un peu plus ces éléments dans le chapitre sur la sécurité applicative.

En 2019, les 3 premières faiblesses ordinaires ont été : 

\begin{itemize}
  \item CWE-119	Improper Restriction of Operations within the Bounds of a Memory Buffer	à 75.56\%
  \item CWE-79	Improper Neutralization of Input During Web Page Generation ('Cross-site Scripting')	à 45.69\%
  \item CWE-20	Improper Input Validation à 43.61\%
\end{itemize}




}
\mode<all>{%-------------------------------------
% Chapitre
% Vulnerability Management
% CERT
% File : chap-Vulman-CERT.tex
%--------------------------------------
\uchap{\jobname}

\subsection{Les services de veille en vulnérabilités}

Vous trouverez quelques éléments sur les CERTs (Computer Emergency Response Team) dans le chapitre sur les vulnérabilités, toutefois le périmètre de fonctions et services des CERTs s'est rapidement élargie ces dernières années. Au delà de la diffusion et alertes sur des vulnérabilités, ils couvrent maintenant avec précision les menaces (analyse et alerte sur codes malveillants, ...) et les incidents. Les CERTs restent les acteurs principaux de cette veille et capacité d'alerte.  
Nous trouvons toutefois de nombreux services de veille en vulnérabilités qui ne sont pas des CERTs, mais qui offrent des services dédiés à des typologies de produits, ou des secteurs ....

\subsubsection{Les CERT de l'ANSSI}


le \uac{aCERT} gouvernemental de l'\uac{aANSSI} (\ulink{https://www.cert.ssi.gouv.fr}{CERT GOUV FR}) publie régulièrement plusieurs types d'information :

\begin{itemize}
  \item Alertes de sécurité;
  \item Rapports sur des menaces et incidents;
  \item Avis de sécurité Indicateur de compromission;
  \item Bulletins et notes d'information.
\end{itemize}

% TODO XXXX


\subsubsection{Les CERTs commerciaux}

\begin{techworkbox}{CERT}
	Les CERTs commerciaux est un bon sujet d'exploration des sociétés qui délivrent des services de veille. Vous trouverez 
	\ulink{https://www.ssi.gouv.fr/agence/cybersécurité/ssi-en-france/les-cert-francais/}{les Certs en France} sur le site de l'ANSSI.  Excellent sujet pour \fichetech.
\end{techworkbox}

\subsubsection{La relation avec un CSIRT Interne }
% Analyse de malware. (Un malware utilise des vulnérabilités, et l’analyse d’une attaque donne aussi les éléments pour comber les fragilités

	Une méthodologie efficace de gestion des vulnérabilités comprend une équipe d’intervention en cas d’incident de sécurité informatique (CSIRT). Le CSIRT est responsable de la publication des avis de sécurité, de la tenue d'informations régulières pour échanger sur les activités malveillantes et des dernières attaques du jour zéro, de la simplification et de la diffusion des alertes de sécurité et de l’élaboration de directives compréhensibles et efficaces en matière de réaction aux incidents pour tous les salariés. De cette manière, chacun sera en mesure de réagir aux indicateurs de compromis potentiels conformément aux pratiques recommandées par l'équipe CSIRT.
	
	
\subsection{les agences de notation/cotation Cyber}

\begin{techworkbox}{Cotation Cyber}
	La cotation des entreprises sur le volet CYBER est un sujet en soi. Des entreprises comme (BiSight, Risk IQ, CyberVadis,..) se sont spécialisées dans la cotation du niveau de sécurité d'une entreprise. Sur la base de méthodes, de référentiels et d'algorithmes souvent propriétaires, elles fournissent une cotation CYBER de l'entreprise permettant à ses clients et partenaires de disposer d'une note ou d'un profil de solidité CYBER de leur infrastructure et organisation. Un excellent sujet pour \fichetech.
\end{techworkbox}

%bitsight, risk IQ, Cybervadis

%TODO (01) COTATION CYBER CYBERVADIS - mettre un exemple}
\mode<all>{%-------------------------------------
% Chapitre
% Vulnerability Management
% Gestion
% File : chap-Vulman-gest.tex
%--------------------------------------
\uchap{Introduction sur les vulnérabilités, chapitre 3.1}
\uchap{\jobname}
%_-_-_-_-_-_-_-_-_-_-_-_-_-_-_-_-_-_-_-_-_-_-_-_-_-_-_

\section{GERER les fragilités}

\label{chap:GESTVUL}

Dans le paysage numérique de plus en plus complexe, nous sommes exposés à des terminologies variées souvent soutenues par les modes du moment. Les termes «analyse de vulnérabilités», «évaluation des vulnérabilités» et «gestion des vulnérabilités» sont souvent utilisés et restent une source de confusion pour nombre d’entre nous. Pour nous assurer de se concentrer sur les tactiques les plus efficaces pour gérer les vulnérabilités, nous donnerons les principales différences entre l’évaluation des vulnérabilités et la gestion des vulnérabilités. Mais en posant comme principe que l'important est d'agir quand sont identifiées des failles dans un système.
 La gestion des vulnérabilités est donc un processus de gestion des risques associés à la présence de vulnérabilités qui se base sur la recherche de celles-ci, l'évaluation de leur impact et qui pilote le calendrier d’application des correctifs disponibles.

\begin{itemize}
\item La gestion des vulnérabilités (\textbf{Vulnerability Management}) est un processus continu servant à identifier, classer, corriger et réduire les vulnérabilités, en particulier dans les logiciels. La gestion des vulnérabilités fait partie intégrante des processus de gestion de la cybersécurité dans l’entreprise. Contrairement au projet d’évaluation ponctuelle des vulnérabilités, une stratégie de gestion des vulnérabilités fait référence à un processus ou programme complet et continu qui vise à gérer les vulnérabilités d’une organisation de manière globale et continue. Nous avons rassemblé quelques caractéristiques et éléments clés d’une approche standard de la gestion des vulnérabilités.
La gestion des vulnérabilités comprend aussi le processus par lequel  les risques associés à ces vulnérabilités sont évalués. Cette évaluation conduit à corriger les vulnérabilités et éliminer le risque ou une acceptation formelle du risque par le gestion d’une organisation (par exemple, au cas où l’impact d’une attaque serait faible ou la le coût de la correction ne dépasse pas les dommages éventuels pour l’organisation).
\item Il est souvent confondu avec l’évaluation des vulnérabilités (\textbf{Vulnerability Assessment}), dont l’objectif est de rechercher les fragilités d’un système ou d’une entreprise. Ces vulnérabilités connues sont recherchées sur le système. Une évaluation de vulnérabilité n'est pas une analyse, c'est un projet ponctuel avec une date de début et une date de fin définies. En règle générale, un consultant externe en sécurité de l'information examine votre environnement d'entreprise et identifie diverses vulnérabilités potentiellement exploitables auxquelles vous êtes exposés dans un rapport détaillé. Le rapport répertoriera non seulement les vulnérabilités identifiées, mais fournira également des recommandations concrètes pour la résolution. Une fois le rapport final préparé, l'évaluation de la vulnérabilité est terminée.
Malgré le fait que les deux sont liés, il existe une différence importante entre les deux. La recherche de vulnérabilités consiste à utiliser par exemple un programme informatique pour identifier les vulnérabilités dans réseaux, infrastructure informatique ou applications. La gestion de la vulnérabilité est le processus entourant ce scan de vulnérabilités, prenant également en compte d’autres aspects tels que acceptation des risques, remédiation, etc. On verra en outre que le scan de vulnérabilités n’est qu’une sous partie de l’évaluation des vulnérabilités.
\item L’analyse des vulnérabilités est un processus de recherche de ces fragilités et des scénario qui vont permettre de les exploiter. Les tests d’intrusion sont un exemple de cette dynamique d’analyse des fragilités afin d’en définir un scénario permettant d’atteindre l’objectif que l’attaquant s’est assigné.
\end{itemize}


% Begin PRZ ===========================
\begin{frame}
\frametitle<presentation>{gestion des vulnérabilités}
% end header PRZ =======================
\upicture{Latex/Sources/EDU/Pictures/img-vul-process}{La gestion des vulnérabilités}{01}{lblvulprocess}
\end{frame}
% end PRZ ===========================


%============================
% La gestion des Vulnérabilités - patch management et co
%--------------------------------------------------
\subsection{Processus de gestion des vulnérabilités}

% décrire le processus de gestion avec Inventaire / Analyse / ....
La gestion des vulnérabilités est un processus continu. Elle apparaît en toile de fond du cycle de vie du Maintien en Condition de sécurité :

% Begin PRZ ============
\begin{frame}
\frametitle<presentation>{Cycle de vie VULMAN}
% end header PRZ ===========
\begin{itemize}
	\item Cartographier, cataloguer l'environnement; 
	\item Identifier les fragilités et les menaces;
	\item Corriger, remédier, améliorer la protection et la défense;
	\item Mesurer et suivre l'efficacité les mesures déployées.
\end{itemize}
\end{frame}
% end PRZ ============


\subsubsection{ISO 27001}

Un chapitre de la norme parle de Veille de la vulnérabilités, que nous pouvons classer dans le domaine de la gestion des vulnérabilités et donne des éléments méthodologiques : 

% Begin PRZ ===========================
\begin{frame}
\frametitle<presentation>{ISO 27001}
% end header PRZ =======================
\begin{itemize}
	\item 1. \textbf{DÉCOUVRIR} : Catalogage de l’existant, des actifs, des ressources du système d’information. 
 	\item 2. \textbf{PRIORISER} : Classifier et attribuer des valeurs quantifiables aux ressources, les hiérarchiser. 
 	\item 3. \textbf{ÉVALUER} : Identifier les vulnérabilités ou les menaces potentielles sur chaque ressource. 
	 \item 4. \textbf{SIGNALER} : Signaler, publier les vulnérabilités découvertes. 
 	\item 5. \textbf{CORRIGER} : Éliminer les vulnérabilités les plus sérieuses des ressources les plus importantes. 
 	\item 6. \textbf{VÉRIFIER} : S’assurer que la vulnérabilité a bien été traitée. 
\end{itemize}

\end{frame}
% end PRZ ===========================


%C'est un processus intégré à la gouvernance de la sécurité.


%	Penser la gestion de vulnérabilités 
	
%	
%Évaluation de la vulnérabilité
%Comme mentionné précédemment, une évaluation de vulnérabilité représente en elle-même un élément crucial d'un cadre de gestion des vulnérabilités et est considérée comme le premier pas vers l'amélioration de votre sécurité informatique. De nombreuses entreprises doivent encore faire face à un vaste pool d’actifs inconnus, de périphériques réseau mal configurés, d’environnements très segmentés, d’outils incompatibles ou tout simplement de trop d’informations à analyser et à traiter. Une évaluation des vulnérabilités présente de nombreux avantages et identifie les actifs clés de votre organisation, détermine les vulnérabilités qui menacent la sécurité de ces actifs, fournit des recommandations pour renforcer votre posture de sécurité et vous aide à réduire les risques, vous permettant ainsi de mieux cibler vos ressources informatiques. .
%
%L'analyse des vulnérabilités  permettra un inventaire complet de tous les logiciels et de leurs versions précises, ainsi que la possibilité de vérifier les configurations de sécurité de base et de détecter les vulnérabilités. Ces analyses de vulnérabilité doivent être annoncées pour permettre le signalement des analyses non autorisées et faciliter la visibilité des modifications apportées au réseau et aux actifs. Les processus d'analyse doivent être documentés et examinés pour favoriser leur maturité.
%
%Une fois les évaluations de vulnérabilité effectuées, il est essentiel de produire des rapports clairs et facilement compréhensibles avec des tâches de correction hiérarchisées. Quel que soit l’outil d’analyse de vulnérabilité utilisé, il devrait permettre de produire des rapports, de marquer les vulnérabilités comme ayant été corrigées ou non détectées, de suivre l’âge des vulnérabilités, etc. Avant de publier les rapports, l’organisation doit s’accorder sur le format du rapport afin de s’assurer que les éléments sont inclus / surlignés et les éléments non pertinents sont supprimés. Comme pour de nombreux autres processus de sécurité critiques, il est vivement recommandé que la direction générale adhère pleinement au processus de signalement et de correction des vulnérabilités.


\subsubsection{Fenêtre d'exposition}

%==============
% Analyse / recherche des Vulnérabilités
%------------------------

Dans un monde idéal où en  temps réel, on accéderait à l'apparition d'une vulnérabilité, ou on pourrait la détecter sur son SI, et patcher avec le correctif publié par l'éditeur, la gestion des vulnérabilités se cantonnerait à constater la fenêtre d’exposition générée par le temps nécessaire à l’éditeur pour publier un correctif immédiatement déployé.

% Begin PRZ ===========================
\begin{frame}
\frametitle<presentation>{Fenêtre d'exposition idéale}
% end header PRZ =======================
\upicture{Latex/Sources/EDU/Pictures/img-vul-patchtimeideal}{Fenêtre d'exposition idéale}{01}{lblvulpatchtimeideal}
\end{frame}
% end PRZ ===========================

L’unique décision pour un RSSI serait  serait une gestion du risque simple liée à la décision de déconnecter ou pas un élément vulnérable du SI pendant la fenêtre d’exposition. Tous les correctifs disponibles seraient appliqués dès qu’ils sont disponibles.

Même si l’information peut être accessible en temps réel, la détection de la présence d’une vulnérabilité sur le SI dépend de la fréquence des « scans » ou des audits de l’organisation, ou la connaissance parfaite du SI via une CMDB à jour. Ainsi, la fenêtre d’exposition apparaît naturellement à plusieurs niveaux quelle que soit la réactivité du service de gestion des vulnérabilités.

 L’application d’un correctif nécessite une fenêtre de maintenance et un arrêt du service. Au mieux, ces fenêtres de maintenance sont prévues et planifiées, au pire, l’application de correctifs est interdite hors mode projet (notamment sur certains réseaux opérationnels et industriels). Par ailleurs, la planification et l’application d’un correctif, même dans le meilleur des cas, est extrêmement chronophage pour les équipes techniques.
 
% Begin PRZ ===========================
\begin{frame}
\frametitle<presentation>{Fenêtre d'exposition}
% end header PRZ =======================
\upicture{Latex/Sources/EDU/Pictures/img-vul-patchtime}{Fenêtre d'exposition}{01}{lblvulpatchtime}
\end{frame}
% end PRZ ===========================

Tous les correctifs  pour une même vulnérabilités ne peuvent pas être déployés au même rythme sur tous les systèmes. La gestion des vulnérabilités se structure donc autour de la gestion des priorités basée sur les risques contextuels évalués sur la base  

\begin{itemize}
  \item du niveau de criticité de la vulnérabilité (c’est à dire ce que son exploit permet de réaliser ou d’obtenir – de nombreux services de veille s’appuient sur le Framework CVSS pour classer les vulnérabilités)
  \item du niveau de criticité pour l’organisation des systèmes et services avérés vulnérables.
\end{itemize}


Il est important de disposer d'un accès à un service  de d'information (renseignement) sur les vulnérabilités et sur les Menaces (cf. CERT commercial) afin de collecter et d’organiser le renseignement extérieur et de le croiser avec le renseignement issu de l’interne (détections, incidents, ...) afin d’identifier l'information prioritaire, de l’enrichir et de la diffuser vers les opérations de sécurité. 
 
Un tel service peut être en mesure d’associer une vulnérabilité à un adversaire, à une phase d’attaque, et à une technique particulière employée dans cette phase d’attaque (Cf. Framework MITRE ATTACK facilite notamment un tel niveau de précision et d’analyse).

%TODO Lien MITRE

Il est ainsi possible de mettre en exergue les vulnérabilités exploitées par les adversaires et les campagnes jugées prioritaires  ciblant l’organisation. Cette information sur l’adversaire peut être alors disséminée à des équipes de gestion des vulnérabilités et utilisée pour prioriser l’application des correctifs.



\subsubsection{Processus d’analyse/recherche des vulnérabilités}

Avant de se lancer dans la dynamique classique des audits de sécurité (qui permettent de trouver des vulnérabilité), on peut positionner des mécanismes de recherches de vulnérabilités dans le cycle de vie d'un système sous la forme des différentes étapes des cycles V et V (Vérification et validation) d'un projet.

% TODO les techniques pour recherches les vulnérabilités techniques (Stack Overflow)
% TODO Le secteurs économique de la vulnérabilité. (les failles une business line).

% Begin PRZ ===========================
\begin{frame}
\frametitle<presentation>{Quand rechercher des vulnérabilités}
% end header PRZ =======================
\begin{itemize}
  \item \textbf{Phase de conception} : recherche des défauts et fragilités de conception avec des techniques d'analyse de risque, de revue de conception avec des analyses de menaces
  \item \textbf{Phase de développement} : pendant la phase de développement il existe de nombreux outils d'audit de code statique qui offre l'assistance aux développeurs pour éviter les erreurs les plus classiques,
  \item \textbf{Phase de validation} : dans cette phase, il est possible d'utiliser des techniques et méthodologies classiques d'audit de sécurité (Pentest, analyse de code, ...)
  \item \textbf{Phase de vérification opérationnelle} en Pré-Production ou en production : dans cette phase c'est généralement de l'audit dynamique de type scan de vulnérabilité et tests d'intrusion.
\end{itemize}
\end{frame}
% end PRZ ===========================

Contrairement à l’évaluation des vulnérabilités, un programme complet de gestion des vulnérabilités n’a pas de date de début ni de fin définie, mais constitue un processus continue.

\subsubsection{Processus d’évaluation des vulnérabilités (Vulnerability Assessement)}

%TODO la recherche de vulérabilité SERODIUM

L'évaluation permet de définir l'impact d'une vulnérabilité sur les risques courus par l'entreprise dans une dynamique d'audit ponctuelle ou récurrente?..


%=====================
% Evaluation des Vulnérabilités 
%--------------------------------------

%. Le vulnerability assessement



% Begin PRZ ===========================
\begin{frame}
\frametitle<presentation>{Rechercher ses vulnérabilités}
% end header PRZ =======================
\upicture{Latex/Sources/EDU/Pictures/img-vulsearch}{Rechercher ses vulnérabilités}{1}{lblvulcycle}
\end{frame}
% end PRZ ===========================


\subsection{Audit sécurité des vulnérabilités}

Les audits de sécurité se conduisent souvent sur une base d'audit de vulnérabilités suivi de l'exploitation des ces vulnérabilités ou fragilités pour construire des scénarios plausibles et de caractériser des risques à forte probabilité ou fort impact.
On distingue cependant :
\begin{itemize}
	\item les scans de vulnérabilités, permettant de manière automatisée à rechercher les vulnérabilités sur un système donné, et qui se base sur des bases de vulnérabilités connues
	\item les audits techniques et pentests, qui peuvent se baser en premier lieu sur des scans pour identifier les vulnérabilités connues, mais qui travaillent aussi sur la recherche de vulnérabilités.
\end{itemize}

\subsubsection{Scan de Vulnérabilités du système}

Les vulnérabilités peuvent être découvertes à l'aide d'un scanner de vulnérabilités, qui analyse un système informatique à la recherche de vulnérabilités connues, telles que les ports ouverts, les configurations logicielles non sécurisées et la vulnérabilité aux infections par logiciels malveillants. Des tests de fuzz peuvent permettre de détecter des vulnérabilités inconnues, telles que le jour zéro, permettant d'identifier certains types de vulnérabilités, telles qu'un débordement de mémoire tampon avec des cas de tests pertinents. Une telle analyse peut être facilitée par l’automatisation des tests. 



\subsubsection{Scan de Vulnérabilités logicielles}

La correction des vulnérabilités peut impliquer de différentes manières l’installation d’un correctif, une modification de la stratégie de sécurité du réseau, la reconfiguration du logiciel ou la formation des utilisateurs à l’ingénierie sociale.


\subsection{La gestion des correctifs}

Le patch management est un processus permettant cette gestion des correctifs de sécurité et leur déploiement en entreprise.

%TODO relecture

En anglais, ce patch management consiste à industrialiser les processus de détection, d'analyse et de déploiement des mises à jours de sécurité logicielles. En effet, lorsqu'un éditeur publie un nouveau patch de sécurité relatif à son produit, ses clients ne sont pas toujours en mesure d'évaluer l'importance de ce dernier ni les risques de son installation. Les solutions de gestion des correctifs proposent alors de stocker localement les correctifs sur un serveur du client, puis d'évaluer l'impact de celui-ci avant éventuellement de le tester puis de l'installer.

\subsubsection{De l'outillage}

%TODO relecture

La force d'un outil de gestion des correctifs vient d'abord de sa base de données. C'est elle qui référence le parc informatique existant et conserve l'historique de l'installation des patchs. Elle permet à tout instant de revenir en arrière après installation d'un correctif ou de déterminer précisément où doivent être installés les patchs de sécurité.
Par cette base de connaissance, les logiciels de gestion de correctifs facilitent la détection d'incompatibilités logicielles ou matérielles avec un correctif de sécurité car chaque patch donne lieu à une validation fonctionnelle. Les correctifs installés ceux des éditeurs et dépendent donc de leurs délais de publication.

\subsection{Les audits}

% Begin PRZ ===========================
\begin{frame}
\frametitle<presentation>{Caractéristiques d'audit}
% end header PRZ =======================

Les audits de vulnérabilités s'inscrivent généralement dans des processus de sécurité d'entreprise ou de projets

Les audits peuvent être de natures différentes :

\begin{itemize}
  \item Audit Organisationnel : pour découvrir les fragilités organisationnelles et humaines
  \item Audit technique : pour découvrir et analyser les fragilités 
\end{itemize}

On peut avoir besoin de ces audits pour des enjeux différents :

\begin{itemize}
  \item Audit de conformité
  \item Audit de vérification et de validation
  \item Audit de contrôle et d'inspection
\end{itemize}

Avec une dynamique d'audit :  

\begin{itemize}
 	\item  Audits ponctuels et campagnes d'audit
	\item  Audit continu
\end{itemize}

\end{frame}
% end PRZ ===========================

\begin{techworkbox}{Les audits}
 Les sujets autour des techniques, méthodologies d'audit informatique liés à la sécurité sont une bonne source de réflexion pour \fichetech.
\end{techworkbox}





}
%---------------------------------------------------------------------
\mode<all>{%-------------------------------------
% Chapitre
% Vulnerability Management
% Pentest Intro
% File : chap-Vulman-pentest-intro.tex
%--------------------------------------
\uchap{\jobname}


\section{RECHERCHER des vulnérabilités}

\subsection{Les tests d'intrusion}

Il y existe de nombreuses activités d'évaluation des vulnérabilités. La littérature identifie par ailleurs des noms de métiers ou de compétences d'activité.

% Begin PRZ ===========================
\begin{frame}
\frametitle<presentation>{Les types de tests}
% end header PRZ =======================
\upicture{Latex/Sources/EDU/Pictures/img-vul-typetests}{Les types de tests de vulnérabilités}{0.5}{lblteamstart1}
\end{frame}
% end PRZ ===========================


% Begin PRZ ===========================
\begin{frame}
\frametitle<presentation>{Les branches du test}
% end header PRZ =======================
\upicture{Latex/Sources/EDU/Pictures/img-teamstar}{Les branches du test}{0.5}{lblteamstart}
\end{frame}
% end PRZ ===========================

%https://danielmiessler.com/study/red-blue-purple-teams/

%https://www.bulletproof.co.uk/blog/what-is-penetration-testing

Un des cadres les plus courant pour conduire des audits est l'audit technique qui comprend les tests d'intrusion : PENTEST.

\subsection{Généralités}

Le terme PENTEST est devenu tellement courant, que l'on oublie quelques fois qu'il est l'abréviation de : \g{Penetration Testing}, qui  veut littéralement dire tests d'intrusion.  L'expression  \g{test de pénétration} est parfois rencontrée, mais les professionnels du PENTEST n'apprécient par trop cette expression. 

Il est toujours un peu complexe de catégoriser l'activité de pentests. Il aura assez rapidement des détracteurs pour soulever le fait que la catégorie n'est pas la bonne, qu'elle n'est pas représentative du métier. 
Je vais donc faire rentrer cette activité dans plusieurs catégories (métier de l'audit, métiers d'expertise, métiers du tests).
Dans le cadre de ce cours, je propose de relier cette activité métier comme un des outils du processus de gestion des vulnérabilités (\emph{Vulnerability Management}).  Mais il est plus courant de  classer les activités de PENTESTS dans les activités d'audit.

\subsection{Le métier de Pentesteur}

Le métier du PENTEST est  lié aux métiers techniques de l'informatique et télécom. Les origines des Pentesteur sont très variées.
Ce sont des métiers qu'il est possible d'exercer avec différents niveaux de formation. 

\subsubsection{ Ethical Hackers}
Etre Ethical Hacker fait partie du mythe de la cybersécurité.

\subsubsection{Peut-on faire confiance à des pentesteurs ?}
Parmi les grandes questions que se posent les \g{commanditaires} de tests d'intrusion se trouve celle de la confiance.
En effet, le principe des tests d'intrusion est d'ouvrir un peu les portes de ses systèmes informatiques à des \g{intrus} qui vont certainement découvrir des fragilités.
Certains, peut être plus paranoïaques que d'autres, peuvent se poser la question de savoir ce que vont devenir ces informations sensibles dans \g{les mains} de Hackers.
Parmi les commanditaires, on trouve bien entendu les RSSI, mais aussi les chefs de projet d'application ou de produits embarquant des technologies de informatique ou de communication (Objets intelligents, connectés).

\subsection{Les sociétés de confiance}

Le niveau de sécurité du système d'information et des  outils permettant de réaliser les audits sont vérifiés et validés par une société de certification (LSTI, AMOSSYS ...). A l'issue de certification, l'ANSSI prononce la qualification de la société d'audit au titre de ce PASSI. Il existe une extension pour les audits liés à la loi de programmation militaire (LPM). c'est à dire pour les audits sur les SIIV (Système d'information d'importance vitale) des Opérateurs d'importance Vitale. (Voir chapitre sur la Cyberdefense)

\subsubsection{Formation des Pentests}

\subsection {Certifications professionnelles}

% Begin PRZ ===========================
\begin{frame}
\frametitle<presentation>{2 Certifications}
\framesubtitle<presentation>{reconnues pour les auditeurs techniques}
% end header PRZ =======================
\includer{inc-certifs-hacking.tex}
\end{frame}
% end PRZ ===========================

\subsection {Les rapport d'audits, et cadre méthodologique}

Au coeur de l'activité \g{professionnelle} des pentesteurs se situe le rapport d'audit.
En effet, si l'objectif est bien d'identifier des fragilités (des vulnérabilités), et les scénarios qui permettent de les exploiter à des fins concrètes et relevant d'une menace présentant un risque pour l'entreprise.
il n'en demeure pas moins important de \g{rendre compte} de ce qui a été trouvé. Ce rapport doit aussi contenir des préconisations, car il est important face à une ou des fragilité(s) de proposer des solutions. 




}
\mode<all>{%-------------------------------------
% Chapitre
% Vulnerability Management
% BugBounty
% File : chap-Vulman-vul-bugbounty.tex
%--------------------------------------

\subsection{Le Bug Bounty}

Un bug bounty est une solution de \g{recherche de bugs et de vulnérabilités} proposée par des entreprises organisatrices qui permet à des personnes de recevoir une prime ou compensation après avoir remonté des bugs, surtout ceux concernant des failles et des exploits associés. Dans le domaine de la cybersécurité, on peut résumer de manière macroscopique cette activité  à un déploiement de chasseurs de prime pour chasser les vulnérabilités \g{techniques} des systèmes informatiques.

La question est de savoir où trouver ces chasseurs de primes et d'organiser cette \g{chasse}. Le principe est d'organiser dans l'entreprise ou la structure un programme ciblé sur un système logiciel.

Un des programmes de Bug Bounty  a été lancé au sein d’une entreprise pionnière de l'Internet, développeur du célèbre navigateur: Netscape. Créée en 1994 et rachetée par AOL en 1998, l’entreprise a disparue en janvier 2003.
A l'époque 90\% des internautes utilisaient ce navigateur et Jarrett Ridlinghafer, un ingénieur du support technique, avait constaté que les membres de la communauté \g{open source} corrigeaient d’eux-mêmes les bugs de l'application sans que l'entreprise ne soit dans la boucle de contrôle validation. C'est ainsi que naquit l'idée d'organiser et piloter cette chasse au bug et d'offrir des goodies au titre de récompenses. 
Aujourd'hui des entreprises intermédiaires proposent de réunir tous les conditions pour réaliser ces chasses au bug.
Elle proposent des plateformes permettant de faire se croiser des \g{hackers} et des \g{éditeurs} autour de chasse au bugs / vulnérabilité.

L'utilisation de ces techniques pour trouver des vulnérabilités dans ses systèmes est conseillée en fin de cycle de recherche de vulnérabilités

\subsection{Quelques entreprises}

\begin{itemize}
    \item  (\ulink{https://www.hackerone.com}{HackerOne})
    \item YesWeHack (\ulink{https://www.yeswehack.com}{YesWeHack})
\end{itemize}


}
\mode<all>{%-------------------------------------
% Chapitre
% Vulnerability Management
% Sec By Design
% File : chap-Vulman-secbydesign.tex
%--------------------------------------
\uchap{Vulnérabilités et SEC By DESIGN (VULMAN-INTRO)}

\section{ANTICIPER et construire solide}

Nous n'allons pas développer les concepts de \UKword{Security By design} qui  englobent de nombreuses thématiques de la sécurité applicative.
Les applications en particulier web et mobiles, de par leur complexité et les temps (et parfois budgets) restreints alloués à leur cycle de développement, contiennent souvent un grand nombre de vulnérabilités.

Comme nous l'avons abordé, tester de manière automatisée les vulnérabilités dans du code applicatif développé se décompose en deux typologies de tests :

% Begin PRZ ===========================
\begin{frame}
\frametitle<presentation>{SEC By DESIGN}
% end header PRZ =======================
\begin{itemize}
  \item  \textbf{Audit de code source automatisé} (SAST - Static Application Security Testing). L’audit du code source (SAST) des applications est important si vous souhaitez détecter et corriger leurs vulnérabilités pendant la phase de développement car en effet plus tôt une vulnérabilité est découverte et moins elle sera coûteuse à corriger.
Un audit SAST est non intrusif par nature. Vous pouvez donc scanner en toute sécurité vos applications les plus critiques sans risque d’impacter leur performance.
  \item   \textbf{Audit dynamique automatisé }(DAST - Dynamic Application Security Testing). Un audit dynamique (DAST) consiste à se servir d’un scanner pour interagir avec l’application (avec des requêtes malicieuses vers l’application auditée)  afin d’y trouver des failles connues.
Un scanner de vulnérabilités DAST est plus à même de détecter des erreurs de configuration au serveur web sur lequel est installée l’application.

\end{itemize}
\end{frame}
% end PRZ ===========================

Toutefois parmi celles-ci se trouve la recherche de vulnérabilités au sein des applications dès la phase de conception, et dans les phases de test et de production. Le terme APPSEC (\UKword{Application Security}) est souvent utilisé.

%\section{CORRIGER les vulnérabilités}

\section{Produits de confiance}

\subsection{Les critères communs} 

\label{chap:ISOCC}

\utodo

\utocomplete

% la correction des vulnérabilités dasn le SEC by design, introduduire les concepts de feeb, de sécurity chapion dans les équipês de DEV. SCRUM master }
\mode<all>{%-------------------------
% Chapitre
% Vulnerability Management
% Divers
% File : chap-Vulman-divers.tex
%--------------------------

%===================================
% Recherche de vulnérabilité mais ou
%-------------------------------------------------------------

\subsection{avant propos}
Pour affiner la gestion des vulnérabilités, il y a bien d'autres points à prendre en compte. Nous avons consignés ici ces points qui sont à développer. Nous ne donnons que des pistes de reflexion.
%----------------------------------	
% des difficultés de la cartographie 
\subsection{Périmètre sous responsabilité de l’entreprise}

\subsubsection{la notion de responsabilité}

L'une des premières étapes d'un programme de gestion des vulnérabilités est un exercice de définition du périmètre de responsabilité et d'inventaire des actifs associés. En particulier au niveau de l'entreprise, les entreprises ont tendance à passer par une multitude de fusions, d'acquisitions et de nouvelles technologies et doivent donc combiner des systèmes incompatibles de manière native ou changer de personnel. Malheureusement, ces circonstances laissent souvent les entreprises confuses quant à la qualité de leur inventaire et beaucoup sont incapables d'identifier tous leurs actifs nécessitant un niveau de protection adéquat. Trop souvent, les entreprises possèdent une multitude d'actifs inconnus dans leur environnement qui pourraient compromettre leur sécurité sur le long terme.

Les experts de la sécurité considèrent que la gestion des actifs doit être confiée à une autorité unique qui assure la découverte pertinente dans tous les réseaux et services locaux, valide régulièrement l’inventaire des actifs et gère la gestion des modifications (dont les actifs nouveaux ou actifs retirés). Une fonction centralisée d’inventaire des actifs peut aider à clarifier l’inventaire des actifs d’une organisation et à renforcer le processus de gestion des vulnérabilités sécurité.


\subsubsection{Inventaire des actifs}

Une des vraie difficulté du déploiement d'une gestion de vulnérabilités efficace est la maîtrise des actifs vulnérables ou devant être surveillés et gérés en vulnérabilités.
La gestion des systèmes d'information et des services informatiques (ITSM) est ainsi devenue un processus essentiel de la transformation digitale, considérée comme un outil privilégié qui va soutenir l’entreprise pour affronter sa propre complexité.

Dans un projet ITSM, le référentiel des actifs s’appelle CMDB (Configuration Management DataBase).

Cette base de données de gestion de configuration intègre  tous les composants d’un système d'information pour  avoir une vision d’ensemble sur l’organisation de ces composants et d’en piloter leur configuration en cas de besoin.
Il est donc important de disposer de ce type d'outil pour pouvoir :

\begin{itemize}
  \item Connecter cette CMDB à une solution de veille en vulnérabilités pour corréler les deux et avertir les bons acteurs sur l'apparition d'une vulnérabilité,
\item Gérer les mécanismes de remédiation et de gestion des correctifs.
\end{itemize}

Il n'en demeure pas moins complexe de disposer d'un CMDB à jour, d'autant plus que des services dans le Cloud ne sont pas encore totalement intégrés dans les principes des CMDB, et que le shadow IT  sévit toujours dans les entreprises.

La maîtrise des actifs passe par des outils de d'autodiscovery et d\g{analyse comportementale} qui permet de découvrir non seulement les usages du SI mais aussi découvrir des composants actifs dans l'environnement.

Au coeur de la gestion des vulnérabilités, la gestion des actifs est aussi une histoire de responsabilité des périmètres  informatiques. :

\begin{itemize}
  \item IT métier
  \item Informatique de gestion
  \item Bureautique communicante.
  \item Réseau
\end{itemize}



\begin{techworkbox}{Périmètre et responsabilité}
Qui a la responsabilité d'un périmètre, et dans quelle mesure ce périmètre est géré en configuration ? Quelles sont les adhérences entre les périmètres, quelle fragilité peut induire des niveaux de maturité différents entre ces périmètres, qui peuvent posséder des RSSI différents ... ? Telles sont les questions qui sont en elles-mêmes des sujets de fond dans la gouvernance de la sécurité.
\end{techworkbox}


L'une des premières étapes d'un programme de gestion des vulnérabilités est un exercice de définition du périmètre de responsabilité et d'inventaire des actifs associés. En particulier au niveau de l'entreprise, les entreprises ont tendance à passer par une multitude de fusions, d'acquisitions et de nouvelles technologies et doivent donc combiner des systèmes incompatibles de manière native ou changer de personnel. Malheureusement, ces circonstances laissent souvent les entreprises confuses quant à la qualité de leur inventaire et beaucoup sont incapables d'identifier tous leurs actifs nécessitant un niveau de protection adéquat. Trop souvent, les entreprises possèdent une multitude d'actifs inconnus dans leur environnement qui pourraient compromettre leur sécurité sur le long terme.

%==========================================
\subsubsection{L’environnement digital externe}

%99 ecollece, 2009 Cloud, 2019 ZTA. 2770:2019 privacy
%NIS SP 800 145 et 800 207

Surveiller les failles dans l'entreprise est fondamental, mais il est aussi nécessaire de surveiller les failles apparaissant dans les services Cloud et dans les réseaux sociaux. Ces failles peuvent avoir un impact sur l'entreprise.

%==========================================
\subsection{Veille et alerte sur les vulnérabilités}

%==========================================
\subsubsection{Abonnement au CERT}

S’abonner à un Cert pour être informé en temps réel des vulnérabilités est devenu un besoin primordial pour réagir au plus vite. Il est important que ces alertes soient contextualisées, mais avec la difficulté de ne recevoir que les vulnérabilités qui nous intéressent, le lien avec le patch et la disposition des solutions de corrections.

%==========================================
\subsubsection{Le marché de la vulnérabilité}

Il existe un marché de la vulnérabilité. Les éditeurs (et les états), achètent des vulnérabilités, en gros ils payent des acteurs pour trouver de vulnérabilités.


%==========================================
% Recherche de vulnérabilité en amont du cycle de vie.
%Corriger une vulnérabilité au plus tot
%--------------------------------------------------------------------------
%\subsection{Le marché de l’insécurité chronique}

%	sécurité applicative
%	et le pire avec les DEVOPS et le WEB (Owasp co)

%----------------------------------	
%\subsection{La chasse aux vulnérabilités}
%	Checking versus Pentest
 

}


%$END_UBODY
}%---------------------------------------------------------------------

%************************************************************
% Chargement  du MODELE
%************************************************************

\umainload



